% Trigonométrie dans le triangle rectangle — Mathématiques 3ème — Cours signé OnBuch+
% Programme officiel MINESEC. Compiler avec : tectonic N04-trigonometrie-triangle-rectangle.tex
\newcommand{\DOCMATIERE}{Mathématiques}
\newcommand{\DOCNIVEAU}{3ème}
\newcommand{\DOCMODULE}{Mathématiques – classe de 3ème}
\newcommand{\DOCLECON}{N04}
\newcommand{\DOCTITRE}{Trigonométrie dans le triangle rectangle}
% OnBuch+ — préambule « cours signés » ENRICHI (compilé avec tectonic / XeLaTeX)
% Système visuel « Pop » d'OnBuch+ : fond crème (jamais blanc), bordures encre
% épaisses, ombres « dures » décalées, titres Archivo Black en capitales.
% Version MATHÉMATIQUES : graphes (pgfplots/tikz), boîtes pédagogiques
% (définition, propriété, méthode, attention, le savais-tu, exercice résolu,
% exercice, TP, bilan), page de garde et signature OnBuch+ ancrée en bas.
%
% Macros attendues dans main.tex AVANT \input{preamble} :
%   \DOCMATIERE  \DOCNIVEAU  \DOCMODULE  \DOCLECON (numéro)  \DOCTITRE
\documentclass[11pt]{article}

\usepackage{fontspec}
\usepackage[french]{babel}
\usepackage[a4paper,margin=2cm,top=3.4cm,bottom=2.8cm,headheight=1.4cm,footskip=1.3cm]{geometry}
\usepackage{amsmath,amssymb,amsfonts}
\usepackage{mathtools} % \xmapsto, \coloneqq... souvent utilisés par les modèles
\usepackage{cancel} % simplifications barrées
% \abs{z} (module, valeur absolue) et \norm{u} (norme), utilisés par les modèles
\providecommand{\abs}{}\let\abs\relax\DeclarePairedDelimiter{\abs}{\lvert}{\rvert}
\providecommand{\norm}{}\let\norm\relax\DeclarePairedDelimiter{\norm}{\lVert}{\rVert}
\usepackage[table]{xcolor} % [table] : fournit aussi \rowcolors (alternance de lignes)
\usepackage{pagecolor}
\usepackage{tikz}
\usetikzlibrary{arrows.meta,positioning,calc,shapes.geometric,shapes.misc,shapes.arrows,decorations.pathmorphing,decorations.pathreplacing,decorations.markings,patterns,shadows,mindmap,trees,fit,backgrounds,matrix,intersections,angles,quotes}
\usepackage{pgfplots}
\usepackage{pgf-pie} % diagrammes circulaires \pie (statistiques)
\pgfplotsset{compat=1.18}
\usepgfplotslibrary{fillbetween,groupplots} % aires entre courbes (calcul intégral)
\usepackage{enumitem}
\usepackage{fancyhdr}
% [most] charge toutes les bibliothèques tcolorbox (listings, minted, raster,
% documentation...) et gonfle inutilement la mémoire TeX sur un document long
% (~300 boîtes) : on ne charge que ce qui est réellement utilisé.
\usepackage[skins,breakable]{tcolorbox}
\usepackage{graphicx}
\usepackage{colortbl}
\usepackage{booktabs}
\usepackage{tabularx}
\usepackage{diagbox} % case d'en-tête diagonale des tableaux à double entrée
% Colonnes extensibles centrée / gauche / droite (C, L, R), souvent utilisées par les modèles
\newcolumntype{C}{>{\centering\arraybackslash}X}
\newcolumntype{L}{>{\raggedright\arraybackslash}X}
\newcolumntype{R}{>{\raggedleft\arraybackslash}X}
\usepackage{array}
\usepackage{multicol}
\usepackage{siunitx}
% parse-numbers=false : accepte aussi \SI{2,5 \times 10^{-3}}{...}
\sisetup{locale=FR,output-decimal-marker={,},group-minimum-digits=4,parse-numbers=false}
\DeclareSIUnit\ohm{\ensuremath{\symup{\Omega}}} % U+2126 absent de la police
\usepackage{qrcode}
% \degree : commande fréquemment utilisée par les modèles pour un angle ou une
% température en dehors de \SI{}{} (non fournie par les paquets chargés).
\providecommand{\degree}{^{\circ}}
% \qed (fin de démonstration), utilisé par les modèles sans amsthm
\providecommand{\qed}{\hfill$\square$}
% \barre : double barre des valeurs interdites dans un tableau de variations (tkz-tab)
\providecommand{\barre}{\ensuremath{\|}}
% environnement proof (amsthm non chargé) utilisé par les modèles
\ifdefined\proof\else\newenvironment{proof}[1][Démonstration]{\par\noindent\textit{#1.}\ }{\qed\par}\fi
\usepackage{lastpage}
\usepackage{hyperref}
\hypersetup{hidelinks}

% ============================================================
% Palette « Pop » OnBuch+ — mêmes hex que la classe P (plus_ui.dart)
% ============================================================
\definecolor{poporange}{HTML}{F59321}
\definecolor{popdark}{HTML}{E07A0C}
\definecolor{popcream}{HTML}{FFF6E8}
\definecolor{popink}{HTML}{1C1714}
\definecolor{poppurple}{HTML}{7A5AE0}
\definecolor{popgreen}{HTML}{2E9E6A}
\definecolor{poppink}{HTML}{E0457B}
\definecolor{popblue}{HTML}{2F7DE1}
\definecolor{popgold}{HTML}{C98A12}
\definecolor{popmuted}{HTML}{8A7F76}
\definecolor{popline}{HTML}{EADFD2}
\definecolor{popgray}{HTML}{8A7F76}
\colorlet{popgrey}{popgray}
% Alias tolérants (le modèle nomme parfois les couleurs différemment)
\colorlet{popwhite}{white}
\colorlet{popblack}{popink}
\colorlet{popred}{poppink}
\colorlet{popyellow}{popgold}
% Teintes claires pour fonds de tableaux / zones
\colorlet{popblueL}{popblue!10!white}
\colorlet{poporangeL}{poporange!14!white}
\colorlet{popgreenL}{popgreen!12!white}
\colorlet{poppurpleL}{poppurple!12!white}
\colorlet{poppinkL}{poppink!12!white}
\colorlet{popgoldL}{popgold!14!white}

\pagecolor{popcream}

% ============================================================
% Typographie
% ============================================================
\defaultfontfeatures{Ligatures=TeX}
\setmainfont{PlusJakartaSans-Medium.ttf}[Path=../../../fonts/,BoldFont=PlusJakartaSans-Bold.ttf]
% Maths en sans-serif (Fira Math) avec chiffres et lettres droites de Plus
% Jakarta, pour que les formules se fondent dans le texte.
\usepackage[math-style=ISO,bold-style=ISO]{unicode-math}
\setmathfont{FiraMath-Regular.otf}[Path=../../../fonts/]
\setmathfont{PlusJakartaSans-Medium.ttf}[Path=../../../fonts/,range={up/num,up/{Latin,latin},\mathup}]
\usepackage{icomma} % virgule décimale sans espace en mode math
% Caractères mathématiques Unicode absents des polices de texte (Plus Jakarta,
% Archivo Black) : rendus par leur équivalent mathématique, en texte comme en
% formule — sinon ils s'affichent en carré vide (ex. « Arithmétique dans ℤ »).
\usepackage{newunicodechar}
\newunicodechar{ℝ}{\ensuremath{\mathbb{R}}}
\newunicodechar{ℤ}{\ensuremath{\mathbb{Z}}}
\newunicodechar{ℂ}{\ensuremath{\mathbb{C}}}
\newunicodechar{ℕ}{\ensuremath{\mathbb{N}}}
\newunicodechar{ℚ}{\ensuremath{\mathbb{Q}}}
\newunicodechar{𝔻}{\ensuremath{\mathbb{D}}}
\newunicodechar{α}{\ensuremath{\alpha}}
\newunicodechar{β}{\ensuremath{\beta}}
\newunicodechar{γ}{\ensuremath{\gamma}}
\newunicodechar{δ}{\ensuremath{\delta}}
\newunicodechar{ε}{\ensuremath{\varepsilon}}
\newunicodechar{θ}{\ensuremath{\theta}}
\newunicodechar{λ}{\ensuremath{\lambda}}
\newunicodechar{μ}{\ensuremath{\mu}}
\newunicodechar{σ}{\ensuremath{\sigma}}
\newunicodechar{τ}{\ensuremath{\tau}}
\newunicodechar{φ}{\ensuremath{\varphi}}
\newunicodechar{ψ}{\ensuremath{\psi}}
\newunicodechar{ω}{\ensuremath{\omega}}
\newunicodechar{Σ}{\ensuremath{\Sigma}}
% majuscules produites par \MakeUppercase dans les titres (α-aminés -> Α)
\newunicodechar{Α}{\ensuremath{\alpha}}
\newunicodechar{Β}{\ensuremath{\beta}}
\newunicodechar{Γ}{\ensuremath{\Gamma}}
\newunicodechar{Δ}{\ensuremath{\Delta}}
\newunicodechar{Ε}{\ensuremath{\varepsilon}}
\newunicodechar{Θ}{\ensuremath{\Theta}}
\newunicodechar{Λ}{\ensuremath{\Lambda}}
\newunicodechar{Μ}{\ensuremath{\mu}}
\newunicodechar{Τ}{\ensuremath{\tau}}
\newunicodechar{Φ}{\ensuremath{\Phi}}
\newunicodechar{Ψ}{\ensuremath{\Psi}}
\newunicodechar{Ω}{\ensuremath{\Omega}}
\newunicodechar{↦}{\ensuremath{\mapsto}}
\newunicodechar{∈}{\ensuremath{\in}}
\newunicodechar{∉}{\ensuremath{\notin}}
\newunicodechar{∧}{\ensuremath{\wedge}}
\newunicodechar{⇔}{\ensuremath{\Leftrightarrow}}
\newunicodechar{⟺}{\ensuremath{\Longleftrightarrow}}
\newunicodechar{⇒}{\ensuremath{\Rightarrow}}
\newunicodechar{⟹}{\ensuremath{\Longrightarrow}}
\newunicodechar{∘}{\ensuremath{\circ}}
\newunicodechar{∪}{\ensuremath{\cup}}
\newunicodechar{∩}{\ensuremath{\cap}}
\newunicodechar{≡}{\ensuremath{\equiv}}
\newunicodechar{⊂}{\ensuremath{\subset}}
\newunicodechar{⊥}{\ensuremath{\perp}}
\newunicodechar{∃}{\ensuremath{\exists}}
\newunicodechar{∀}{\ensuremath{\forall}}
\newunicodechar{∛}{\ensuremath{\sqrt[3]{\,}}}
\newunicodechar{∜}{\ensuremath{\sqrt[4]{\,}}}
\newunicodechar{⃗}{\ensuremath{{}^{\to}}}
\newunicodechar{ⁿ}{\ifmmode^{n}\else\textsuperscript{\ensuremath{n}}\fi}
\newunicodechar{ˣ}{\ifmmode^{x}\else\textsuperscript{\ensuremath{x}}\fi}
\newunicodechar{ᵇ}{\ifmmode^{b}\else\textsuperscript{\ensuremath{b}}\fi}
\newunicodechar{ʳ}{\ifmmode^{r}\else\textsuperscript{\ensuremath{r}}\fi}
\newunicodechar{ᵘ}{\ifmmode^{u}\else\textsuperscript{\ensuremath{u}}\fi}
\newunicodechar{ʸ}{\ifmmode^{y}\else\textsuperscript{\ensuremath{y}}\fi}
\newunicodechar{ᵃ}{\ifmmode^{a}\else\textsuperscript{\ensuremath{a}}\fi}
\newunicodechar{ᵉ}{\ifmmode^{e}\else\textsuperscript{\ensuremath{e}}\fi}
\newunicodechar{ᵏ}{\ifmmode^{k}\else\textsuperscript{\ensuremath{k}}\fi}
\newunicodechar{ᵖ}{\ifmmode^{p}\else\textsuperscript{\ensuremath{p}}\fi}
\newunicodechar{ᴺ}{\ifmmode^{N}\else\textsuperscript{\ensuremath{N}}\fi}
\newunicodechar{ᶜ}{\ifmmode^{c}\else\textsuperscript{\ensuremath{c}}\fi}
\newunicodechar{ʰ}{\ifmmode^{h}\else\textsuperscript{\ensuremath{h}}\fi}
\newunicodechar{ᵠ}{\ifmmode^{q}\else\textsuperscript{\ensuremath{q}}\fi}
\newunicodechar{ₙ}{\ifmmode_{n}\else\textsubscript{\ensuremath{n}}\fi}
\newunicodechar{ₐ}{\ifmmode_{a}\else\textsubscript{\ensuremath{a}}\fi}
\newunicodechar{ᵢ}{\ifmmode_{i}\else\textsubscript{\ensuremath{i}}\fi}
\newunicodechar{ᵣ}{\ifmmode_{r}\else\textsubscript{\ensuremath{r}}\fi}
\newunicodechar{ₖ}{\ifmmode_{k}\else\textsubscript{\ensuremath{k}}\fi}
\newunicodechar{ᵦ}{\ifmmode_{\beta}\else\textsubscript{\ensuremath{\beta}}\fi}
\newunicodechar{ₓ}{\ifmmode_{x}\else\textsubscript{\ensuremath{x}}\fi}
\newfontfamily\popdisplay{ArchivoBlack-Regular.ttf}[Path=../../../fonts/]
\newfontfamily\poplogo{SpaceGrotesk-Bold.ttf}[Path=../../../fonts/]
\newfontfamily\popsemi{PlusJakartaSans-SemiBold.ttf}[Path=../../../fonts/]
\newfontfamily\popxbold{PlusJakartaSans-ExtraBold.ttf}[Path=../../../fonts/]
\newfontfamily\popxboldls{PlusJakartaSans-ExtraBold.ttf}[Path=../../../fonts/,LetterSpace=3.0]

\setlist[enumerate]{leftmargin=1.7em,itemsep=5pt,topsep=4pt}
\setlist[itemize]{leftmargin=1.7em,itemsep=4pt,topsep=4pt,label={\color{poporange}\textbullet}}
\setlength{\parindent}{0pt}
\setlength{\parskip}{5pt}
\renewcommand{\arraystretch}{1.25}

% pgfplots : style Pop par défaut
\pgfplotsset{
  every axis/.append style={axis line style={popink,line width=1pt},tick style={popink},
    grid style={popline,line width=0.5pt},label style={font=\small},tick label style={font=\footnotesize}},
  popcurve/.style={poporange,line width=1.8pt,no markers,smooth},
}
% Même style disponible dans un \draw TikZ simple (hors pgfplots)
\tikzset{popcurve/.style={poporange,line width=1.8pt}}

% ============================================================
% Pill / squiggle
% ============================================================
\newcommand{\poppill}[3]{%
  \tikz[baseline=-0.55ex]\node[draw=popink,line width=1.2pt,rounded corners=2.6pt,%
    fill=#2,inner xsep=6.5pt,inner ysep=3pt]{%
    \popxboldls\fontsize{7}{7}\selectfont\color{#3}\MakeUppercase{#1}};%
}
\newcommand{\popsquiggle}[1]{%
  \tikz[baseline=-0.4ex]{
    \draw[#1,line width=2.6pt,line cap=round]
      plot [smooth] coordinates {(0,0.09) (0.34,0.01) (0.68,0.21) (1.02,0.01) (1.36,0.10)};
  }%
}
% Mot-clé surligné (marqueur orange sous le texte)
\newcommand{\cle}[1]{{\popxbold\color{popdark}#1}}

% ============================================================
% En-tête / pied
% ============================================================
\pagestyle{fancy}
\fancyhf{}
\renewcommand{\headrulewidth}{0pt}
\renewcommand{\footrulewidth}{0pt}
\lhead{\raisebox{-0.15ex}{\poplogo\fontsize{13}{13}\selectfont\color{popink}OnBuch\color{poporange}+}}
\rhead{\poppill{\DOCMATIERE\ \textbullet\ \DOCNIVEAU\ \textbullet\ Leçon \DOCLECON}{white}{popink}}
\lfoot{\popxbold\fontsize{7.6}{7.6}\selectfont\color{popmuted}\MakeUppercase{Cours signé OnBuch+}\ \textbullet\ {\popsemi\DOCTITRE}}
\rfoot{\tikz[baseline=-0.6ex]\node[rounded corners=8.5pt,fill=popink,minimum height=17pt,minimum width=17pt,inner xsep=5pt,inner sep=0]{\popxbold\fontsize{8}{8}\selectfont\color{popcream}\thepage\,/\,\pageref*{LastPage}};}

% ============================================================
% Titres
% ============================================================
\newcommand{\chaptitre}[2]{%
  \begin{tcolorbox}[enhanced,colback=poporange,colframe=popink,boxrule=2.6pt,arc=11pt,
      drop shadow=popink,left=15pt,right=15pt,top=12pt,bottom=11pt,before skip=3pt,after skip=18pt]
    \poppill{#2}{popink}{popcream}\par\vspace{9pt}
    {\raggedright\hyphenpenalty=10000\exhyphenpenalty=10000\popdisplay\fontsize{22}{24}\selectfont\color{popcream}\MakeUppercase{#1}\par}
  \end{tcolorbox}
}
% Section numérotée : pastille orange avec le numéro + titre Archivo
\newcounter{popsec}
\newcommand{\coursec}[1]{%
  \stepcounter{popsec}\setcounter{popsub}{0}%
  \par\vspace{16pt}\noindent
  \begin{minipage}{\linewidth}
  \tikz[baseline=-0.7ex]\node[draw=popink,line width=1.6pt,fill=poporange,rounded corners=4pt,minimum size=22pt,inner sep=0]{\popdisplay\fontsize{12}{12}\selectfont\color{popcream}\thepopsec};%
  \hspace{8pt}{\popdisplay\fontsize{15}{17}\selectfont\color{popink}\MakeUppercase{#1}}\par
  \vspace{4pt}\noindent{\color{poporange}\rule{36pt}{3.6pt}}
  \end{minipage}\par\nopagebreak\vspace{8pt}
}
\newcounter{popsub}
\newcommand{\courssub}[1]{%
  \stepcounter{popsub}%
  \par\vspace{9pt}\noindent{\popxbold\fontsize{11.5}{13}\selectfont\color{popink}{\color{poporange}\thepopsec.\thepopsub}\hspace{6pt}#1}\par\nopagebreak\vspace{4pt}
}

% ============================================================
% Boîtes pédagogiques — même recette : fond blanc, bordure épaisse colorée,
% ombre dure encre, badge plein. Titre optionnel [#1].
% ============================================================
\tcbset{popbase/.style={breakable,enhanced,colback=white,boxrule=2.3pt,arc=9pt,
  shadow={3pt}{-3pt}{0pt}{popink},left=14pt,right=14pt,top=11pt,bottom=11pt,before skip=11pt,after skip=15pt,
  fontupper=\popsemi\fontsize{10.6}{14.5}\selectfont\color{popink},
  fontlower=\popsemi\fontsize{10.6}{14.5}\selectfont\color{popink}}}
% Coche dessinée (le glyphe ✓ n'existe pas dans les polices du document)
\newcommand{\popcheck}[1]{\tikz[baseline=-0.5ex]\draw[#1,line width=1.6pt,line cap=round,line join=round](-0.13,0.0)--(-0.04,-0.09)--(0.13,0.1);}
\newcommand{\popboxhead}[3]{\poppill{#1}{#2}{white}\ifx\relax#3\relax\else\hspace{6pt}{\popxbold\fontsize{10.6}{12}\selectfont\color{popink}#3}\fi\par\vspace{7pt}}

\newtcolorbox{aretenir}[1][]{popbase,colframe=popgreen,before upper={\popboxhead{À retenir}{popgreen}{#1}}}
\newtcolorbox{exemplebox}[1][]{popbase,colframe=poporange,before upper={\popboxhead{Exemple}{poporange}{#1}}}
\newtcolorbox{definition}[1][]{popbase,colframe=popblue,before upper={\popboxhead{Définition}{popblue}{#1}}}
\newtcolorbox{propriete}[1][]{popbase,colframe=poppurple,before upper={\popboxhead{Propriété}{poppurple}{#1}}}
\newtcolorbox{methode}[1][]{popbase,colframe=popgold,before upper={\popboxhead{Méthode}{popgold}{#1}}}
\newtcolorbox{attention}[1][]{popbase,colframe=poppink,before upper={\popboxhead{Attention}{poppink}{#1}}}
\newtcolorbox{experience}[1][]{popbase,colframe=popblue,colback=popblueL,before upper={\popboxhead{Expérience}{popblue}{#1}}}
% « Le savais-tu ? » : carte violette pleine (contexte, culture, Cameroun)
\newtcolorbox{savaistu}[1][]{popbase,colback=poppurple,colframe=popink,
  fontupper=\popsemi\fontsize{10.4}{14.2}\selectfont\color{white},
  before upper={\poppill{Le savais-tu\,?}{white}{poppurple}\ifx\relax#1\relax\else\hspace{6pt}{\popxbold\color{white}#1}\fi\par\vspace{7pt}}}
% Exercice résolu : énoncé en haut, \tcblower puis la solution sur fond crème
\newcounter{popexo}\newcounter{popexr}
\newtcolorbox{exoresolu}[1][]{popbase,colframe=poporange,colbacklower=popcream,
  segmentation style={popink,line width=1.2pt,dashed},
  before upper={\stepcounter{popexr}\popboxhead{Exercice résolu \thepopexr}{poporange}{#1}},
  before lower={\poppill{Solution}{popink}{popcream}\par\vspace{6pt}}}
% Exercice (non résolu en place) — la correction est dans la partie Corrigés
\newtcolorbox{exercice}[1][]{popbase,colframe=popink,
  before upper={\stepcounter{popexo}\popboxhead{Exercice \thepopexo}{popink}{#1}}}
% Corrigé
\newtcolorbox{corrige}[1][]{popbase,colframe=popgreen,colback=popgreenL,
  before upper={\popboxhead{Corrigé}{popgreen}{#1}}}
% Objectifs / prérequis / bilan
\newtcolorbox{objectifs}{popbase,colframe=popink,colback=poporangeL,
  before upper={\popboxhead{Objectifs de la leçon}{popink}{}}}
\newtcolorbox{prerequis}{popbase,colframe=popink,colback=popblueL,
  before upper={\popboxhead{Prérequis}{popblue}{}}}
\newtcolorbox{bilan}{popbase,colframe=popink,colback=white,boxrule=2.8pt,
  before upper={\popboxhead{Fiche bilan}{poporange}{L'essentiel à connaître pour l'examen}}}
% Figure encadrée avec légende
\newtcolorbox{popfigure}[1][]{enhanced,breakable=false,colback=white,colframe=popline,boxrule=1.4pt,arc=8pt,
  left=10pt,right=10pt,top=10pt,bottom=8pt,before skip=10pt,after skip=12pt,halign=center,
  lower separated=false,halign lower=center,
  fontlower=\popsemi\fontsize{9}{11.5}\selectfont\color{popmuted},
  #1}
\newcommand{\legende}[1]{\par\vspace{6pt}{\popsemi\fontsize{9}{11.5}\selectfont\color{popmuted}#1}}

% ============================================================
% Signature OnBuch+
% ============================================================
\newcommand{\poplogoinline}[1]{{\poplogo\fontsize{#1}{#1}\selectfont\color{popink}OnBuch\color{poporange}+}}
\newcommand{\signatureonbuch}{%
  \begin{tcolorbox}[enhanced,colback=popink,colframe=popink,boxrule=2.2pt,arc=10pt,
      drop shadow=poporange,left=14pt,right=14pt,top=10pt,bottom=10pt,before skip=0pt,after skip=0pt]
    \tikz[baseline=-0.7ex]\node[circle,fill=popgreen,draw=popcream,line width=1.3pt,minimum size=22pt,inner sep=0]{\popcheck{white}};%
    \hspace{9pt}\begin{minipage}[c]{0.62\linewidth}
      {\popxbold\fontsize{12}{13}\selectfont\color{popcream}Signé\ \textbullet\ {\poplogo OnBuch\color{poporange}+}}\par\vspace{2pt}
      {\raggedright\popsemi\fontsize{8}{10.5}\selectfont\color{popline}\MakeUppercase{Équipe pédagogique OnBuch+}\\\MakeUppercase{Conforme au programme officiel MINESEC}\par}
    \end{minipage}\hfill
    {\poplogo\fontsize{22}{22}\selectfont\color{popcream}OnBuch\color{poporange}+}
  \end{tcolorbox}%
}

% ============================================================
% Page de garde
% #1 titre · #2 matière · #3 niveau · #4 module · #5 n° leçon · #6 durée ·
% #7 description / objectifs (texte ou itemize)
% ============================================================
\newcommand{\pagegarde}[7]{%
  \thispagestyle{empty}
  \newgeometry{a4paper,margin=1.7cm,top=1.6cm,bottom=1.4cm}%
  \begin{tikzpicture}[remember picture,overlay]
    \fill[poporange!18] ([xshift=-3.2cm,yshift=-3.6cm]current page.north east) circle (4.6cm);
    \fill[poppurple!14] ([xshift=2.4cm,yshift=7.6cm]current page.south west) circle (3.8cm);
  \end{tikzpicture}%
  {\poplogo\fontsize{30}{30}\selectfont\color{popink}OnBuch\color{poporange}+}\hfill
  \raisebox{0.6ex}{\poppill{Cours signé \textbullet\ Premium}{popink}{popcream}}\par
  \vspace{1pt}\popsquiggle{poppurple}\par
  \vspace{8pt}
  \begin{tcolorbox}[enhanced,colback=poporange,colframe=popink,boxrule=2.8pt,arc=13pt,
      drop shadow=popink,left=18pt,right=18pt,top=12pt,bottom=13pt,after skip=10pt]
    \poppill{#2 \textbullet\ #3}{popink}{popcream}\hspace{5pt}\poppill{#4}{white}{popink}\par\vspace{8pt}
    {\popdisplay\fontsize{14}{14}\selectfont\color{popink}LEÇON #5}\par\vspace{4pt}
    {\raggedright\hyphenpenalty=10000\exhyphenpenalty=10000\popdisplay\fontsize{25}{27}\selectfont\color{popcream}\MakeUppercase{#1}\par}
    \vspace{8pt}
    {\popxbold\fontsize{9.5}{9.5}\selectfont\color{popink}\MakeUppercase{Durée conseillée : #6}}
  \end{tcolorbox}
  \begin{tcolorbox}[enhanced,colback=white,colframe=popink,boxrule=2.2pt,arc=10pt,
      drop shadow=popink,left=16pt,right=16pt,top=10pt,bottom=10pt,after skip=8pt]
    \poppill{Dans cette leçon}{poppurple}{white}\par\vspace{5pt}
    \popsemi\fontsize{9.3}{12.6}\selectfont\color{popink}#7
  \end{tcolorbox}
  \noindent\begin{minipage}[t]{0.487\linewidth}
    \begin{tcolorbox}[enhanced,colback=popgreen,colframe=popink,boxrule=2.2pt,arc=10pt,
        drop shadow=popink,left=11pt,right=11pt,top=10pt,bottom=10pt,height=3.1cm,valign=center]
      \tcbox[colback=white,colframe=popink,boxrule=1.3pt,arc=5pt,boxsep=0pt,nobeforeafter,
        left=3pt,right=3pt,top=3pt,bottom=3pt,valign=center]{\qrcode[height=1.9cm]{https://play.google.com/store/apps/details?id=com.luvvix.onbuch&hl=fr}}%
      \hspace{8pt}\begin{minipage}[c]{0.5\linewidth}
        {\popdisplay\fontsize{11}{12.5}\selectfont\color{white}\MakeUppercase{Télécharge OnBuch}}\par\vspace{4pt}
        {\raggedright\popsemi\fontsize{8.4}{11}\selectfont\color{white}Gratuit \textbullet\ Google Play\\Tuteur IA, annales, concours, résultats\par}
      \end{minipage}
    \end{tcolorbox}
  \end{minipage}\hfill
  \begin{minipage}[t]{0.487\linewidth}
    \begin{tcolorbox}[enhanced,colback=poppurple,colframe=popink,boxrule=2.2pt,arc=10pt,
        drop shadow=popink,left=11pt,right=11pt,top=10pt,bottom=10pt,height=3.1cm,valign=center]
      \tikz[baseline=-0.7ex]\node[circle,fill=white,draw=popink,line width=1.3pt,minimum size=1.6cm,inner sep=0]{\popdisplay\fontsize{24}{24}\selectfont\color{poppurple}+};%
      \hspace{8pt}\begin{minipage}[c]{0.6\linewidth}
        {\popdisplay\fontsize{11}{12.5}\selectfont\color{white}\MakeUppercase{Passe à OnBuch+}}\par\vspace{4pt}
        {\raggedright\popsemi\fontsize{8.4}{11}\selectfont\color{white}Tous les cours signés, Tuteur IA illimité, fiches PDF — onglet OnBuch+ de l'app\par}
      \end{minipage}
    \end{tcolorbox}
  \end{minipage}
  \vspace{8pt}
  \popcontenu
  \vfill
  \signatureonbuch
  \restoregeometry
  \newpage
}

% Rangée « Ce cours contient » (page de garde)
\newcommand{\poptile}[3]{% #1 couleur #2 gros texte #3 légende
  \begin{tcolorbox}[enhanced,colback=#1,colframe=popink,boxrule=2pt,arc=9pt,shadow={2.5pt}{-2.5pt}{0pt}{popink},
      left=6pt,right=6pt,top=9pt,bottom=9pt,halign=center,valign=center,height=1.75cm,width=\linewidth,nobeforeafter]
    {\popdisplay\fontsize{12.5}{13}\selectfont\color{white}\MakeUppercase{#2}}\par\vspace{3pt}
    {\centering\popsemi\fontsize{7.8}{9.5}\selectfont\color{white}#3\par}
  \end{tcolorbox}}
\newcommand{\popcontenu}{%
  \par\noindent{\popxboldls\fontsize{7.5}{7.5}\selectfont\color{popmuted}CE COURS SIGNÉ CONTIENT}\par\vspace{7pt}
  \noindent\begin{minipage}[t]{0.235\linewidth}\poptile{popblue}{Cours}{complet, illustré, pas à pas}\end{minipage}\hfill
  \begin{minipage}[t]{0.235\linewidth}\poptile{poporange}{Méthodes}{et exercices résolus}\end{minipage}\hfill
  \begin{minipage}[t]{0.235\linewidth}\poptile{poppink}{Exercices}{type examen + corrigés}\end{minipage}\hfill
  \begin{minipage}[t]{0.235\linewidth}\poptile{popgreen}{Fiche bilan}{l'essentiel à retenir}\end{minipage}\par}

% Sommaire
\newtcolorbox{sommaire}{enhanced,colback=white,colframe=popink,boxrule=2.2pt,arc=10pt,
  drop shadow=popink,left=18pt,right=18pt,top=13pt,bottom=13pt,before skip=6pt,after skip=20pt,
  before upper={\poppill{Sommaire}{poppurple}{white}\par\vspace{9pt}\popsemi\fontsize{10.6}{14.5}\selectfont\color{popink}}}

% Signature de fin de cours — poussée tout en bas de la dernière page
\newcommand{\findecours}{%
  \par\vfill\nopagebreak
  \begin{center}\popsquiggle{poporange}\end{center}\vspace{4pt}
  \signatureonbuch
}
\AtBeginDocument{\def\llbracket{\mathopen{[\mkern-3.5mu[}}\def\rrbracket{\mathclose{]\mkern-3.5mu]}}} % intervalles d'entiers (glyphe ⟦ absent de la police)
% Glyphes absents de Fira Math : redéfinis à partir de glyphes disponibles
\AtBeginDocument{%
  \def\setminus{\mathbin{\backslash}}\let\smallsetminus\setminus
  \def\vdots{\mathord{\vbox{\baselineskip3.2pt\lineskiplimit0pt\kern4pt\hbox{.}\hbox{.}\hbox{.}}}}%
  \def\ddots{\mathinner{\mkern1mu\raise6.5pt\hbox{.}\mkern2mu\raise3.6pt\hbox{.}\mkern2mu\raise0.7pt\hbox{.}\mkern1mu}}%
  \def\triangle{\mathord{\scalebox{0.9}{$\symup{\Delta}$}}}%
  \def\nabla{\mathord{\raisebox{\depth}{\scalebox{1}[-1]{$\symup{\Delta}$}}}}%
  \def\checkmark{\popcheck{popdark}}%
  \def\star{\mathbin{\ast}}\def\bigstar{\mathord{\ast}}%
  \def\diamond{\mathbin{\scalebox{0.6}{$\Diamond$}}}%
  \def\bigcup{\mathop{\mathchoice{\raisebox{-0.3em}{\scalebox{1.7}{$\cup$}}}{\scalebox{1.25}{$\cup$}}{\cup}{\cup}}\displaylimits}%
  \def\bigcap{\mathop{\mathchoice{\raisebox{-0.3em}{\scalebox{1.7}{$\cap$}}}{\scalebox{1.25}{$\cap$}}{\cap}{\cap}}\displaylimits}%
  \def\lesseqqgtr{\mathrel{\lessgtr}}\def\bigoplus{\mathop{\oplus}\displaylimits}%
}
\providecommand{\ssi}{si et seulement si} % abréviation fréquente
\AtBeginDocument{\providecommand{\wideparen}{\overparen}} % arc de cercle

%% FIGFIX-BEGIN v5 — correctifs globaux des figures (docs/onbuch-generation/tools/figfix)
\usepackage{adjustbox}\usepackage{etoolbox}\tcbuselibrary{raster}
% toute figure TikZ/pgfplots plus large que la ligne est réduite à la largeur disponible
\BeforeBeginEnvironment{tikzpicture}{\begin{adjustbox}{max width=\linewidth}}
\AfterEndEnvironment{tikzpicture}{\end{adjustbox}}
% légendes pgfplots lisibles par-dessus les courbes
\pgfplotsset{every axis/.append style={legend style={fill=white,fill opacity=0.92,text opacity=1,draw=black!15,inner sep=3pt}}}
% étiquettes d'arêtes (syntaxe edge["..."]) avec fond blanc
\tikzset{every edge quotes/.append style={fill=white,inner sep=1.2pt}}
%% FIGFIX-END
\begin{document}

\pagegarde{Trigonométrie dans le triangle rectangle}{Mathématiques}{3ème}{Mathématiques – classe de 3ème}{N04}{3 séances (fiche de progression harmonisée)}{Cette leçon introduit les rapports trigonométriques (cosinus, sinus, tangente) dans le triangle rectangle pour calculer des longueurs ou des angles aigus. Elle vise l'autonomie dans la modélisation de situations réelles (hauteurs inaccessibles, pentes, orientations) et la rédaction rigoureuse de démonstrations.
\vspace{4pt}
\begin{multicols}{2}\small
\begin{itemize}
\item Identifier les côtés (hypoténuse, adjacent, opposé) par rapport à un angle aigu dans un triangle rectangle.
\item Définir et calculer le cosinus, le sinus et la tangente d'un angle aigu à l'aide des rapports de longueurs.
\item Utiliser la calculatrice pour déterminer une valeur trigonométrique ou un angle à partir de son rapport (fonctions inverses).
\item Connaître et appliquer les valeurs exactes des rapports pour les angles 30°, 45° et 60°.
\item Calculer une longueur inconnue dans un triangle rectangle connaissant un angle aigu et un côté.
\item Calculer la mesure d'un angle aigu connaissant deux côtés d'un triangle rectangle.
\end{itemize}\end{multicols}}

\begin{sommaire}
\begin{multicols}{2}
\begin{enumerate}
\item 1. Repérage et vocabulaire dans le triangle rectangle
\item 2. Définition des rapports trigonométriques : Cosinus, Sinus, Tangente
\item 3. Angles remarquables (30°, 45°, 60°) et valeurs exactes
\item 4. Calculer une longueur dans un triangle rectangle
\item 5. Calculer un angle aigu dans un triangle rectangle
\item 6. Résolution de triangle rectangle \& Problèmes concrets (Intégration)
\item Activité d'intégration
\item Méthodes et exercices résolus
\item Exercices
\item Corrigés des exercices
\item Fiche bilan
\end{enumerate}
\end{multicols}
\end{sommaire}

\begin{objectifs}
\begin{itemize}
\item Identifier les côtés (hypoténuse, adjacent, opposé) par rapport à un angle aigu dans un triangle rectangle.
\item Définir et calculer le cosinus, le sinus et la tangente d'un angle aigu à l'aide des rapports de longueurs.
\item Utiliser la calculatrice pour déterminer une valeur trigonométrique ou un angle à partir de son rapport (fonctions inverses).
\item Connaître et appliquer les valeurs exactes des rapports pour les angles 30°, 45° et 60°.
\item Calculer une longueur inconnue dans un triangle rectangle connaissant un angle aigu et un côté.
\item Calculer la mesure d'un angle aigu connaissant deux côtés d'un triangle rectangle.
\item Résoudre complètement un triangle rectangle (trouver tous les côtés et angles) à partir de données minimales.
\item Modéliser et résoudre un problème concret de la vie courante (hauteur d'arbre, pente de toit, distance infranchissable) en justifiant chaque étape.
\end{itemize}
\end{objectifs}

\begin{prerequis}
\begin{itemize}
\item Théorème de Pythagore et sa réciproque (4ème).
\item Théorème de Thalès et propriétés de similitude (4ème).
\item Notion d'angle aigu, triangle rectangle, hypoténuse (6ème/5ème).
\item Calcul littéral : manipulation de fractions, produit en croix, arrondis (4ème).
\item Utilisation basique de la calculatrice (touches DEG, SIN, COS, TAN, touches inverses).
\end{itemize}
\end{prerequis}

\newpage

\coursec{Repérage et vocabulaire dans le triangle rectangle}

\courssub{Une situation concrète pour ouvrir la leçon}

Monsieur Fouda, menuisier à Yaoundé, vient de recevoir une commande importante : il doit fabriquer une \cle{rampe d'accès} pour personnes à mobilité réduite devant l'entrée d'une école du quartier Nkolbisson. La marche d'entrée fait \SI{40}{cm} de haut. La norme d'accessibilité impose une \cle{inclinaison maximale} de $15^{\circ}$ par rapport à l'horizontale.

Monsieur Fouda se pose alors deux questions très concrètes :
\begin{itemize}
\item Quelle doit être la \textbf{longueur minimale} de la planche qui formera la rampe ?
\item À quelle \textbf{distance du mur} doit-il poser le pied de la rampe ?
\end{itemize}

En dessinant la situation, on s'aperçoit que la rampe, le sol et le mur forment un \textbf{triangle rectangle} : le sol est horizontal, le mur est vertical, et la rampe joue le rôle du côté incliné. On connaît un côté (la hauteur, \SI{40}{cm}) et un angle ($15^{\circ}$), et l'on cherche les deux autres côtés. C'est exactement le genre de problème que la \cle{trigonométrie} permet de résoudre.

\begin{center}
\begin{popfigure}
\begin{tikzpicture}[scale=0.9]
\coordinate (A) at (0,0);
\coordinate (B) at (6,0);
\coordinate (C) at (6,1.6);
\draw[thick, popdark] (A)--(B)--(C)--cycle;
\draw[pattern=north east lines, pattern color=popmuted] (6,0) rectangle (6.5,1.6);
\draw[popblue, thick] (0.9,0) arc (0:15:0.9);
\node[popblue] at (1.35,0.28) {$15^{\circ}$};
\node[below] at (3,0) {distance au sol ?};
\node[right, popdark] at (6.25,0.8) {$40$ cm};
\node[above left, poporange] at (2.9,0.95) {rampe ?};
\draw[popdark] (5.7,0) rectangle (6,0.3);
\node[below left] at (A) {pied};
\node[below right] at (B) {mur};
\end{tikzpicture}
\legende{La rampe de Monsieur Fouda modélisée par un triangle rectangle.}
\end{popfigure}
\end{center}

\textbf{Problématique de la leçon.} Dans un triangle rectangle, existe-t-il des relations entre les \emph{angles} et les \emph{longueurs des côtés} qui permettraient de calculer des distances inaccessibles à la règle ? Pour y répondre, il faut d'abord apprendre à \textbf{nommer correctement} les côtés d'un triangle rectangle : c'est l'objet de cette première section.

\courssub{Rappel : le triangle rectangle et son hypoténuse}

Tu as rencontré le triangle rectangle dès la classe de 6\textsuperscript{e}. Rappelons sa définition.

\begin{definition}[Triangle rectangle et hypoténuse]
Un \cle{triangle rectangle} est un triangle qui possède un \textbf{angle droit} ($90^{\circ}$).
Le côté \textbf{opposé à l'angle droit} s'appelle l'\cle{hypoténuse}. C'est toujours le \textbf{plus long} des trois côtés.
Les deux autres angles du triangle sont \textbf{aigus} (leur mesure est comprise entre $0^{\circ}$ et $90^{\circ}$) et leur somme vaut $90^{\circ}$ : on dit qu'ils sont \textbf{complémentaires}.
\end{definition}

\begin{exemplebox}[Vérification rapide]
Dans un triangle $ABC$ rectangle en $A$ :
\begin{itemize}
\item l'angle droit est $\widehat{BAC} = 90^{\circ}$ ;
\item l'hypoténuse est le côté $[BC]$, celui qui ne touche pas le sommet $A$ ;
\item les angles aigus sont $\widehat{ABC}$ et $\widehat{ACB}$, et $\widehat{ABC} + \widehat{ACB} = 90^{\circ}$.
\end{itemize}
\end{exemplebox}

\begin{attention}[L'hypoténuse ne change jamais]
Dans un triangle rectangle donné, l'hypoténuse est \textbf{toujours le même côté} : celui qui fait face à l'angle droit. Peu importe l'angle aigu que l'on étudie, l'hypoténuse reste l'hypoténuse. C'est le seul côté dont le nom ne dépend pas de l'angle choisi.
\end{attention}

\courssub{Nommer les côtés par rapport à un angle aigu}

Voici l'idée centrale de toute la trigonométrie. Imagine que tu te places « dans » l'un des angles aigus du triangle, comme un observateur. De ton point de vue :
\begin{itemize}
\item un côté est juste \textbf{en face de toi} : c'est le \emph{côté opposé} ;
\item un côté \textbf{te touche} et monte vers l'angle droit : c'est le \emph{côté adjacent} (adjacent signifie « voisin ») ;
\item le dernier côté te touche aussi, mais c'est l'\emph{hypoténuse}, déjà nommée.
\end{itemize}

\begin{definition}[Côtés adjacent et opposé à un angle aigu]
Dans un triangle rectangle, par rapport à un \textbf{angle aigu} donné :
\begin{itemize}
\item le côté opposé à l'\textbf{angle droit} est l'\cle{hypoténuse} ;
\item le côté qui \textbf{forme l'angle} avec l'hypoténuse (autrement dit le côté de l'angle qui n'est pas l'hypoténuse) est le \cle{côté adjacent} à cet angle ;
\item le troisième côté, qui est \textbf{en face} de l'angle, est le \cle{côté opposé} à cet angle.
\end{itemize}
\end{definition}

Le point essentiel à comprendre est que les mots « adjacent » et « opposé » sont \textbf{relatifs} : ils dépendent de l'angle aigu que l'on choisit d'étudier. Le même côté peut être opposé à un angle et adjacent à l'autre !

\begin{center}
\begin{popfigure}
\begin{tikzpicture}[scale=1.1]
\coordinate (A) at (0,0);
\coordinate (B) at (4,0);
\coordinate (C) at (0,3);
\draw[very thick, popdark] (A)--(B)--(C)--cycle;
\draw[popdark] (0,0.35) rectangle (0.35,0);
\draw[popblue, very thick] (A)--(B);
\draw[poporange, very thick] (A)--(C);
\draw[poppurple, very thick] (B)--(C);
\draw[popgreen, thick] (0.9,0) arc (0:36.87:0.9);
\node[popgreen] at (1.3,0.32) {$\widehat{B}$};
\node[below left, popdark] at (A) {$A$};
\node[below right, popdark] at (B) {$B$};
\node[above left, popdark] at (C) {$C$};
\node[below, popblue] at (2,0) {côté adjacent à $\widehat{B}$ : $[AB]$};
\node[left, poporange, rotate=90] at (-0.15,1.5) {côté opposé à $\widehat{B}$ : $[AC]$};
\node[above right, poppurple, rotate=-36.87] at (2.1,1.75) {hypoténuse : $[BC]$};
\end{tikzpicture}
\legende{Repérage des côtés par rapport à l'angle $\widehat{B}$ dans le triangle $ABC$ rectangle en $A$.}
\end{popfigure}
\end{center}

\begin{center}
\begin{popfigure}
\begin{tikzpicture}[scale=1.1]
\coordinate (A) at (0,0);
\coordinate (B) at (4,0);
\coordinate (C) at (0,3);
\draw[very thick, popdark] (A)--(B)--(C)--cycle;
\draw[popdark] (0,0.35) rectangle (0.35,0);
\draw[poporange, very thick] (A)--(C);
\draw[popblue, very thick] (A)--(B);
\draw[poppurple, very thick] (B)--(C);
\draw[popgreen, thick] (0,2.1) arc (-90:-53.13:0.9);
\node[popgreen] at (0.42,2.05) {$\widehat{C}$};
\node[below left, popdark] at (A) {$A$};
\node[below right, popdark] at (B) {$B$};
\node[above left, popdark] at (C) {$C$};
\node[below, popblue] at (2,0) {côté opposé à $\widehat{C}$ : $[AB]$};
\node[left, poporange, rotate=90] at (-0.15,1.5) {côté adjacent à $\widehat{C}$ : $[AC]$};
\node[above right, poppurple, rotate=-36.87] at (2.1,1.75) {hypoténuse : $[BC]$};
\end{tikzpicture}
\legende{Même triangle, mais repérage par rapport à l'angle $\widehat{C}$ : les rôles de $[AB]$ et $[AC]$ sont échangés !}
\end{popfigure}
\end{center}

Observe bien les deux figures : le côté $[AB]$ est \textbf{adjacent} à l'angle $\widehat{B}$ mais \textbf{opposé} à l'angle $\widehat{C}$. Le côté $[AC]$ fait le chemin inverse. Seule l'hypoténuse $[BC]$ garde son nom dans les deux cas.

\begin{methode}[Repérer les côtés par rapport à un angle aigu]
Pour nommer correctement les côtés d'un triangle rectangle par rapport à un angle aigu :
\begin{enumerate}
\item \textbf{Repérer l'angle droit} : le côté qui lui fait face est l'\textbf{hypoténuse}.
\item \textbf{Regarder l'angle aigu étudié} : le côté qui \emph{forme} cet angle avec l'hypoténuse est le \textbf{côté adjacent}.
\item Le côté restant, \emph{en face} de l'angle, est le \textbf{côté opposé}.
\end{enumerate}
Astuce : sur ta figure, marque l'angle étudié d'un arc de couleur, puis entoure les deux côtés qui le forment : l'un est l'hypoténuse, l'autre est l'adjacent.
\end{methode}

\begin{exoresolu}[Exercice de vocabulaire : changer d'angle]
\textbf{Énoncé.} Soit $ABC$ un triangle rectangle en $A$ tel que $AB = 4$ cm, $AC = 3$ cm et $BC = 5$ cm.
\begin{enumerate}
\item Par rapport à l'angle $\widehat{ABC}$, nommer l'hypoténuse, le côté adjacent et le côté opposé, et donner leurs longueurs.
\item Reprendre la même question par rapport à l'angle $\widehat{ACB}$.
\end{enumerate}
\tcblower
\textbf{Solution détaillée.}
\begin{enumerate}
\item \textbf{Par rapport à l'angle $\widehat{B}$ :}
\begin{itemize}
\item l'angle droit est en $A$, donc l'\textbf{hypoténuse} est $[BC]$, de longueur $5$ cm ;
\item les côtés qui forment l'angle $\widehat{B}$ sont $[BA]$ et $[BC]$ ; comme $[BC]$ est l'hypoténuse, le \textbf{côté adjacent} est $[AB]$, de longueur $4$ cm ;
\item le côté en face de $\widehat{B}$ est $[AC]$ : c'est le \textbf{côté opposé}, de longueur $3$ cm.
\end{itemize}
\item \textbf{Par rapport à l'angle $\widehat{C}$ :}
\begin{itemize}
\item l'\textbf{hypoténuse} est toujours $[BC]$, de longueur $5$ cm (elle ne change pas !) ;
\item les côtés qui forment l'angle $\widehat{C}$ sont $[CA]$ et $[CB]$ ; le \textbf{côté adjacent} est donc $[AC]$, de longueur $3$ cm ;
\item le côté en face de $\widehat{C}$ est $[AB]$ : c'est le \textbf{côté opposé}, de longueur $4$ cm.
\end{itemize}
\end{enumerate}
\begin{center}
\begin{tabular}{|l|c|c|c|}
\hline
\rowcolor{popblueL} Angle étudié & Hypoténuse & Côté adjacent & Côté opposé \\
\hline
$\widehat{B}$ & $BC = 5$ & $AB = 4$ & $AC = 3$ \\
\hline
$\widehat{C}$ & $BC = 5$ & $AC = 3$ & $AB = 4$ \\
\hline
\end{tabular}
\end{center}
On constate que $[AB]$ et $[AC]$ échangent leurs rôles quand on change d'angle, tandis que l'hypoténuse reste $[BC]$.
\end{exoresolu}

\begin{attention}[Les trois pièges classiques]
\begin{itemize}
\item \textbf{Confondre côté adjacent et hypoténuse.} L'angle aigu est formé par \emph{deux} côtés : l'hypoténuse et l'adjacent. Commence toujours par identifier l'hypoténuse (face à l'angle droit) pour éliminer la confusion.
\item \textbf{Croire que l'adjacent est fixe.} Le côté adjacent \emph{dépend de l'angle choisi} : il change quand on passe de $\widehat{B}$ à $\widehat{C}$. Seule l'hypoténuse est invariante.
\item \textbf{Oublier que l'hypoténuse est le plus grand côté.} Si ton calcul donne un côté adjacent ou opposé plus long que l'hypoténuse, c'est forcément faux : vérifie ton repérage.
\end{itemize}
\end{attention}

\courssub{Pourquoi ces noms ? Le lien avec les rapports de côtés (Thalès)}

On pourrait se demander : pourquoi prendre tant de soin à nommer les côtés ? La réponse est profonde, et c'est le théorème de Thalès qui nous la donne.

\begin{experience}[Activité de découverte : des triangles qui se ressemblent]
Trace un angle $\widehat{XBY} = 35^{\circ}$. Sur la demi-droite $[BX$, place trois points $A$, $A'$, $A''$ tels que $BA = 3$ cm, $BA' = 5$ cm, $BA'' = 7$ cm. Par chacun de ces points, trace la perpendiculaire à $(BX)$ ; elle coupe la demi-droite $[BY$ respectivement en $C$, $C'$, $C''$.
On obtient ainsi trois triangles rectangles : $ABC$ (rectangle en $A$), $A'B'C'$ (rectangle en $A'$) et $A''B''C''$ (rectangle en $A''$), qui ont tous le même angle aigu en $B$.
\begin{enumerate}
\item Mesure les côtés opposés $AC$, $A'C'$, $A''C''$.
\item Calcule pour chaque triangle le rapport $\dfrac{\text{côté opposé à } \widehat{B}}{\text{côté adjacent à } \widehat{B}}$.
\item Que constates-tu ? Recommence avec le rapport $\dfrac{\text{côté opposé}}{\text{hypoténuse}}$.
\end{enumerate}
Tu dois trouver, aux erreurs de mesure près, \textbf{le même nombre} pour les trois triangles. Ce n'est pas un hasard : expliquons-le.
\end{experience}

\begin{propriete}[Constance des rapports de côtés]
Si deux triangles rectangles possèdent un angle aigu de \textbf{même mesure}, alors ils ont des côtés \textbf{proportionnels}. En particulier, dans tout triangle rectangle, le rapport de deux côtés (opposé sur adjacent, opposé sur hypoténuse, adjacent sur hypoténuse) \textbf{ne dépend que de la mesure de l'angle aigu} considéré, et pas de la taille du triangle.
\end{propriete}

\textbf{Justification par le théorème de Thalès.} Considérons deux triangles rectangles $ABC$ et $A'B'C'$, rectangles respectivement en $A$ et $A'$, qui partagent le même angle aigu en $B$ (on a donc $B = B'$). Les points $A'$ et $C'$ sont choisis sur les segments $[BA]$ et $[BC]$.
Les droites $(AC)$ et $(A'C')$ sont toutes deux perpendiculaires à la droite $(BA)$ : elles sont donc \textbf{parallèles}. D'après le théorème de Thalès appliqué au triangle $ABC$ et à la parallèle $(A'C')$ :
\[ \frac{BA'}{BA} = \frac{BC'}{BC} = \frac{A'C'}{AC}. \]
On en déduit, par exemple, que $\dfrac{A'C'}{BA'} = \dfrac{AC}{BA}$ : le rapport $\dfrac{\text{opposé}}{\text{adjacent}}$ est le \textbf{même} dans les deux triangles. Le même raisonnement vaut pour les deux autres rapports. \hfill $\square$

\begin{aretenir}[L'idée fondatrice de la trigonométrie]
À chaque angle aigu (par exemple $30^{\circ}$, $45^{\circ}$, $15^{\circ}$...) sont associés des \textbf{nombres constants} : les rapports entre les côtés de n'importe quel triangle rectangle possédant cet angle. Ces nombres sont si utiles qu'on leur a donné des \textbf{noms} : \cle{sinus}, \cle{cosinus} et \cle{tangente}. C'est pour pouvoir écrire ces rapports sans ambiguïté qu'il faut d'abord savoir nommer les côtés : \emph{hypoténuse}, \emph{adjacent}, \emph{opposé}.
\end{aretenir}

\begin{savaistu}[Un mot d'histoire]
Le mot « trigonométrie » vient du grec \emph{trigônon} (triangle) et \emph{metron} (mesure) : c'est littéralement la « mesure des triangles ». Les astronomes de l'Antiquité (comme Hipparque, puis les savants arabes comme Al-Khwarizmi et Al-Battani) ont construit les premières tables de rapports trigonométriques pour prédire les positions des astres. Aujourd'hui encore, ces mêmes rapports servent aux géomètres pour cartographier le terrain, aux ingénieurs de CAMTEL pour orienter les antennes relais, et aux menuisiers comme Monsieur Fouda pour construire des rampes aux normes !
\end{savaistu}

\begin{exercice}[À toi de jouer]
Soit $MNP$ un triangle rectangle en $M$.
\begin{enumerate}
\item Quelle est son hypoténuse ?
\item Par rapport à l'angle $\widehat{N}$, nommer le côté adjacent et le côté opposé.
\item Par rapport à l'angle $\widehat{P}$, nommer le côté adjacent et le côté opposé.
\item Que peut-on dire des angles $\widehat{N}$ et $\widehat{P}$ ?
\end{enumerate}
\end{exercice}

\begin{corrige}[Exercice 1]
\begin{enumerate}
\item L'angle droit est en $M$, donc l'hypoténuse est le côté opposé à $M$ : $[NP]$.
\item Par rapport à $\widehat{N}$ : les côtés formant l'angle sont $[NM]$ et $[NP]$ ; l'adjacent est $[MN]$ et l'opposé est $[MP]$.
\item Par rapport à $\widehat{P}$ : l'adjacent est $[MP]$ et l'opposé est $[MN]$.
\item $\widehat{N} + \widehat{P} = 90^{\circ}$ : ce sont des angles complémentaires.
\end{enumerate}
\end{corrige}

\textbf{En résumé de cette section.} Tu sais désormais repérer, dans tout triangle rectangle, l'hypoténuse (fixe, face à l'angle droit), le côté adjacent et le côté opposé (qui dépendent de l'angle aigu choisi). Tu sais aussi \emph{pourquoi} ce vocabulaire existe : grâce à Thalès, les rapports entre ces côtés ne dépendent que de l'angle. Dans la prochaine section, nous donnerons un nom à ces rapports — sinus, cosinus, tangente — et nous pourrons enfin répondre à la question de Monsieur Fouda.

\coursec{Définition des rapports trigonométriques : Cosinus, Sinus, Tangente}

Dans la section précédente, nous avons appris à \cle{repérer} dans un triangle rectangle l'hypoténuse, le côté adjacent et le côté opposé à un angle aigu donné. Nous savons donc maintenant nommer correctement les trois côtés \emph{relativement à un angle aigu choisi}. Mais à quoi cela sert-il ?

Imagine la situation suivante : un technicien de CAMTEL doit vérifier l'inclinaison d'un pylône, ou un géomètre à Yaoundé veut connaître la hauteur d'une colline sans la gravir. Mesurer des angles est facile (avec un rapporteur ou un théodolite) ; mesurer certaines longueurs est parfois impossible. La \cle{trigonométrie} est précisément l'outil mathématique qui relie les \textbf{angles} et les \textbf{longueurs} dans un triangle rectangle. Le mot lui-même le dit : \emph{trigono} (triangle) et \emph{métrie} (mesure). Dans cette section, nous définissons les trois rapports fondamentaux : le \cle{cosinus}, le \cle{sinus} et la \cle{tangente}.

\courssub{Une intuition avant les formules}

Considérons une échelle posée contre un mur. Si l'échelle est presque verticale, son pied est très proche du mur ; si elle est très inclinée, son pied s'éloigne beaucoup du mur. Le rapport $\dfrac{\text{distance du pied au mur}}{\text{longueur de l'échelle}}$ décrit donc \emph{l'inclinaison} de l'échelle : il ne dépend pas de la longueur de l'échelle, mais seulement de l'angle qu'elle fait avec le sol.

C'est exactement l'idée de la trigonométrie : dans un triangle rectangle, les \textbf{rapports de longueurs} entre les côtés caractérisent les angles aigus. Deux triangles rectangles qui ont le même angle aigu $\alpha$ ont la même forme (ils sont \emph{semblables}), donc les rapports de leurs côtés sont égaux, même si leurs tailles diffèrent.

\courssub{Définition des trois rapports trigonométriques}

Plaçons-nous dans un triangle $ABC$ rectangle en $A$, et intéressons-nous à l'angle aigu $\widehat{ABC}$ que nous noterons $\alpha$. Par rapport à cet angle :
\begin{itemize}
\item l'\textbf{hypoténuse} est $[BC]$ (le côté opposé à l'angle droit, le plus long) ;
\item le \textbf{côté adjacent} à $\alpha$ est $[AB]$ (le côté de l'angle droit qui \emph{touche} $\alpha$) ;
\item le \textbf{côté opposé} à $\alpha$ est $[AC]$ (le côté qui est \emph{en face} de $\alpha$).
\end{itemize}

\begin{definition}[Cosinus, sinus et tangente d'un angle aigu]
Soit un triangle rectangle et $\alpha$ un de ses angles aigus. On définit trois nombres, appelés \cle{rapports trigonométriques} de l'angle $\alpha$ :
\[ \cos(\alpha) = \dfrac{\text{côté adjacent à } \alpha}{\text{hypoténuse}} \qquad ; \qquad \sin(\alpha) = \dfrac{\text{côté opposé à } \alpha}{\text{hypoténuse}} \]
\[ \tan(\alpha) = \dfrac{\text{côté opposé à } \alpha}{\text{côté adjacent à } \alpha} \]
Dans le triangle $ABC$ rectangle en $A$, avec $\alpha = \widehat{ABC}$ :
\[ \cos(\alpha) = \dfrac{AB}{BC} \qquad ; \qquad \sin(\alpha) = \dfrac{AC}{BC} \qquad ; \qquad \tan(\alpha) = \dfrac{AC}{AB} \]
\end{definition}

\begin{attention}[Bien choisir l'angle de référence]
Les mots \og adjacent \fg{} et \og opposé \fg{} n'ont de sens que \textbf{par rapport à l'angle choisi}. Dans le triangle $ABC$ rectangle en $A$, le côté $[AC]$ est opposé à l'angle $\widehat{ABC}$, mais il est adjacent à l'angle $\widehat{ACB}$. Avant d'écrire un rapport, commence toujours par dire : \og je travaille avec l'angle ... \fg{}, puis repère les trois côtés par rapport à cet angle. L'hypoténuse, elle, ne change jamais.
\end{attention}

\begin{popfigure}
\begin{center}
\begin{tikzpicture}[scale=1.6, line cap=round, line join=round]
\coordinate (A) at (0,0);
\coordinate (B) at (4,0);
\coordinate (C) at (0,3);
\draw[thick, popdark] (A) -- (B) -- (C) -- cycle;
% Right angle mark at A
\draw[popdark] (0.3,0) -- (0.3,0.3) -- (0,0.3);
% Angle alpha at B
\draw[popblue, thick] (4,0) ++(180:1.1) arc (180:143.13:1.1);
\node[popblue] at (3.0,0.35) {$\alpha$};
% Labels
\node[below, popdark] at (2,0) {\textbf{côté adjacent} : $AB$};
\node[left, poporange, align=right] at (-0.1,1.5) {\textbf{côté opposé} : $AC$};
\node[above right, poppurple, rotate=-36.87] at (2,1.5) {\textbf{hypoténuse} : $BC$};
\node[below left, popdark] at (A) {$A$};
\node[below right, popdark] at (B) {$B$};
\node[above left, popdark] at (C) {$C$};
\end{tikzpicture}
\end{center}
\legende{Triangle $ABC$ rectangle en $A$ : les trois côtés vus depuis l'angle aigu $\alpha = \widehat{ABC}$.}
\end{popfigure}

\begin{popfigure}
\begin{center}
\begin{tikzpicture}[scale=1.35, line cap=round, line join=round]
\coordinate (A) at (0,0);
\coordinate (B) at (4,0);
\coordinate (C) at (0,3);
\draw[thick, popdark] (A) -- (B) -- (C) -- cycle;
\draw[popdark] (0.3,0) -- (0.3,0.3) -- (0,0.3);
\draw[popblue, thick] (4,0) ++(180:1.1) arc (180:143.13:1.1);
\node[popblue] at (3.0,0.35) {$\alpha$};
\node[below, popblue] at (2,0) {adjacent};
\node[left, poporange] at (-0.1,1.5) {opposé};
\node[above right, poppurple, rotate=-36.87] at (2,1.5) {hypoténuse};
\node[popblue, align=left] at (-1.5,3.5) {$\cos(\alpha) = \dfrac{\text{adjacent}}{\text{hypoténuse}}$};
\node[poporange, align=left] at (-1.5,2.8) {$\sin(\alpha) = \dfrac{\text{opposé}}{\text{hypoténuse}}$};
\node[poppurple, align=left] at (-1.5,2.1) {$\tan(\alpha) = \dfrac{\text{opposé}}{\text{adjacent}}$};
\end{tikzpicture}
\end{center}
\legende{Les trois formules de trigonométrie rattachées aux côtés du triangle rectangle.}
\end{popfigure}

\courssub{Le moyen mnémotechnique : CAH--SOH--TOA}

Retenir trois formules peut sembler difficile. Heureusement, il existe une phrase-clé célèbre dans le monde entier : \textbf{CAH -- SOH -- TOA}. Chaque syllabe résume une formule (la lettre du milieu est le numérateur, la dernière le dénominateur) :

\begin{aretenir}[CAH -- SOH -- TOA]
\begin{center}
\begin{tabularx}{\linewidth}{|c|X|X|X|}\hline
\rowcolor{popblueL} \textbf{Formule} & \textbf{CAH} & \textbf{SOH} & \textbf{TOA} \\ \hline
Rapport & $\cos(\alpha)$ & $\sin(\alpha)$ & $\tan(\alpha)$ \\ \hline
Numérateur & \textbf{A}djacent & \textbf{O}pposé & \textbf{O}pposé \\ \hline
Dénominateur & \textbf{H}ypoténuse & \textbf{H}ypoténuse & \textbf{A}djacent \\ \hline
\end{tabularx}
\end{center}
On retient : \textbf{C}osinus = \textbf{A}djacent / \textbf{H}ypoténuse ; \textbf{S}inus = \textbf{O}pposé / \textbf{H}ypoténuse ; \textbf{T}angente = \textbf{O}pposé / \textbf{A}djacent. Certains élèves mémorisent la phrase : \og \emph{Certains Aiment Histoire, Sans Oublier Histoire, Toutefois On Aime} \fg{}... invente la tienne !
\end{aretenir}

\begin{exemplebox}[Lecture directe des trois rapports]
Soit $MNP$ un triangle rectangle en $M$ tel que $MN = 6$ cm, $MP = 8$ cm et $NP = 10$ cm. On considère l'angle $\alpha = \widehat{MNP}$.

Par rapport à $\alpha$ : l'hypoténuse est $NP = 10$ ; le côté adjacent est $MN = 6$ ; le côté opposé est $MP = 8$. Donc :
\[ \cos(\alpha) = \dfrac{MN}{NP} = \dfrac{6}{10} = 0,6 \quad ; \quad \sin(\alpha) = \dfrac{MP}{NP} = \dfrac{8}{10} = 0,8 \quad ; \quad \tan(\alpha) = \dfrac{MP}{MN} = \dfrac{8}{6} = \dfrac{4}{3} \]
Remarque : ici les longueurs sont exactes, donc les rapports sont des valeurs \textbf{exactes} et on utilise le signe $=$.
\end{exemplebox}

\courssub{Un rapport qui ne dépend que de l'angle}

Voici la propriété qui donne tout son sens à la trigonométrie. Traçons deux triangles rectangles $ABC$ et $AB'C'$ ayant le même angle $\alpha$ en $B$, l'un deux fois plus grand que l'autre. Les côtés sont différents, mais les \emph{rapports} sont identiques : si $B'C' = 2\,BC$ et $AC' = 2\,AC$, alors $\dfrac{A'C'}{B'C'} = \dfrac{2\,AC}{2\,BC} = \dfrac{AC}{BC}$. Le sinus est le même !

\begin{propriete}[Théorème admis — invariance des rapports trigonométriques]
La valeur du cosinus, du sinus ou de la tangente d'un angle aigu \textbf{ne dépend que de la mesure de cet angle}, et non des dimensions du triangle rectangle considéré.

Autrement dit : deux triangles rectangles ayant un angle aigu de même mesure $\alpha$ donnent exactement les mêmes valeurs de $\cos(\alpha)$, $\sin(\alpha)$ et $\tan(\alpha)$. Ce résultat est une conséquence du théorème de Thalès (les triangles sont semblables) ; il est \textbf{admis} en classe de 3ème.
\end{propriete}

\begin{savaistu}[D'où viennent ces noms ?]
Le mot \emph{sinus} vient d'une longue histoire : les mathématiciens indiens utilisaient la \og demi-corde \fg{} d'un arc de cercle ; leur terme fut traduit en arabe par \emph{jayb} (qui signifie \og pli \fg{} ou \og poche \fg{}), puis ce mot fut traduit en latin par \emph{sinus} (pli, courbe). Le \emph{cosinus} est le \og sinus de l'angle complémentaire \fg{} (\emph{complementi sinus} en latin) : en effet, dans un triangle rectangle, le cosinus d'un angle aigu est le sinus de l'autre angle aigu, car les deux angles aigus sont complémentaires (leur somme vaut $90°$). Enfin, \emph{tangente} vient du latin \emph{tangere}, \og toucher \fg{} : ce nom est lié à la droite tangente au cercle trigonométrique, que tu découvriras au lycée.
\end{savaistu}

\courssub{Domaine de variation des trois rapports}

Dans un triangle rectangle, l'hypoténuse est toujours le \textbf{plus long} des trois côtés. Par conséquent, pour un angle aigu $\alpha$ :
\begin{itemize}
\item le côté adjacent est plus petit que l'hypoténuse, donc $\cos(\alpha) = \dfrac{\text{adjacent}}{\text{hypoténuse}}$ est un nombre compris entre $0$ et $1$ ;
\item le côté opposé est plus petit que l'hypoténuse, donc $\sin(\alpha)$ est aussi compris entre $0$ et $1$ ;
\item en revanche, le côté opposé peut être \emph{plus grand} que le côté adjacent (pense à une échelle presque verticale !), donc $\tan(\alpha)$ peut dépasser $1$.
\end{itemize}

\begin{aretenir}[Encadrement des rapports trigonométriques]
Pour tout angle aigu $\alpha$ (c'est-à-dire $0° < \alpha < 90°$) :
\[ 0 < \cos(\alpha) < 1 \qquad ; \qquad 0 < \sin(\alpha) < 1 \qquad ; \qquad \tan(\alpha) > 0 \]
La tangente n'est \textbf{pas} bornée par $1$ : plus l'angle se rapproche de $90°$, plus sa tangente devient grande.
\end{aretenir}

\begin{attention}[Un contrôle rapide pour détecter les erreurs]
Si tu trouves, dans un calcul, $\cos(\alpha) = 1,4$ ou $\sin(\alpha) = -0,3$, c'est \textbf{forcément faux} : un cosinus ou un sinus d'angle aigu est toujours compris entre $0$ et $1$. En revanche, $\tan(\alpha) = 2,5$ est tout à fait possible. Utilise ce réflexe pour vérifier tes résultats au BEPC.
\end{attention}

\courssub{La relation fondamentale : tangente = sinus sur cosinus}

Les trois rapports ne sont pas indépendants : la tangente s'exprime à l'aide du sinus et du cosinus. La démonstration est un simple calcul littéral sur les fractions, à ta portée.

\begin{propriete}[Relation entre tangente, sinus et cosinus]
Pour tout angle aigu $\alpha$ :
\[ \tan(\alpha) = \dfrac{\sin(\alpha)}{\cos(\alpha)} \]
\end{propriete}

\begin{experience}[Démonstration guidée de $\tan(\alpha) = \dfrac{\sin(\alpha)}{\cos(\alpha)}$]
On se place dans un triangle $ABC$ rectangle en $A$, avec $\alpha = \widehat{ABC}$. Notons pour simplifier : $o = AC$ (côté opposé), $a = AB$ (côté adjacent), $h = BC$ (hypoténuse).

\textbf{Étape 1.} Écris les définitions : $\sin(\alpha) = \dfrac{o}{h}$ et $\cos(\alpha) = \dfrac{a}{h}$.

\textbf{Étape 2.} Forme le quotient $\dfrac{\sin(\alpha)}{\cos(\alpha)}$ :
\[ \dfrac{\sin(\alpha)}{\cos(\alpha)} = \dfrac{\dfrac{o}{h}}{\dfrac{a}{h}} \]

\textbf{Étape 3.} Diviser par une fraction, c'est multiplier par son inverse :
\[ \dfrac{\dfrac{o}{h}}{\dfrac{a}{h}} = \dfrac{o}{h} \times \dfrac{h}{a} = \dfrac{o \times h}{h \times a} = \dfrac{o}{a} \]
(le facteur $h$, non nul, se simplifie).

\textbf{Étape 4.} Or $\dfrac{o}{a} = \dfrac{\text{opposé}}{\text{adjacent}} = \tan(\alpha)$. Conclusion :
\[ \dfrac{\sin(\alpha)}{\cos(\alpha)} = \tan(\alpha) \qquad \text{CQFD.} \]
\end{experience}

\begin{exemplebox}[Vérification numérique]
Reprenons le triangle $MNP$ de l'exemple précédent : $\sin(\alpha) = 0,8$ et $\cos(\alpha) = 0,6$. Alors :
\[ \dfrac{\sin(\alpha)}{\cos(\alpha)} = \dfrac{0,8}{0,6} = \dfrac{8}{6} = \dfrac{4}{3} = \tan(\alpha) \]
On retrouve bien la valeur de la tangente calculée directement. La relation est vérifiée.
\end{exemplebox}

\courssub{Utilisation de la calculatrice}

En pratique, on ne connaît pas les longueurs des côtés : c'est la calculatrice qui fournit les valeurs des rapports trigonométriques à partir de la mesure de l'angle. Mais attention : elle doit être réglée dans la bonne unité d'angle.

\begin{methode}[Calculer un rapport trigonométrique à la calculatrice]
Pour calculer par exemple $\cos(35°)$ :
\begin{enumerate}
\item \textbf{Vérifier le mode} : la calculatrice doit être en mode \textbf{DEG} (degrés). Cherche l'indication \og D \fg{} ou \og DEG \fg{} en haut de l'écran. Si tu vois \og RAD \fg{} ou \og G \fg{}, change le mode (touche \textsc{mode} ou \textsc{setup}).
\item Appuyer sur la touche du rapport voulu : \textsc{cos}, \textsc{sin} ou \textsc{tan}.
\item Taper la mesure de l'angle, par exemple $35$, puis fermer la parenthèse si elle s'est ouverte.
\item Appuyer sur \textsc{=} ou \textsc{EXE}.
\item \textbf{Arrondir} selon la consigne : en général au millième ($10^{-3}$) ou au dix-millième ($10^{-4}$).
\end{enumerate}
Exemple : $\cos(35°) \approx 0,8192$ (arrondi à $10^{-4}$) ; $\sin(35°) \approx 0,5736$ ; $\tan(35°) \approx 0,7002$.
\end{methode}

\begin{attention}[Rédaction : parenthèses, degrés et symbole $\approx$]
Trois règles de rédaction sont exigées au BEPC :
\begin{itemize}
\item Écris toujours l'angle \textbf{entre parenthèses} et \textbf{avec le symbole degré} : $\cos(30°)$, et jamais \og cos 30 \fg{}.
\item La calculatrice donne une valeur \textbf{approchée} : utilise le symbole $\approx$, pas le signe $=$. On écrit $\cos(30°) \approx 0,866$, et non $\cos(30°) = 0,86$.
\item Respecte l'arrondi demandé par l'énoncé ($10^{-2}$, $10^{-3}$ ou $10^{-4}$) et ne donne pas tous les chiffres de l'écran.
\end{itemize}
\end{attention}

\begin{exoresolu}[Lecture de rapports et vérification à la calculatrice]
\textbf{Énoncé.} Soit $RST$ un triangle rectangle en $R$ tel que $RS = 5$ cm et $RT = 12$ cm. On note $\beta = \widehat{RST}$.

1) Calculer la longueur $ST$.
2) Donner les valeurs exactes de $\cos(\beta)$, $\sin(\beta)$ et $\tan(\beta)$.
3) À l'aide de la calculatrice, vérifier que $\dfrac{\sin(\beta)}{\cos(\beta)}$ redonne bien $\tan(\beta)$, avec des valeurs arrondies à $10^{-4}$.
\tcblower
\textbf{Solution détaillée.}

\textbf{1)} Le triangle $RST$ est rectangle en $R$, donc d'après le théorème de Pythagore :
\begin{align*}
ST^2 &= RS^2 + RT^2 \\
ST^2 &= 5^2 + 12^2 = 25 + 144 = 169 \\
ST &= \sqrt{169} = 13 \text{ cm}
\end{align*}

\textbf{2)} Par rapport à l'angle $\beta = \widehat{RST}$ : l'hypoténuse est $ST = 13$ ; le côté adjacent est $RS = 5$ ; le côté opposé est $RT = 12$. D'où :
\[ \cos(\beta) = \dfrac{RS}{ST} = \dfrac{5}{13} \qquad ; \qquad \sin(\beta) = \dfrac{RT}{ST} = \dfrac{12}{13} \qquad ; \qquad \tan(\beta) = \dfrac{RT}{RS} = \dfrac{12}{5} \]

\textbf{3)} Valeurs approchées à $10^{-4}$ : $\sin(\beta) = \dfrac{12}{13} \approx 0,9231$ et $\cos(\beta) = \dfrac{5}{13} \approx 0,3846$. Alors :
\[ \dfrac{\sin(\beta)}{\cos(\beta)} \approx \dfrac{0,9231}{0,3846} \approx 2,4000 \qquad \text{et} \qquad \tan(\beta) = \dfrac{12}{5} = 2,4 \]
La relation $\tan(\beta) = \dfrac{\sin(\beta)}{\cos(\beta)}$ est bien vérifiée.
\end{exoresolu}

\begin{exercice}[À toi de jouer]
Soit $EFG$ un triangle rectangle en $E$ tel que $EF = 9$ cm et $EG = 12$ cm. On note $\gamma = \widehat{EFG}$.
\begin{enumerate}
\item Calculer $FG$.
\item Donner les valeurs exactes de $\cos(\gamma)$, $\sin(\gamma)$ et $\tan(\gamma)$.
\item Vérifier que $\cos(\gamma)$ et $\sin(\gamma)$ sont bien compris entre $0$ et $1$, puis vérifier la relation $\tan(\gamma) = \dfrac{\sin(\gamma)}{\cos(\gamma)}$.
\item À l'aide de la calculatrice (mode DEG), donner $\cos(50°)$, $\sin(50°)$ et $\tan(50°)$ arrondis à $10^{-3}$, puis vérifier la relation fondamentale pour l'angle $50°$.
\end{enumerate}
\end{exercice}

\begin{corrige}[Exercice : À toi de jouer]
\textbf{1)} D'après le théorème de Pythagore : $FG^2 = EF^2 + EG^2 = 81 + 144 = 225$, donc $FG = 15$ cm.

\textbf{2)} Par rapport à $\gamma$ : hypoténuse $FG = 15$, adjacent $EF = 9$, opposé $EG = 12$.
\[ \cos(\gamma) = \dfrac{9}{15} = \dfrac{3}{5} \quad ; \quad \sin(\gamma) = \dfrac{12}{15} = \dfrac{4}{5} \quad ; \quad \tan(\gamma) = \dfrac{12}{9} = \dfrac{4}{3} \]

\textbf{3)} $\cos(\gamma) = 0,6$ et $\sin(\gamma) = 0,8$ : tous deux sont bien compris entre $0$ et $1$. De plus $\dfrac{\sin(\gamma)}{\cos(\gamma)} = \dfrac{0,8}{0,6} = \dfrac{4}{3} = \tan(\gamma)$ : la relation est vérifiée.

\textbf{4)} En mode DEG : $\cos(50°) \approx 0,643$ ; $\sin(50°) \approx 0,766$ ; $\tan(50°) \approx 1,192$. Vérification : $\dfrac{0,766}{0,643} \approx 1,191 \approx \tan(50°)$ (aux arrondis près). On note au passage que $\tan(50°) > 1$, ce qui est possible car l'angle dépasse $45°$.
\end{corrige}

\begin{aretenir}[L'essentiel de la section]
\begin{itemize}
\item Dans un triangle rectangle, pour un angle aigu $\alpha$ : $\cos(\alpha) = \dfrac{\text{adjacent}}{\text{hypoténuse}}$, $\sin(\alpha) = \dfrac{\text{opposé}}{\text{hypoténuse}}$, $\tan(\alpha) = \dfrac{\text{opposé}}{\text{adjacent}}$ (moyen mnémotechnique : \textbf{CAH--SOH--TOA}).
\item Ces rapports ne dépendent \textbf{que de la mesure de l'angle}, pas de la taille du triangle.
\item Pour un angle aigu : $0 < \cos(\alpha) < 1$, $0 < \sin(\alpha) < 1$ et $\tan(\alpha) > 0$.
\item Relation fondamentale : $\tan(\alpha) = \dfrac{\sin(\alpha)}{\cos(\alpha)}$.
\item Calculatrice : mode \textbf{DEG}, écriture $\cos(30°) \approx 0,866$ avec parenthèses, degré et symbole $\approx$.
\end{itemize}
\end{aretenir}

Maintenant que les trois rapports sont définis et que tu sais les lire sur ta calculatrice, une question naturelle se pose : existe-t-il des angles dont les valeurs trigonométriques sont \emph{exactes} et remarquables ? C'est l'objet de la section suivante, consacrée aux angles de $30°$, $45°$ et $60°$.

\coursec{Angles remarquables (30\textsuperscript{ème}, 45\textsuperscript{ème}, 60\textsuperscript{ème}) et valeurs exactes}

Dans la section précédente, nous avons appris à utiliser les rapports \cle{cosinus}, \cle{sinus} et \cle{tangente} pour calculer des longueurs ou des angles dans un triangle rectangle, à l'aide d'une calculatrice. Mais il existe des angles particuliers — les \textbf{angles remarquables} — pour lesquels ces rapports prennent des \textbf{valeurs exactes} que l'on peut écrire sans approximation décimale, uniquement avec des entiers et des racines carrées. Connaître ces valeurs par cœur (ou savoir les retrouver instantanément) est un atout majeur : cela permet de travailler \emph{sans calculatrice}, de simplifier des expressions littérales, et c'est souvent exigé aux examens (BEPC, Probatoire) avec la mention \og valeur exacte \fg{} ou \og forme simplifiée \fg{}.

\courssub{Le triangle rectangle isocèle : l'angle de 45\textsuperscript{ème}}

Commençons par le cas le plus simple. Imagine un triangle rectangle qui est aussi \emph{isocèle}. Ses deux angles aigus sont donc égaux ; comme leur somme vaut 90\textsuperscript{ème}, chacun mesure \textbf{45\textsuperscript{ème}}.

\begin{popfigure}
\begin{tikzpicture}[scale=2.5, >=Stealth]
% Triangle ABC rectangle en A, AB=AC=1
\coordinate (A) at (0,0);
\coordinate (B) at (1,0);
\coordinate (C) at (0,1);
\draw[popblue, thick] (A) -- node[below] {$1$} (B) -- node[right, sloped] {$\sqrt{2}$} (C) -- node[left] {$1$} (A) -- cycle;
% Angle droit
\draw[popdark] (0.1,0) -- (0.1,0.1) -- (0,0.1);
% Angles 45°
\draw[poporange, thick] (1,0) ++(-0.25,0) arc (180:135:0.25) node[midway, left=2pt, poporange] {$45^\circ$};
\draw[poporange, thick] (0,1) ++(0,-0.25) arc (-90:-45:0.25) node[midway, below=2pt, poporange] {$45^\circ$};
% Noms des sommets
\node[below left, popdark] at (A) {$A$};
\node[below right, popdark] at (B) {$B$};
\node[above left, popdark] at (C) {$C$};
% Angle droit label
\node[above right, popdark] at (0.1,0.1) {$90^\circ$};
\end{tikzpicture}
\legende{Triangle rectangle isocèle $ABC$ rectangle en $A$, $AB=AC=1$, $BC=\sqrt{2}$ (Théorème de Pythagore).}
\end{popfigure}

Appliquons les définitions \cle{CAH-SOH-TOA} à l'angle $\widehat{B}=45^\circ$ (côté adjacent $AB=1$, côté opposé $AC=1$, hypoténuse $BC=\sqrt{2}$) :
\[
\cos 45^\circ = \frac{AB}{BC} = \frac{1}{\sqrt{2}} = \frac{\sqrt{2}}{2},\quad
\sin 45^\circ = \frac{AC}{BC} = \frac{1}{\sqrt{2}} = \frac{\sqrt{2}}{2},\quad
\tan 45^\circ = \frac{AC}{AB} = \frac{1}{1} = 1.
\]

\begin{attention}[Forme rationnalisée exigée]
Au BEPC et au Probatoire, on \textbf{ne laisse jamais} de racine carrée au dénominateur.
\begin{itemize}
    \item On écrit $\dfrac{\sqrt{2}}{2}$ et \textbf{pas} $\dfrac{1}{\sqrt{2}}$.
    \item On écrit $\dfrac{\sqrt{3}}{3}$ et \textbf{pas} $\dfrac{1}{\sqrt{3}}$.
\end{itemize}
C'est une convention mathématique universelle : on \textbf{rationnalise le dénominateur}.
\end{attention}

\courssub{Le triangle rectangle demi-équilatéral : les angles de 30\textsuperscript{ème} et 60\textsuperscript{ème}}

Prenons maintenant un triangle \textbf{équilatéral} de côté $2$. Tracez sa hauteur. Elle coupe le triangle en deux triangles rectangles identiques. Dans chacun de ces triangles :
\begin{itemize}
    \item L'hypoténuse vaut $2$ (côté de l'équilatéral).
    \item Le côté adjacent à l'angle de $60^\circ$ (demi-base) vaut $1$.
    \item L'angle au sommet de l'équilatéral est coupé en deux : $60^\circ \to 30^\circ$.
    \item Le troisième angle vaut $180^\circ - 90^\circ - 30^\circ = 60^\circ$ (ou $180^\circ - 90^\circ - 60^\circ = 30^\circ$).
\end{itemize}
Le troisième côté (la hauteur) se calcule par Pythagore : $h^2 = 2^2 - 1^2 = 3 \implies h = \sqrt{3}$.

\begin{popfigure}
\begin{tikzpicture}[scale=2.5, >=Stealth]
% Triangle équilatéral complet en pointillés (côté 2)
\coordinate (O) at (0,0);
\coordinate (P) at (2,0);
\coordinate (Q) at (1,1.732); % sqrt(3)
\draw[popmuted, dashed, thick] (O) -- (P) -- (Q) -- cycle;
\node[popmuted, above] at (Q) {Triangle équilatéral de côté 2};
% Triangle rectangle demi-équilatéral (gauche) en plein
\coordinate (A) at (0,0);
\coordinate (B) at (1,0);
\coordinate (C) at (1,1.732);
\draw[popblue, very thick] (A) -- node[below] {$1$} (B) -- node[right, sloped] {$2$} (C) -- node[left, sloped] {$\sqrt{3}$} (A) -- cycle;
% Angle droit en B
\draw[popdark] (0.9,0) -- (0.9,0.1) -- (1,0.1);
% Angles
\draw[poporange, thick] (0,0) ++(0.3,0) arc (0:60:0.3) node[midway, right=2pt, poporange] {$60^\circ$};
\draw[poporange, thick] (1,1.732) ++(-0.3,-0.52) arc (-120:-90:0.3) node[midway, left=2pt, poporange] {$30^\circ$};
% Noms
\node[below left, popdark] at (A) {$A$};
\node[below right, popdark] at (B) {$B$};
\node[above right, popdark] at (C) {$C$};
% Cotes du grand triangle
\node[below, popmuted] at (1,0) {$2$};
\node[left, popmuted] at (0.5,0.866) {$2$};
\end{tikzpicture}
\legende{Construction des angles 30\textsuperscript{ème} et 60\textsuperscript{ème} par la hauteur d'un triangle équilatéral de côté 2. Le triangle rectangle $ABC$ (rectangle en $B$) a pour côtés $1$, $\sqrt{3}$, $2$.}
\end{popfigure}

Calculons les rapports pour l'angle de \textbf{60\textsuperscript{ème}} (sommet $A$ : adjacent $=1$, opposé $=\sqrt{3}$, hyp $=2$) :
\[
\cos 60^\circ = \frac{1}{2},\quad \sin 60^\circ = \frac{\sqrt{3}}{2},\quad \tan 60^\circ = \frac{\sqrt{3}}{1} = \sqrt{3}.
\]

Calculons les rapports pour l'angle de \textbf{30\textsuperscript{ème}} (sommet $C$ : adjacent $=\sqrt{3}$, opposé $=1$, hyp $=2$) :
\[
\cos 30^\circ = \frac{\sqrt{3}}{2},\quad \sin 30^\circ = \frac{1}{2},\quad \tan 30^\circ = \frac{1}{\sqrt{3}} = \frac{\sqrt{3}}{3}.
\]

\begin{propriete}[Valeurs exactes des rapports trigonométriques des angles remarquables]
Voici le tableau \textbf{à connaître par cœur} (valeurs exactes, dénominateurs rationnalisés) :

\begin{center}
\begin{tabular}{|c|c|c|c|}\hline
\rowcolor{popblueL}
\textbf{Angle} & \textbf{$\boldsymbol{\cos}$} & \textbf{$\boldsymbol{\sin}$} & \textbf{$\boldsymbol{\tan}$} \\ \hline
$30^\circ$ & $\dfrac{\sqrt{3}}{2}$ & $\dfrac{1}{2}$ & $\dfrac{\sqrt{3}}{3}$ \\ \hline
$45^\circ$ & $\dfrac{\sqrt{2}}{2}$ & $\dfrac{\sqrt{2}}{2}$ & $1$ \\ \hline
$60^\circ$ & $\dfrac{1}{2}$ & $\dfrac{\sqrt{3}}{2}$ & $\sqrt{3}$ \\ \hline
\end{tabular}
\end{center}
\end{propriete}

\begin{aretenir}[Relations complémentaires (Astuce mémoire)]
Les angles $30^\circ$ et $60^\circ$ sont \textbf{complémentaires} ($30^\circ+60^\circ=90^\circ$). Dans un triangle rectangle, les deux angles aigus sont complémentaires. Donc :
\begin{itemize}
    \item $\cos 30^\circ = \sin 60^\circ = \dfrac{\sqrt{3}}{2}$ \quad (le côté adjacent à 30\textsuperscript{ème} est opposé à 60\textsuperscript{ème})
    \item $\cos 60^\circ = \sin 30^\circ = \dfrac{1}{2}$
    \item $\tan 30^\circ = \dfrac{1}{\tan 60^\circ} = \dfrac{\sqrt{3}}{3}$
\end{itemize}
\emph{Pour 45\textsuperscript{ème}, les deux angles aigus sont égaux, d'où $\cos 45^\circ = \sin 45^\circ$ et $\tan 45^\circ = 1$.}
\end{aretenir}

\courssub{Méthode pour retrouver les valeurs sans les apprendre par cœur}

\begin{methode}[Reconstruire le tableau en 30 secondes (Technique du prof)]
Si tu as un trou de mémoire le jour de l'examen, ne panique pas : \textbf{redessine les deux triangles de base}.
\begin{enumerate}
    \item \textbf{Pour 45\textsuperscript{ème} :} Dessine un triangle rectangle isocèle. Mets $1$ et $1$ sur les côtés de l'angle droit. L'hypoténuse vaut $\sqrt{2}$ (Pythagore). Applique CAH-SOH-TOA $\to$ tu retrouves $\frac{\sqrt{2}}{2}$ et $1$.
    \item \textbf{Pour 30\textsuperscript{ème} et 60\textsuperscript{ème} :} Dessine un triangle équilatéral de côté $2$. Trace la hauteur. Tu obtiens un triangle rectangle de côtés $1$, $\sqrt{3}$, $2$. Identifie où sont 30\textsuperscript{ème} et 60\textsuperscript{ème} (le petit angle est face au petit côté $1$). Applique CAH-SOH-TOA.
    \item \textbf{Rationnalise :} $\frac{1}{\sqrt{2}} \to \frac{\sqrt{2}}{2}$ ; $\frac{1}{\sqrt{3}} \to \frac{\sqrt{3}}{3}$.
\end{enumerate}
Cette méthode est \textbf{infaillible} et montre au correcteur que tu comprends la géométrie derrière les formules.
\end{methode}

\courssub{Applications : Calculer des longueurs exactes et simplifier des expressions}

L'intérêt principal de ces valeurs exactes est de pouvoir donner des résultats \textbf{sans calculatrice}, sous forme de nombres exacts (entiers, fractions, racines carrées). C'est la norme en classe de 3\textsuperscript{ème} et au-delà.

\begin{exemplebox}[Calcul de longueurs en valeur exacte]
Soit $ABC$ un triangle rectangle en $A$ tel que $\widehat{B} = 30^\circ$ et $AB = \SI{5}{\centi\meter}$.
Calculer $AC$ et $BC$ en \textbf{valeur exacte}.

\medskip
\textbf{Solution.}
On se place par rapport à l'angle $30^\circ$ en $B$.
\begin{itemize}
    \item Côté adjacent : $AB = \SI{5}{\centi\meter}$.
    \item Côté opposé : $AC$ (inconnu).
    \item Hypoténuse : $BC$ (inconnu).
\end{itemize}

\noindent\textbf{1. Calcul de $AC$ (côté opposé) :} On utilise la tangente.
\[
\tan 30^\circ = \frac{AC}{AB} \iff AC = AB \times \tan 30^\circ = 5 \times \frac{\sqrt{3}}{3} = \frac{5\sqrt{3}}{3}\,\si{\centi\meter}.
\]

\noindent\textbf{2. Calcul de $BC$ (hypoténuse) :} On utilise le cosinus.
\[
\cos 30^\circ = \frac{AB}{BC} \iff BC = \frac{AB}{\cos 30^\circ} = \frac{5}{\frac{\sqrt{3}}{2}} = 5 \times \frac{2}{\sqrt{3}} = \frac{10}{\sqrt{3}} = \frac{10\sqrt{3}}{3}\,\si{\centi\meter}.
\]

\noindent\textbf{Vérification (Pythagore) :}
$AB^2 + AC^2 = 25 + \left(\frac{5\sqrt{3}}{3}\right)^2 = 25 + \frac{75}{9} = 25 + \frac{25}{3} = \frac{100}{3}$.
$BC^2 = \left(\frac{10\sqrt{3}}{3}\right)^2 = \frac{300}{9} = \frac{100}{3}$. \quad $\checkmark$

\medskip
\textbf{Résultat :} $AC = \dfrac{5\sqrt{3}}{3}\,\si{\centi\meter}$ et $BC = \dfrac{10\sqrt{3}}{3}\,\si{\centi\meter}$.
\end{exemplebox}

\begin{exoresolu}[Simplification d'expression trigonométrique]
Simplifier l'expression suivante en valeur exacte :
\[
E = \frac{\sin 60^\circ \times \cos 30^\circ}{\tan 45^\circ} - \sin^2 30^\circ.
\]

\tcblower
\textbf{Solution détaillée.}
On remplace chaque rapport par sa valeur exacte du tableau :
\begin{align*}
\sin 60^\circ &= \frac{\sqrt{3}}{2}, \quad
\cos 30^\circ = \frac{\sqrt{3}}{2}, \quad
\tan 45^\circ = 1, \quad
\sin 30^\circ = \frac{1}{2}.
\end{align*}
On substitue :
\[
E = \frac{\frac{\sqrt{3}}{2} \times \frac{\sqrt{3}}{2}}{1} - \left(\frac{1}{2}\right)^2
= \frac{\frac{3}{4}}{1} - \frac{1}{4}
= \frac{3}{4} - \frac{1}{4}
= \frac{2}{4}
= \frac{1}{2}.
\]
\textbf{Résultat :} $E = \dfrac{1}{2}$.
\end{exoresolu}

\courssub{Problème concret : L'ombre du poteau électrique (Contexte Cameroun)}

\begin{exoresolu}[Hauteur d'un poteau CAMTEL]
À Yaoundé, un poteau électrique CAMTEL est planté verticalement sur un sol horizontal. À un certain moment de la journée, l'ombre du poteau mesure $4\sqrt{3}\,\si{\meter}$ et le soleil forme un angle de $60^\circ$ avec le sol. Quelle est la hauteur exacte du poteau ?

\tcblower
\textbf{Solution.}
Modélisation : Triangle rectangle $POT$ rectangle en $O$ (pied du poteau).
\begin{itemize}
    \item $PO$ = hauteur du poteau (inconnue).
    \item $OT = 4\sqrt{3}\,\si{\meter}$ (longueur de l'ombre, côté adjacent à l'angle de $60^\circ$).
    \item $\widehat{T} = 60^\circ$ (angle du soleil au sol).
\end{itemize}
On cherche $PO$ (côté opposé à l'angle $60^\circ$). On utilise la tangente :
\[
\tan 60^\circ = \frac{PO}{OT} \iff PO = OT \times \tan 60^\circ = 4\sqrt{3} \times \sqrt{3} = 4 \times 3 = \SI{12}{\meter}.
\]
\textbf{Le poteau mesure exactement 12 mètres.}
\end{exoresolu}

\begin{savaistu}[Pourquoi "valeur exacte" aux examens ?]
Au BEPC et au Probatoire, les sujets tombent souvent sur des angles remarquables ($30^\circ$, $45^\circ$, $60^\circ$). La consigne \og \textbf{Donner la valeur exacte} \fg{} ou \og \textbf{Écrire sous la forme $a\frac{\sqrt{b}}{c}$} \fg{} signifie : \textbf{interdit d'utiliser la touche \texttt{sin}/\texttt{cos}/\texttt{tan} de la calculatrice} pour donner un nombre décimal comme $0,866...$ ou $1,732...$.
Le correcteur attend $\frac{\sqrt{3}}{2}$ ou $\sqrt{3}$.
\begin{itemize}
    \item \textbf{Gain de temps :} Pas besoin de taper sur la machine, tu écris direct.
    \item \textbf{Zéro erreur d'arrondi :} $0,8660254...$ n'est pas $\frac{\sqrt{3}}{2}$.
    \item \textbf{Simplification algébrique :} Comme dans l'exercice résolu ci-dessus, les racines se simplifient entre elles ($\sqrt{3} \times \sqrt{3} = 3$), ce qui est impossible avec des décimaux.
\end{itemize}
Maîtriser ce tableau, c'est s'assurer \textbf{4 à 6 points faciles} sur une épreuve de trigonométrie.
\end{savaistu}

\courssub{Exercices d'entraînement (Niveau BEPC)}

\begin{exercice}[Exercice 1 : Application directe]
Dans un triangle rectangle $MNP$ rectangle en $N$, on sait que $\widehat{M} = 60^\circ$ et $MN = \SI{6}{\centi\meter}$.
Calculer en valeur exacte les longueurs $NP$ et $MP$.
\end{exercice}

\begin{exercice}[Exercice 2 : Simplification]
Calculer la valeur exacte de :
\[
F = 2\cos^2 30^\circ - \tan 45^\circ + \frac{\sin 60^\circ}{\cos 30^\circ}.
\]
\end{exercice}

\begin{exercice}[Exercice 3 : Problème géométrique]
$ABC$ est un triangle rectangle en $A$ tel que $AB = \SI{4}{\centi\meter}$ et $\widehat{B} = 45^\circ$.
Le point $H$ est le projeté orthogonal de $A$ sur $[BC]$ (la hauteur issue de $A$).
Calculer en valeur exacte $AH$, $BH$ et $HC$.
\end{exercice}

\begin{corrige}[Exercice 1]
Triangle $MNP$ rectangle en $N$, $\widehat{M}=60^\circ$, $MN=\SI{6}{\centi\meter}$ (adjacent à 60\textsuperscript{ème}).
\begin{itemize}
    \item $NP$ (opposé) : $\tan 60^\circ = \frac{NP}{6} \implies NP = 6 \times \sqrt{3} = 6\sqrt{3}\,\si{\centi\meter}$.
    \item $MP$ (hyp) : $\cos 60^\circ = \frac{6}{MP} \implies MP = \frac{6}{1/2} = \SI{12}{\centi\meter}$.
\end{itemize}
\end{corrige}

\begin{corrige}[Exercice 2]
$\cos 30^\circ = \frac{\sqrt{3}}{2} \implies \cos^2 30^\circ = \frac{3}{4}$.
$\tan 45^\circ = 1$.
$\frac{\sin 60^\circ}{\cos 30^\circ} = \frac{\frac{\sqrt{3}}{2}}{\frac{\sqrt{3}}{2}} = 1$.
$F = 2 \times \frac{3}{4} - 1 + 1 = \frac{3}{2}$.
\end{corrige}

\begin{corrige}[Exercice 3]
$ABC$ rectangle en $A$, $\widehat{B}=45^\circ \implies \widehat{C}=45^\circ$ (triangle rectangle isocèle).
$AB = AC = \SI{4}{\centi\meter}$.
$BC = 4\sqrt{2}\,\si{\centi\meter}$ (Pythagore ou $\frac{4}{\cos 45^\circ}$).
$AH$ est hauteur dans un triangle rectangle isocèle $\implies H$ est milieu de $[BC]$.
$BH = HC = \frac{BC}{2} = 2\sqrt{2}\,\si{\centi\meter}$.
$AH = BH = 2\sqrt{2}\,\si{\centi\meter}$ (triangle $AHB$ rectangle isocèle en $H$ car $\widehat{B}=45^\circ$).
\end{corrige}

\coursec{Calculer une longueur dans un triangle rectangle}

Après avoir maîtrisé les valeurs exactes des angles remarquables (30°, 45°, 60°), nous disposons maintenant de tous les outils pour \emph{résoudre} un triangle rectangle : connaître un angle aigu et un côté suffit à déterminer toutes les autres mesures. C'est la puissance de la trigonométrie : transformer un problème géométrique en un calcul algébrique simple. Dans cette section, nous apprenons la démarche rigoureuse pour \textbf{calculer une longueur inconnue} lorsque l'on connaît un angle aigu et un côté.

\courssub{Stratégie générale : de l'angle au côté}

Le point de départ est toujours le même : \textbf{un angle aigu connu} (appelons-le $\alpha$) et \textbf{un côté connu}. Le côté cherché peut être l'hypoténuse, le côté adjacent à $\alpha$ ou le côté opposé à $\alpha$. La clé est de choisir le \textbf{bon rapport trigonométrique} — celui qui \emph{relie} le côté connu et le côté cherché.

\begin{popfigure}
\centering
\begin{tikzpicture}[
    node distance=1.5cm and 2cm,
    decision/.style={diamond, draw=popblue, thick, fill=popblueL, text width=3.2cm, align=center, inner sep=4pt},
    process/.style={rectangle, draw=popdark, thick, fill=popcream, text width=3.2cm, align=center, inner sep=4pt, rounded corners=4pt},
    startstop/.style={ellipse, draw=popgreen, thick, fill=popgreenL, text width=2.8cm, align=center, inner sep=4pt},
    arrow/.style={-Stealth, thick, popdark},
    yesno/.style={midway, fill=white, inner sep=2pt, font=\small\bfseries}
]
\node[startstop] (start) {Je connais un angle aigu $\alpha$ et un côté};
\node[process, below=of start, align=center] (schema){1. Faire un schéma annoté\\Repérer Hyp, Adj, Opp par rapport à $\alpha$};
\node[decision, below=of schema, align=center] (choix){Quel côté est connu ?\\Quel côté est cherché ?};
\node[process, below left=of choix, align=center] (cas1){Côté connu = Hyp\\Cherché = Adj $\rightarrow$ $\cos\alpha$\\Cherché = Opp $\rightarrow$ $\sin\alpha$};
\node[process, below=of choix, align=center] (cas2){Côté connu = Adj\\Cherché = Hyp $\rightarrow$ $\cos\alpha$\\Cherché = Opp $\rightarrow$ $\tan\alpha$};
\node[process, below right=of choix, align=center] (cas3){Côté connu = Opp\\Cherché = Hyp $\rightarrow$ $\sin\alpha$\\Cherché = Adj $\rightarrow$ $\tan\alpha$};
\node[process, below=of cas2, align=center] (calcul){2. Écrire l'égalité\\3. Isoler l'inconnue (produit en croix)\\4. Calculer (calculatrice en mode DEG)\\5. Conclure avec unité};
\node[startstop, below=of calcul] (end) {Longueur trouvée !};

\draw[arrow] (start) -- (schema);
\draw[arrow] (schema) -- (choix);
\draw[arrow] (choix) -- node[yesno, left] {Hyp connu} (cas1);
\draw[arrow] (choix) -- node[yesno, above] {Adj connu} (cas2);
\draw[arrow] (choix) -- node[yesno, right] {Opp connu} (cas3);
\draw[arrow] (cas1) -- (calcul);
\draw[arrow] (cas2) -- (calcul);
\draw[arrow] (cas3) -- (calcul);
\draw[arrow] (calcul) -- (end);
\end{tikzpicture}
\legende{Arbre de décision : choisir le bon rapport trigonométrique selon le côté connu et le côté cherché.}
\end{popfigure}

\courssub{Méthode pas à pas}

\begin{methode}[Calculer une longueur dans un triangle rectangle]
\textbf{Étape 1 :} Faire un \textbf{schéma annoté} (même approximatif) du triangle rectangle. Indiquer l'angle aigu connu $\alpha$, le côté connu (avec sa valeur et son unité) et le côté cherché (avec une lettre, par exemple $x$).

\textbf{Étape 2 :} Choisir l'\textbf{angle aigu de référence} (celui dont on connaît la mesure). Nommer les trois côtés \emph{par rapport à cet angle} :
\begin{itemize}
    \item \textbf{Hyp} = hypoténuse (côté opposé à l'angle droit, toujours le plus grand) ;
    \item \textbf{Adj} = côté adjacent à $\alpha$ (celui qui « touche » $\alpha$ et n'est pas l'hypoténuse) ;
    \item \textbf{Opp} = côté opposé à $\alpha$ (celui qui ne touche pas $\alpha$).
\end{itemize}

\textbf{Étape 3 :} Choisir la formule qui contient \textbf{le côté connu et le côté cherché} :
\begin{align*}
    \cos\alpha &= \frac{\text{Adj}}{\text{Hyp}} & \sin\alpha &= \frac{\text{Opp}}{\text{Hyp}} & \tan\alpha &= \frac{\text{Opp}}{\text{Adj}}
\end{align*}

\textbf{Étape 4 :} Écrire l'égalité en remplaçant $\alpha$ par sa valeur, le côté connu par sa mesure, et le côté cherché par $x$.

\textbf{Étape 5 :} Isoler $x$ par un \textbf{produit en croix} (multiplier les deux membres par le dénominateur).

\textbf{Étape 6 :} Calculer sur la calculatrice (\textbf{mode DEGRÉS} impératif !) ou utiliser une valeur exacte si $\alpha \in \{30°,\,45°,\,60°\}$. Arrondir selon la consigne (souvent au millimètre ou au dixième de cm).

\textbf{Étape 7 :} Écrire la conclusion : « Le côté [...] mesure environ ... cm (ou m). »
\end{methode}

\courssub{Propriété fondamentale}

\begin{propriete}[Calcul des côtés d'un triangle rectangle]
Dans un triangle rectangle, si l'on connaît la mesure d'un angle aigu $\alpha$ ($0°<\alpha<90°$) et la longueur d'un côté, on peut calculer les deux autres côtés.
\begin{itemize}
    \item Si l'on connaît l'hypoténuse $H$ : $\quad \text{Adjacent} = H \times \cos\alpha \quad;\quad \text{Opposé} = H \times \sin\alpha$.
    \item Si l'on connaît le côté adjacent $A$ : $\quad \text{Hypoténuse} = \frac{A}{\cos\alpha} \quad;\quad \text{Opposé} = A \times \tan\alpha$.
    \item Si l'on connaît le côté opposé $O$ : $\quad \text{Hypoténuse} = \frac{O}{\sin\alpha} \quad;\quad \text{Adjacent} = \frac{O}{\tan\alpha}$.
\end{itemize}
\emph{Preuve :} Cela découle directement des définitions $\cos\alpha=\frac{\text{Adj}}{\text{Hyp}}$, $\sin\alpha=\frac{\text{Opp}}{\text{Hyp}}$, $\tan\alpha=\frac{\text{Opp}}{\text{Adj}}$ en isolant chaque côté.
\end{propriete}

\courssub{Exemple détaillé : l'échelle contre le mur}

\begin{exemplebox}[Échelle de 6 m — Calcul de la hauteur atteinte]
Un artisan pose une échelle de longueur $6\,\text{m}$ contre un mur vertical. Le pied de l'échelle est à $1,5\,\text{m}$ du mur. On suppose que le sol est horizontal et que l'échelle, le mur et le sol forment un triangle rectangle.
\begin{enumerate}
    \item Calculer l'angle $\alpha$ que l'échelle fait avec le sol.
    \item Calculer la hauteur $h$ atteinte sur le mur.
    \item Vérifier $h$ par le théorème de Pythagore.
\end{enumerate}
\end{exemplebox}

\begin{popfigure}
\centering
\begin{tikzpicture}[scale=0.7, >=Stealth]
% Sol
\draw[popline, thick] (-1,0) -- (8,0) node[right] {Sol};
% Mur
\draw[popline, thick] (1.5,0) -- (1.5,6.5) node[above] {Mur};
% Échelle (hypoténuse)
\draw[poporange, very thick] (0,0) -- (1.5,5.81) node[midway, above left=2pt, poporange] {\textbf{6 m}};
% Côté adjacent (au sol)
\draw[popblue, thick, <->] (0,-0.4) -- (1.5,-0.4) node[midway, below] {\textbf{1,5 m}};
% Côté opposé (hauteur)
\draw[popgreen, thick, <->] (1.7,0) -- (1.7,5.81) node[midway, right] {\textbf{$h$}};
% Angle au sol
\draw[poppurple, thick] (0.3,0) arc (0:75.5:0.3) node[right=2pt, poppurple] {$\alpha$};
% Angle droit
\draw[popdark, thick] (1.5,0) rectangle (1.7,0.2);
% Points
\fill[popdark] (0,0) circle (3pt) node[below left] {$B$ (pied)};
\fill[popdark] (1.5,0) circle (3pt) node[below right] {$C$};
\fill[popdark] (1.5,5.81) circle (3pt) node[above right] {$A$ (sommet)};
\end{tikzpicture}
\legende{Schéma annoté : triangle rectangle $ABC$ rectangle en $C$, $\widehat{B}=\alpha$, $AB=6\,\text{m}$ (hyp), $BC=1,5\,\text{m}$ (adj), $AC=h$ (opp).}
\end{popfigure}

\textbf{Résolution :}

\begin{enumerate}
    \item \textbf{Angle $\alpha$ :} Par rapport à $\alpha$, le côté connu $BC=1,5\,\text{m}$ est \emph{adjacent}, l'hypoténuse $AB=6\,\text{m}$ est connue. Le rapport qui lie Adj et Hyp est le \textbf{cosinus}.
    \[
    \cos\alpha = \frac{\text{Adj}}{\text{Hyp}} = \frac{BC}{AB} = \frac{1,5}{6} = 0,25
    \]
    Sur calculatrice (mode DEG) : $\alpha = \arccos(0,25) \approx 75,522...°$. \textbf{Arrondi au dixième : $\alpha \approx 75,5°$.}

    \item \textbf{Hauteur $h$ (côté opposé) :} On connaît maintenant $\alpha \approx 75,5°$ et l'hypoténuse $6\,\text{m}$. Le rapport qui lie Opp et Hyp est le \textbf{sinus}.
    \[
    \sin\alpha = \frac{\text{Opp}}{\text{Hyp}} = \frac{h}{6} \quad\Rightarrow\quad h = 6 \times \sin(75,5°)
    \]
    Calculatrice : $6 \times \sin(75,5°) \approx 5,809...$ \textbf{Arrondi au cm : $h \approx 5,81\,\text{m}$.}

    \item \textbf{Vérification par Pythagore :} Dans le triangle rectangle $ABC$,
    \[
    AB^2 = AC^2 + BC^2 \quad\Rightarrow\quad AC = \sqrt{AB^2 - BC^2} = \sqrt{6^2 - 1,5^2} = \sqrt{36 - 2,25} = \sqrt{33,75} \approx 5,809...
    \]
    On retrouve $h \approx 5,81\,\text{m}$. \checkmark \quad La cohérence est validée.
\end{enumerate}

\courssub{Cas particulier : l'angle connu ne « touche » pas le côté connu}

Il arrive que l'angle aigu connu ne soit pas adjacent au côté connu. Dans un triangle rectangle $ABC$ rectangle en $A$, si l'on connaît l'angle $\widehat{B}$ et le côté $AC$ (opposé à $\widehat{B}$), tout va bien : $AC$ est l'opposé par rapport à $\widehat{B}$. Mais si l'on connaît $\widehat{B}$ et le côté $AB$ (l'hypoténuse), c'est encore direct.

Le piège survient quand on connaît \textbf{l'angle $C$} (l'autre angle aigu) et un côté adjacent à $B$. \textbf{Règle d'or :} \emph{Les noms « adjacent » et « opposé » dépendent de l'angle de référence.} Si vous changez d'angle, les noms changent.

\begin{attention}[Piège : changer d'angle sans renommer les côtés]
Dans un triangle rectangle $ABC$ rectangle en $A$, supposons qu'on connaisse $\widehat{C}=30°$ et le côté $AB=10\,\text{cm}$.
\begin{itemize}
    \item \textbf{Erreur fréquente :} Écrire $\cos 30° = \frac{AB}{AC}$ en croyant que $AB$ est adjacent à $\widehat{C}$. \textbf{Faux !} Par rapport à $\widehat{C}$, le côté $AB$ est \emph{opposé}, et $AC$ est l'\emph{hypoténuse}.
    \item \textbf{Bonne méthode :} Soit on utilise l'angle $\widehat{C}$ et on nomme correctement : $\sin 30° = \frac{AB}{AC}$ (Opp/Hyp). Soit on utilise l'angle complémentaire $\widehat{B}=60°$ : par rapport à $\widehat{B}$, $AB$ devient \emph{adjacent} et $AC$ reste l'hypoténuse, donc $\cos 60° = \frac{AB}{AC}$. Les deux donnent $AC = \frac{10}{\sin 30°} = \frac{10}{\cos 60°} = 20\,\text{cm}$.
\end{itemize}
\textbf{Astuce :} Les deux angles aigus sont \textbf{complémentaires} ($\widehat{B}+\widehat{C}=90°$). On a toujours $\sin\widehat{B}=\cos\widehat{C}$ et $\cos\widehat{B}=\sin\widehat{C}$. Choisissez l'angle qui rend le repérage le plus simple.
\end{attention}

\courssub{Exemple avec angle complémentaire}

\begin{exoresolu}[Toit d'une case — Angle au faîtage connu]
Le toit d'une case traditionnelle forme un triangle isocèle. La moitié du toit est un triangle rectangle $ABC$ rectangle en $A$ (sommet du faîtage), avec $\widehat{B}=30°$ (angle au faîtage) et la hauteur du faîtage au plafond $AC = 2,5\,\text{m}$. Calculer la longueur de la panne (hypoténuse $BC$) et la demi-largeur de la case $AB$.
\end{exoresolu}

\textbf{Solution :}

\begin{popfigure}
\centering
\begin{tikzpicture}[scale=0.8, >=Stealth]
% Triangle rectangle
\coordinate (A) at (0,0);
\coordinate (B) at (4.33,0);
\coordinate (C) at (0,2.5);
\draw[popdark, very thick] (A) -- (B) node[midway, below] {$AB$ (demi-largeur)};
\draw[popdark, very thick] (A) -- (C) node[midway, left] {$AC = 2,5\,\text{m}$};
\draw[poporange, very thick] (B) -- (C) node[midway, above right=2pt, poporange] {$BC$ (panne)};
% Angle droit
\draw[popdark, thick] (0,0) rectangle (0.2,0.2);
% Angle B = 30°
\draw[poppurple, thick] (4.03,0) arc (180:210:0.3) node[below left=2pt, poppurple] {$30°$};
% Points
\fill[popdark] (A) circle (3pt) node[below left] {$A$};
\fill[popdark] (B) circle (3pt) node[below right] {$B$};
\fill[popdark] (C) circle (3pt) node[above left] {$C$};
\end{tikzpicture}
\legende{Triangle rectangle $ABC$ rectangle en $A$, $\widehat{B}=30°$, $AC=2,5\,\text{m}$ (opposé à $\widehat{B}$).}
\end{popfigure}

Par rapport à $\widehat{B}=30°$ :
\begin{itemize}
    \item $AC = 2,5\,\text{m}$ est le côté \textbf{opposé} (il ne touche pas $\widehat{B}$).
    \item $BC$ est l'\textbf{hypoténuse}.
    \item $AB$ est le côté \textbf{adjacent}.
\end{itemize}

\textbf{1. Calcul de $BC$ (hypoténuse) :} On connaît l'opposé, on cherche l'hypoténuse $\rightarrow$ \textbf{sinus}.
\[
\sin 30° = \frac{AC}{BC} = \frac{2,5}{BC}
\]
Valeur exacte : $\sin 30° = \frac{1}{2}$. Donc $\frac{1}{2} = \frac{2,5}{BC} \;\Rightarrow\; BC = 2 \times 2,5 = \textbf{5 m}$.
\emph{(Note : dans un triangle rectangle, le côté opposé à 30° vaut la moitié de l'hypoténuse. Vérification immédiate !)}

\textbf{2. Calcul de $AB$ (adjacent) :} On connaît l'opposé, on cherche l'adjacent $\rightarrow$ \textbf{tangente}.
\[
\tan 30° = \frac{AC}{AB} = \frac{2,5}{AB}
\]
Valeur exacte : $\tan 30° = \frac{\sqrt{3}}{3}$. Donc $\frac{\sqrt{3}}{3} = \frac{2,5}{AB} \;\Rightarrow\; AB = \frac{2,5 \times 3}{\sqrt{3}} = \frac{7,5}{\sqrt{3}} = 2,5\sqrt{3} \approx \textbf{4,33 m}$.
\emph{Autre méthode :} Pythagore $AB = \sqrt{BC^2 - AC^2} = \sqrt{25 - 6,25} = \sqrt{18,75} = 2,5\sqrt{3}$. \checkmark

\courssub{Vérifications de cohérence (réflexes à acquérir)}

\begin{aretenir}[Contrôles rapides avant de rendre sa copie]
\begin{enumerate}
    \item \textbf{L'hypoténuse est le plus grand côté.} Si vous trouvez un côté plus grand que l'hypoténuse, il y a une erreur (mauvaise formule, mode RAD au lieu de DEG, inversion produit en croix).
    \item \textbf{Côté opposé à 30° = $\frac{1}{2}$ hypoténuse.} Si $\alpha=30°$, vérifiez que $\text{Opp} = \frac{\text{Hyp}}{2}$.
    \item \textbf{Côté opposé à 45° = Côté adjacent.} Si $\alpha=45°$, les deux côtés de l'angle droit sont égaux.
    \item \textbf{Ordre de grandeur :} $\sin\alpha$ et $\cos\alpha$ sont entre $0$ et $1$ ; $\tan\alpha$ peut être $>1$ (si $\alpha>45°$) ou $<1$ (si $\alpha<45°$).
    \item \textbf{Unités :} La réponse doit porter la même unité que la donnée (cm, m, km...).
\end{enumerate}
\end{aretenir}

\courssub{Le savais-tu ? : La règle du 3-4-5 sur les chantiers}

\begin{savaistu}[Des pyramides aux toits de Yaoundé]
Bien avant les calculatrices, les bâtisseurs égyptiens utilisaient une corde à 13 nœuds équidistants pour tracer des angles droits parfaits : en formant un triangle de côtés 3, 4 et 5 (unités quelconques), l'angle opposé au côté 5 est \textbf{exactement un angle droit} car $3^2+4^2=5^2$ (réciproque de Pythagore).

Aujourd'hui encore, au Cameroun comme ailleurs, les maçons, menuisiers et couvreurs utilisent la \textbf{règle du 3-4-5} (ou 30-40-50 cm, 3-4-5 m) pour vérifier l'équerrage d'un mur, d'une dalle ou d'une charpente. La trigonométrie que vous apprenez généralise cette astuce : elle permet de tracer \emph{n'importe quel angle}, pas seulement 90°, et de calculer des longueurs impossibles à mesurer directement (hauteur d'un pylône MTN, largeur d'un fleuve, pente d'une route d'électrification rurale).
\end{savaistu}

\courssub{Exercices d'application}

\begin{exercice}[Entraînement — Calcul de longueurs]
\begin{enumerate}
    \item Dans un triangle rectangle $ABC$ rectangle en $A$, on donne $\widehat{B}=40°$ et $AB=8\,\text{cm}$. Calculer $AC$ et $BC$ (arrondi au mm).
    \item Une antenne-relais $OT$ verticale de hauteur $30\,\text{m}$ est haubanée par un câble $TA$ attaché au sol en $A$. Le câble fait un angle de $60°$ avec le sol. Calculer la distance $OA$ (pied de l'antenne au point d'ancrage) et la longueur du câble $TA$ (valeurs exactes puis arrondies au cm).
    \item Un bateau quitte le port $P$ et navigue cap au Nord sur $12\,\text{km}$ jusqu'au point $B$. Il vire alors à l'Est sur $5\,\text{km}$ jusqu'au point $C$. Calculer l'angle que fait la route directe $PC$ avec le Nord (arrondi au degré) et la distance $PC$.
    \item \textbf{Problème concret :} Pour installer un panneau solaire sur un toit incliné à $25°$, un technicien doit fixer une barre de renfort le long de la pente. Si la hauteur verticale à couvrir est de $1,8\,\text{m}$, quelle doit être la longueur de la barre (arrondi au cm) ?
\end{enumerate}
\end{exercice}

\begin{corrige}[Exercice 1]
Triangle $ABC$ rectangle en $A$, $\widehat{B}=40°$, $AB=8\,\text{cm}$ (adjacent à $\widehat{B}$).
\begin{itemize}
    \item $AC$ (opposé) : $\tan 40° = \frac{AC}{8} \;\Rightarrow\; AC = 8 \times \tan 40° \approx 8 \times 0,8391 \approx \textbf{6,71 cm}$.
    \item $BC$ (hypoténuse) : $\cos 40° = \frac{8}{BC} \;\Rightarrow\; BC = \frac{8}{\cos 40°} \approx \frac{8}{0,7660} \approx \textbf{10,44 cm}$.
\end{itemize}
Vérification : $BC > AB$ et $BC > AC$ \checkmark. Pythagore : $\sqrt{8^2+6,71^2} \approx \sqrt{64+45,0} \approx \sqrt{109,0} \approx 10,44$ \checkmark.
\end{corrige}

\begin{corrige}[Exercice 2]
Triangle $OAT$ rectangle en $O$, $\widehat{A}=60°$, $OT=30\,\text{m}$ (opposé à $\widehat{A}$).
\begin{itemize}
    \item $OA$ (adjacent) : $\tan 60° = \frac{30}{OA} \;\Rightarrow\; OA = \frac{30}{\tan 60°} = \frac{30}{\sqrt{3}} = 10\sqrt{3} \approx \textbf{17,32 m}$.
    \item $TA$ (hypoténuse) : $\sin 60° = \frac{30}{TA} \;\Rightarrow\; TA = \frac{30}{\sin 60°} = \frac{30}{\frac{\sqrt{3}}{2}} = 20\sqrt{3} \approx \textbf{34,64 m}$.
\end{itemize}
Note : $TA = 2 \times OA$ car dans un triangle rectangle à $30°$-$60°$, l'hypoténuse est double du côté opposé à $30°$ (ici $\widehat{O}=30°$).
\end{corrige}

\begin{corrige}[Exercice 3]
Triangle $PBC$ rectangle en $B$ ($PB=12\,\text{km}$ Nord, $BC=5\,\text{km}$ Est). $PC$ est l'hypoténuse.
\begin{itemize}
    \item $\tan(\widehat{P}) = \frac{BC}{PB} = \frac{5}{12} \approx 0,4167 \;\Rightarrow\; \widehat{P} \approx \arctan(0,4167) \approx \textbf{23°}$.
    \item $PC = \sqrt{12^2+5^2} = \sqrt{144+25} = \sqrt{169} = \textbf{13 km}$ (triplet pythagoricien 5-12-13).
\end{itemize}
\end{corrige}

\begin{corrige}[Exercice 4]
Triangle rectangle : angle au sol $\alpha=25°$, hauteur $h=1,8\,\text{m}$ (opposé à $\alpha$), barre $L$ = hypoténuse.
\[
\sin 25° = \frac{1,8}{L} \;\Rightarrow\; L = \frac{1,8}{\sin 25°} \approx \frac{1,8}{0,4226} \approx 4,259\,\text{m} \approx \textbf{4,26 m} \text{ (soit 426 cm)}.
\]
\end{corrige}

\courssub{Synthèse de la section}

\begin{aretenir}[À retenir pour le BEPC]
\begin{itemize}
    \item \textbf{Schéma annoté} obligatoire : angle de référence, côtés nommés Hyp/Adj/Opp.
    \item \textbf{Choix de la formule} : celui qui contient le côté connu et le côté cherché.
    \item \textbf{Produit en croix} pour isoler l'inconnue.
    \item \textbf{Calculatrice en mode DEG} (vérifiez l'affichage « D » ou « DEG »).
    \item \textbf{Vérifications} : hypoténuse $>$ autres côtés ; côté opposé à 30° = $\frac{1}{2}$ hypoténuse ; cohérence par Pythagore.
    \item \textbf{Conclusion} avec unité et arrondi demandé.
\end{itemize}
\end{aretenir}

La section suivante (« 5. Calculer un angle aigu dans un triangle rectangle ») utilisera les mêmes rapports, mais cette fois l'inconnue sera l'angle : nous apprendrons à manipuler les fonctions \textbf{arccos}, \textbf{arcsin} et \textbf{arctan} sur la calculatrice.

\coursec{Calculer un angle aigu dans un triangle rectangle}

Dans la section précédente, nous avons appris à \emph{calculer une longueur} dans un triangle rectangle lorsqu'on connaît un angle aigu et un côté : on choisissait le bon rapport (sinus, cosinus ou tangente), puis on isolait la longueur inconnue. Mais la vie courante pose souvent le problème \emph{en sens inverse} : un maçon à Yaoundé connaît la base et la hauteur du toit d'une case et veut connaître \textbf{l'angle de pente} ; un technicien de CAMTEL connaît la longueur d'un câble et la distance au poteau et cherche \textbf{l'angle d'inclinaison} du câble. Bref : on connaît \emph{des longueurs}, et on cherche \emph{un angle}. C'est exactement l'objet de cette section : \cle{remonter d'un rapport trigonométrique à l'angle} grâce aux \cle{fonctions inverses} de la calculatrice.

\courssub{De la valeur du rapport à l'angle : l'idée de fonction inverse}

\textbf{Une question de sens : « dans l'autre sens »}\par

Reprenons la situation de la section précédente. Si un triangle rectangle a un angle aigu de $30°$, alors le rapport $\dfrac{\text{côté opposé}}{\text{hypoténuse}}$ vaut $\sin(30°) = 0,5$. La calculatrice répond à la question : \emph{« connaissant l'angle, quel est le rapport ? »}

Posons maintenant la question inverse : \emph{« je sais que le rapport côté opposé / hypoténuse vaut $0,5$ ; quel est l'angle ? »} La réponse est $30°$. On dit qu'on applique la \textbf{fonction inverse} du sinus, notée $\sin^{-1}$ (ou $\arcsin$).

\begin{popfigure}
\begin{center}
\begin{tikzpicture}[scale=1, font=\small]
  % Boîte aller
  \draw[rounded corners, fill=popblueL, draw=popblue] (-4.6,-0.6) rectangle (-1.4,0.6);
  \node at (-3,0.18) {\textbf{Sens direct}};
  \node at (-3,-0.25) {angle $\longrightarrow$ rapport};
  \node[popblue] at (-3,-0.85) {$\cos(60°) = 0,5$};
  % Boîte retour
  \draw[rounded corners, fill=poporangeL, draw=poporange] (1.4,-0.6) rectangle (4.6,0.6);
  \node at (3,0.18) {\textbf{Sens inverse}};
  \node at (3,-0.25) {rapport $\longrightarrow$ angle};
  \node[poporange] at (3,-0.85) {$\cos^{-1}(0,5) = 60°$};
  % Flèches
  \draw[-{Stealth}, very thick, popblue] (-1.3,0.25) -- (1.3,0.25) node[midway, above, black]{$\cos$};
  \draw[-{Stealth}, very thick, poporange] (1.3,-0.25) -- (-1.3,-0.25) node[midway, below, black]{$\cos^{-1}$};
\end{tikzpicture}
\legende{La fonction cosinus et sa fonction inverse $\cos^{-1}$ : deux questions, deux sens de lecture.}
\end{center}
\end{popfigure}

\textbf{Pourquoi l'angle est-il unique ?}\par

On pourrait craindre qu'une même valeur de rapport corresponde à plusieurs angles aigus. Il n'en est rien, et c'est une propriété essentielle (admise au programme de 3ème) :

\begin{propriete}[Unicité de l'angle aigu (admise)]
Sur l'intervalle $[0°\,;\,90°]$, les fonctions cosinus, sinus et tangente sont \cle{strictement monotones} :
\begin{itemize}
  \item le \textbf{cosinus} est strictement \emph{décroissant} de $1$ à $0$ ;
  \item le \textbf{sinus} est strictement \emph{croissant} de $0$ à $1$ ;
  \item la \textbf{tangente} est strictement \emph{croissante} à partir de $0$ et prend de grandes valeurs lorsque l'angle se rapproche de $90°$.
\end{itemize}
Par conséquent, à une valeur donnée du rapport correspond \textbf{un seul} angle aigu : la correspondance entre l'angle aigu et la valeur de chaque rapport est \cle{biunivoque} (à chaque angle une seule valeur, et à chaque valeur un seul angle).
\end{propriete}

Concrètement : si $\cos(\alpha) = 0,5$ avec $\alpha$ aigu, alors $\alpha$ ne peut valoir \emph{que} $60°$. C'est cette unicité qui autorise la calculatrice à afficher une réponse unique quand on tape $\cos^{-1}(0,5)$.

\begin{popfigure}
\begin{center}
\begin{tikzpicture}[scale=1]
\begin{axis}[
    width=0.85\linewidth, height=5.2cm,
    axis lines=middle,
    xlabel={angle (degrés)}, ylabel={valeur},
    xmin=0, xmax=95, ymin=0, ymax=1.15,
    xtick={0,30,45,60,90},
    ytick={0,0.5,0.7071,0.866,1},
    yticklabels={$0$,$0{,}5$,$\frac{\sqrt2}{2}$,$\frac{\sqrt3}{2}$,$1$},
    grid=both, grid style={popline!40},
    legend style={at={(0.03,0.97)}, anchor=north west, font=\small}
]
  \addplot[popcurve, domain=0:90, samples=120, popblue, very thick] {cos(x)};
  \addlegendentry{$\cos(\alpha)$ (décroît)}
  \addplot[popcurve, domain=0:90, samples=120, poporange, very thick] {sin(x)};
  \addlegendentry{$\sin(\alpha)$ (croît)}
  % Lecture inverse : cos = 0.5 -> 60°
  \addplot[dashed, popdark] coordinates {(0,0.5) (60,0.5)};
  \addplot[dashed, popdark] coordinates {(60,0.5) (60,0)};
  \addplot[only marks, mark=*, popdark] coordinates {(60,0.5)};
  \node[popdark, anchor=south west, font=\small] at (axis cs:2,0.52) {$\cos(\alpha)=0{,}5 \Rightarrow \alpha=60°$};
\end{axis}
\end{tikzpicture}
\legende{Lecture graphique : chaque valeur entre $0$ et $1$ est atteinte une seule fois par le cosinus (et une seule fois par le sinus) sur $[0°;90°]$.}
\end{center}
\end{popfigure}

\courssub{Les fonctions inverses : définition et notation}

\begin{definition}[Fonctions trigonométriques inverses]
\begin{itemize}
  \item Pour tout nombre $x \in [0\,;\,1]$, on note $\arccos(x)$ ou $\cos^{-1}(x)$ l'\textbf{unique} angle $\alpha \in [0°\,;\,90°]$ tel que $\cos(\alpha) = x$.
  \item Pour tout nombre $x \in [0\,;\,1]$, on note $\arcsin(x)$ ou $\sin^{-1}(x)$ l'\textbf{unique} angle $\alpha \in [0°\,;\,90°]$ tel que $\sin(\alpha) = x$.
  \item Pour tout nombre $x \geq 0$, on note $\arctan(x)$ ou $\tan^{-1}(x)$ l'\textbf{unique} angle $\alpha \in [0°\,;\,90°[$ tel que $\tan(\alpha) = x$.
\end{itemize}
\end{definition}

\begin{aretenir}[Les équivalences à connaître par cœur]
Pour un angle aigu $\alpha$ :
\[ \cos(\alpha) = x \iff \alpha = \cos^{-1}(x) \qquad\qquad \sin(\alpha) = x \iff \alpha = \sin^{-1}(x) \]
\[ \tan(\alpha) = x \iff \alpha = \tan^{-1}(x) \]
On « fait passer » la fonction de l'autre côté en la remplaçant par sa fonction inverse.
\end{aretenir}

\begin{attention}[Piège majeur : $\cos^{-1}$ ne veut PAS dire $\dfrac{1}{\cos}$]
Sur la calculatrice, la touche $\cos^{-1}$ (obtenue le plus souvent par \textbf{SHIFT} ou \textbf{2nd} puis \textbf{COS}) désigne la \textbf{fonction réciproque} du cosinus, et \emph{surtout pas} l'inverse $\dfrac{1}{\cos(\alpha)}$ (qui s'appelle la sécante et qui est hors programme). Avant de taper ton quotient, vérifie que l'écran affiche bien \texttt{cos\textsuperscript{-1}(} et non \texttt{cos(}. Autre piège : ta calculatrice doit être en \textbf{mode degrés} (DEG affiché à l'écran), sinon le résultat sera faux !
\end{attention}

\begin{popfigure}
\begin{center}
\begin{tikzpicture}[font=\small]
  % Corps calculatrice
  \draw[rounded corners=6pt, fill=popdark!85, draw=popdark] (-2.3,-3.4) rectangle (2.3,2.6);
  % Écran
  \draw[rounded corners=2pt, fill=popgreenL, draw=popdark] (-1.9,1.2) rectangle (1.9,2.3);
  \node[anchor=west, font=\ttfamily\small] at (-1.75,1.95) {cos\textasciicircum -1(0.75)};
  \node[anchor=west, font=\ttfamily\small] at (-1.75,1.5) {41.4096...};
  \node[anchor=east, font=\tiny\ttfamily, popdark] at (1.85,2.18) {DEG};
  % Touches
  \foreach \x/\y in {-1.5/0.4, -0.5/0.4, 0.5/0.4, 1.5/0.4}{
    \draw[rounded corners=2pt, fill=popmuted!25, draw=popmuted] (\x-0.38,\y-0.22) rectangle (\x+0.38,\y+0.22);
  }
  \node[font=\tiny\bfseries, poporange] at (-1.5,0.4) {SHIFT};
  % Rangée sin cos tan
  \foreach \x/\t in {-1.5/{sin}, -0.5/{cos}, 0.5/{tan}}{
    \draw[rounded corners=2pt, fill=popmuted!25, draw=popmuted] (\x-0.38,-0.5) rectangle (\x+0.38,-0.1);
    \node[font=\tiny] at (\x,-0.3) {\t};
    \node[font=\tiny, poporange] at (\x,-0.62) {\t\textasciicircum -1};
  }
  \foreach \x/\t in {-1.5/{7}, -0.5/{8}, 0.5/{9}, 1.5/{(}}{
    \draw[rounded corners=2pt, fill=popmuted!15, draw=popmuted] (\x-0.38,-1.3) rectangle (\x+0.38,-0.9);
    \node[font=\tiny] at (\x,-1.1) {\t};
  }
  \foreach \x/\t in {-1.5/{4}, -0.5/{5}, 0.5/{6}, 1.5/{)}}{
    \draw[rounded corners=2pt, fill=popmuted!15, draw=popmuted] (\x-0.38,-2.1) rectangle (\x+0.38,-1.7);
    \node[font=\tiny] at (\x,-1.9) {\t};
  }
  \foreach \x/\t in {-1.5/{1}, -0.5/{2}, 0.5/{3}, 1.5/{=}}{
    \draw[rounded corners=2pt, fill=popmuted!15, draw=popmuted] (\x-0.38,-2.9) rectangle (\x+0.38,-2.5);
    \node[font=\tiny] at (\x,-2.7) {\t};
  }
  % Flèche SHIFT -> cos
  \draw[-{Stealth}, thick, poporange] (-1.5,0.62) .. controls (-1.2,1.0) and (-0.8,0.9) .. (-0.55,-0.05);
  \node[poporange, font=\tiny, align=center] at (-2.9,0.9) {1. SHIFT\\puis 2. cos};
\end{tikzpicture}
\legende{Calculatrice (schéma) : pour obtenir $\cos^{-1}$, on appuie d'abord sur SHIFT (ou 2nd), puis sur COS. L'écran affiche \texttt{cos\textasciicircum -1(} et le mode DEG doit être actif.}
\end{center}
\end{popfigure}

\courssub{Protocole de calcul d'un angle aigu}

\begin{methode}[Calculer un angle aigu connaissant deux côtés]
\begin{enumerate}
  \item \textbf{Faire un schéma} du triangle rectangle, en nommant les sommets et en marquant l'angle droit.
  \item \textbf{Repérer les deux côtés connus} par rapport à l'angle $\alpha$ cherché : hypoténuse, côté adjacent, côté opposé.
  \item \textbf{Choisir le rapport} qui fait intervenir ces deux côtés (SOH -- CAH -- TOA) : par exemple adjacent et hypoténuse $\Rightarrow$ cosinus.
  \item \textbf{Écrire l'égalité} et calculer le quotient : $\cos(\alpha) = \dfrac{\text{adjacent}}{\text{hypoténuse}}$.
  \item \textbf{Appliquer la fonction inverse} : $\alpha = \cos^{-1}(\text{quotient})$, à la calculatrice en mode degrés.
  \item \textbf{Vérifier} que le résultat est cohérent : $0° < \alpha < 90°$.
  \item \textbf{Conclure} par une phrase, avec l'arrondi demandé (au degré près ou au dixième de degré).
\end{enumerate}
\end{methode}

\begin{attention}[Conditions d'application à ne jamais oublier]
\begin{itemize}
  \item Le quotient d'un cosinus ou d'un sinus doit être compris \textbf{entre 0 et 1} : si tu obtiens $\cos(\alpha) = 1,4$, c'est qu'il y a une erreur (côtés mal identifiés, ou hypoténuse au numérateur !). Rappelle-toi : l'hypoténuse est le plus grand côté.
  \item La tangente, elle, peut dépasser $1$ (angle supérieur à $45°$) : ce n'est pas une erreur.
  \item Précise toujours le triangle dans lequel tu travailles : « Dans le triangle $ABC$ rectangle en $A$... ».
\end{itemize}
\end{attention}

\begin{exoresolu}[Angle d'inclinaison d'une échelle]
Une échelle de $5$ m est posée contre un mur ; son pied est à $1,5$ m du mur. On note $\alpha$ l'angle entre l'échelle et le sol. Calculer $\alpha$ au dixième de degré près.
\tcblower
\textbf{Solution.}
\begin{enumerate}
  \item Schéma : le mur, le sol et l'échelle forment un triangle rectangle ; l'angle droit est entre le mur et le sol.
  \item Par rapport à $\alpha$ : le côté adjacent est la distance pied--mur ($1,5$ m) ; l'hypoténuse est l'échelle ($5$ m).
  \item Adjacent + hypoténuse $\Rightarrow$ \textbf{cosinus}.
  \item $\cos(\alpha) = \dfrac{1,5}{5} = 0,3$.
  \item $\alpha = \cos^{-1}(0,3) \approx 72,5°$ (calculatrice en mode DEG).
  \item Vérification : $0° < 72,5° < 90°$ : cohérent (échelle très « dressée », angle proche de $90°$).
  \item \textbf{Conclusion :} l'angle entre l'échelle et le sol vaut environ $\alpha \approx 72,5°$.
\end{enumerate}
\end{exoresolu}

\begin{exemplebox}[La pente du toit d'une case]
Un charpentier construit le toit d'une case : la base (largeur) mesure $4$ m et la hauteur du faîtage au-dessus de la base est $1,5$ m. Le faîtage est au milieu, donc la \textbf{demi-base} vaut $2$ m. Quel est l'angle de pente $\alpha$ du toit ?

Par rapport à $\alpha$ : le côté adjacent est la demi-base ($2$ m), le côté opposé est la hauteur ($1,5$ m). Opposé + adjacent $\Rightarrow$ \textbf{tangente} :
\[ \tan(\alpha) = \frac{1,5}{2} = 0,75 \qquad\text{donc}\qquad \alpha = \tan^{-1}(0,75) \approx 36,9°. \]
\textbf{Vérification :} l'autre angle aigu du triangle vaut $90° - 36,9° = 53,1°$, et on a bien $90° + 36,9° + 53,1° = 180°$. Le toit a une pente d'environ $36,9°$ : raisonnable pour évacuer les pluies de la saison humide !
\end{exemplebox}

\begin{popfigure}
\begin{center}
\begin{tikzpicture}[scale=1.15, font=\small]
  \coordinate (A) at (0,0);
  \coordinate (B) at (4,0);
  \coordinate (C) at (4,1.5);
  \draw[thick, popdark] (A) -- (B) -- (C) -- cycle;
  \draw[popdark] (3.7,0) -- (3.7,0.3) -- (4,0.3);
  \node[below] at (2,0) {demi-base $= 2$ m (adjacent)};
  \node[right, rotate=90] at (4.15,0.75) {\;hauteur $= 1{,}5$ m (opposé)};
  \draw[poporange, thick] (1.1,0) arc (0:36.87:1.1);
  \node[poporange] at (1.55,0.22) {$\alpha = ?$};
  % Flèche vers formule
  \draw[-{Stealth}, very thick, popblue] (5.0,0.75) -- (6.1,0.75);
  \node[anchor=west, popblue, align=left] at (6.2,0.75) {$\tan(\alpha)=\dfrac{1{,}5}{2}=0{,}75$\\[2pt] $\alpha=\tan^{-1}(0{,}75)\approx 36{,}9°$};
\end{tikzpicture}
\legende{Deux côtés connus, un angle inconnu : on choisit la tangente, puis on applique $\tan^{-1}$.}
\end{center}
\end{popfigure}

\courssub{Sans calculatrice : reconnaître les valeurs exactes}

Pour les angles remarquables $30°$, $45°$ et $60°$ étudiés dans la section sur les angles remarquables, on peut retrouver l'angle \emph{sans calculatrice} dès qu'on reconnaît une valeur exacte du tableau :

\begin{center}
\begin{tabular}{|c|c|c|c|}
\hline
\rowcolor{popblueL} $\alpha$ & $30°$ & $45°$ & $60°$ \\
\hline
$\cos(\alpha)$ & $\dfrac{\sqrt{3}}{2}$ & $\dfrac{\sqrt{2}}{2}$ & $\dfrac{1}{2}$ \\
\hline
$\sin(\alpha)$ & $\dfrac{1}{2}$ & $\dfrac{\sqrt{2}}{2}$ & $\dfrac{\sqrt{3}}{2}$ \\
\hline
$\tan(\alpha)$ & $\dfrac{\sqrt{3}}{3}$ & $1$ & $\sqrt{3}$ \\
\hline
\end{tabular}
\end{center}

\begin{exemplebox}[Lecture à l'envers du tableau]
\begin{itemize}
  \item Si $\cos(\alpha) = \dfrac{1}{2}$, alors $\alpha = 60°$ (le cosinus \emph{décroît} : petite valeur $\Rightarrow$ grand angle).
  \item Si $\sin(\alpha) = \dfrac{1}{2}$, alors $\alpha = 30°$.
  \item Si $\tan(\alpha) = 1$, alors $\alpha = 45°$ (triangle rectangle isocèle).
  \item Si $\tan(\alpha) = \sqrt{3}$, alors $\alpha = 60°$.
\end{itemize}
\textbf{Astuce de sécurité :} $\cos$ et $\sin$ d'un même angle ne donnent jamais la même ligne du tableau : $\dfrac{1}{2}$ correspond à $60°$ pour le cosinus mais à $30°$ pour le sinus. Vérifie toujours \emph{quel} rapport tu as calculé !
\end{exemplebox}

\courssub{Vérification et cohérence des résultats}

Une bonne copie de BEPC se termine toujours par une vérification. Deux réflexes :

\begin{aretenir}[Deux vérifications rapides]
\begin{enumerate}
  \item \textbf{Somme des angles :} dans un triangle rectangle, les deux angles aigus sont \cle{complémentaires} : $\alpha + \beta = 90°$, car $90° + \alpha + \beta = 180°$. Si tu as trouvé $\alpha \approx 36,9°$, l'autre angle aigu vaut $90° - 36,9° = 53,1°$.
  \item \textbf{Ordre de grandeur :} dans un triangle rectangle, le plus grand angle aigu est \emph{opposé au plus long des deux côtés de l'angle droit}. Si le côté opposé à $\alpha$ est plus petit que le côté adjacent, alors $\alpha < 45°$.
\end{enumerate}
\end{aretenir}

\begin{exoresolu}[Un câble d'ancrage]
Pour stabiliser un poteau électrique, un technicien tend un câble entre le sommet du poteau et le sol. Le câble mesure $8$ m et le point d'ancrage au sol est à $4$ m du pied du poteau. Calculer l'angle $\beta$ que fait le câble avec le poteau (arrondir au degré près), puis vérifier.
\tcblower
\textbf{Solution.} Le poteau est vertical et le sol horizontal : ils forment un angle droit. Par rapport à l'angle $\beta$ (entre le câble et le poteau) :
\begin{itemize}
  \item le côté \textbf{opposé} à $\beta$ est la distance au sol : $4$ m ;
  \item l'\textbf{hypoténuse} est le câble : $8$ m.
\end{itemize}
Opposé + hypoténuse $\Rightarrow$ \textbf{sinus} :
\[ \sin(\beta) = \frac{4}{8} = 0,5 \quad\Longrightarrow\quad \beta = \sin^{-1}(0,5) = 30° \quad\text{(valeur exacte !)} \]
\textbf{Vérification :} l'angle entre le câble et le sol vaut $90° - 30° = 60°$, et $90° + 30° + 60° = 180°$. De plus $0,5 = \sin(30°)$ est une valeur remarquable : cohérent. \textbf{Conclusion :} $\beta = 30°$.
\end{exoresolu}

\begin{exercice}[À toi de jouer]
\begin{enumerate}
  \item Dans un triangle $MNP$ rectangle en $M$, on donne $MN = 6$ cm et $NP = 10$ cm. Calculer l'angle $\widehat{MNP}$ au dixième de degré près.
  \item Une rampe d'accès pour personnes à mobilité réduite s'élève de $0,6$ m sur une longueur horizontale de $7,2$ m. Calculer son angle d'inclinaison au degré près.
  \item Sans calculatrice : dans un triangle rectangle, $\cos(\widehat{A}) = \dfrac{\sqrt{2}}{2}$. Que vaut $\widehat{A}$ ? Que vaut l'autre angle aigu ?
\end{enumerate}
\end{exercice}

\begin{corrige}[Exercice « À toi de jouer »]
\begin{enumerate}
  \item Par rapport à $\widehat{MNP}$ : adjacent $= MN = 6$, hypoténuse $= NP = 10$. $\cos(\widehat{MNP}) = \dfrac{6}{10} = 0,6$, donc $\widehat{MNP} = \cos^{-1}(0,6) \approx 53,1°$.
  \item $\tan(\alpha) = \dfrac{0,6}{7,2} \approx 0,0833$, donc $\alpha = \tan^{-1}(0,0833) \approx 4,8°$, soit environ $5°$ au degré près.
  \item $\widehat{A} = 45°$ (valeur remarquable) ; l'autre angle aigu vaut $90° - 45° = 45°$ : le triangle est rectangle isocèle.
\end{enumerate}
\end{corrige}

\begin{savaistu}[S'orienter grâce aux angles : navigation, topographie et GPS]
En navigation et en topographie, on utilise le \textbf{relèvement} : l'angle mesuré entre la direction du Nord et la direction d'un objectif. Savoir calculer un angle à partir de distances permet de s'orienter en forêt, de tracer une route ou d'implanter une parcelle agricole. Les arpenteurs qui délimitent les terrains à Douala ou à Bafoussam utilisent en permanence la trigonométrie du triangle rectangle ! Quant au GPS de ton téléphone, il repose sur une généralisation de ces idées à trois dimensions : la \emph{trigonométrie sphérique}, qui permet de calculer ta position sur la Terre à partir des signaux de satellites. Les formules que tu apprends aujourd'hui en sont le tout premier pas.
\end{savaistu}

\begin{aretenir}[L'essentiel de la section]
\begin{itemize}
  \item À une valeur de $\cos$, $\sin$ ou $\tan$ correspond \textbf{un unique} angle aigu (monotonie stricte sur $[0°;90°]$).
  \item On remonte à l'angle par la \textbf{fonction inverse} : $\alpha = \cos^{-1}(x)$, $\sin^{-1}(x)$ ou $\tan^{-1}(x)$ (touche SHIFT + cos/sin/tan, calculatrice en mode \textbf{DEG}).
  \item $\cos^{-1}$ n'est \textbf{pas} $\dfrac{1}{\cos}$ : c'est la fonction réciproque.
  \item Le quotient d'un cosinus ou d'un sinus est toujours compris entre $0$ et $1$.
  \item Les valeurs remarquables ($\tfrac12$, $\tfrac{\sqrt2}{2}$, $\tfrac{\sqrt3}{2}$, $1$, $\sqrt 3$) permettent de conclure \textbf{sans calculatrice}.
  \item Toujours vérifier : $0° < \alpha < 90°$ et $\alpha + \beta = 90°$ pour les deux angles aigus.
\end{itemize}
\end{aretenir}

\coursec{Résolution de triangle rectangle et Problèmes concrets (Intégration)}

\courssub{Que signifie « résoudre un triangle rectangle » ?}

Dans les sections précédentes, nous avons appris à calculer une longueur ou un angle isolé. Aujourd'hui, nous allons plus loin : \textbf{résoudre complètement un triangle rectangle}. Cela signifie retrouver \emph{toutes} ses mesures inconnues : les trois longueurs des côtés et les trois mesures des angles (dont l'angle droit, qui vaut toujours $90^\circ$).

\begin{definition}[Résolution d'un triangle rectangle]
Soit un triangle rectangle $ABC$ rectangle en $A$. Le résoudre, c'est déterminer :
\begin{itemize}
    \item les longueurs des trois côtés $AB$, $AC$, $BC$ ;
    \item les mesures des trois angles $\widehat{BAC}=90^\circ$, $\widehat{ABC}$ et $\widehat{ACB}$.
\end{itemize}
\end{definition}

Pour y parvenir, nous n'avons pas besoin de tout connaître au départ. Le programme de 3ème nous garantit que \textbf{deux données bien choisies suffisent} :
\begin{itemize}
    \item \textbf{Cas 1 :} un angle aigu et un côté (n'importe lequel) ;
    \item \textbf{Cas 2 :} deux côtés (n'importe lesquels).
\end{itemize}

\courssub{Stratégie unifiée : la méthode en 6 étapes}

Face à un problème, ne te précipite pas sur la calculatrice. Adopte la démarche du topographe : rigoureuse, ordonnée, sans faille.

\begin{methode}[Résolution complète d'un triangle rectangle]
\textbf{Étape 1 : Schéma réaliste.} Dessine le triangle (même schématique) en respectant l'angle droit. Place les données connues (longueurs, angles) et nomme les inconnues (ex: $x$, $y$, $\alpha$).

\textbf{Étape 2 : Inventaire.} Liste : \textit{Données} (ce que tu sais) / \textit{Cherchés} (ce que tu dois trouver).

\textbf{Étape 3 : Choix de l'angle de travail.} Si un angle aigu est connu, utilise-le. Si tu as deux côtés, choisis l'angle aigu qui te permet d'utiliser directement les côtés connus (adjacent, opposé, hypoténuse).

\textbf{Étape 4 : Chaîne de calculs.}
\begin{itemize}
    \item \textbf{Si 1 angle + 1 côté :} Calcule le 2ème angle aigu ($90^\circ - \alpha$), puis les 2 côtés restants (sinus, cosinus, tangente).
    \item \textbf{Si 2 côtés :} Calcule le 3ème côté (Théorème de Pythagore), puis les 2 angles aigus (fonctions trigonométriques inverses : $\arccos$, $\arcsin$, $\arctan$).
\end{itemize}

\textbf{Étape 5 : Vérification globale.}
\begin{itemize}
    \item Somme des angles : $\alpha + \beta + 90^\circ = 180^\circ$ ?
    \item Pythagore : $c^2 = a^2 + b^2$ (avec les valeurs arrondies, l'égalité est approchée) ?
\end{itemize}

\textbf{Étape 6 : Phrase de réponse.} Conclus par une phrase complète, avec unités et arrondi contextualisé (ex: « La hauteur de l'arbre est d'environ 10,1 m »).
\end{methode}

\begin{attention}[Piège mortel : l'arrondi intermédiaire]
Ne \textbf{jamais} arrondir une valeur intermédiaire pour la réinjecter dans le calcul suivant. Tu perdrais en précision et ton résultat final serait faux.
\emph{Astuce de pro :} Utilise la touche \texttt{ANS} (ou la mémoire) de ta calculatrice. Garde la valeur exacte (ou non arrondie) en mémoire, et n'arrondis \textbf{que le résultat final} demandé par l'énoncé.
\end{attention}

\courssub{Modélisation : des situations concrètes camerounaises}

Les mathématiques ne vivent pas seulement dans les cahiers. Elles sont l'outil des ingénieurs, géomètres, architectes. Voici trois classiques que tu rencontreras au BEPC et dans la vie.

\begin{popfigure}
\begin{tikzpicture}[scale=0.8, every node/.style={font=\small}]
% Rive gauche
\draw[popblue, very thick] (-3,0) -- (5,0) node[right] {Rive gauche};
\draw[popblue!50, dashed] (-3,-0.2) -- (5,-0.2);
% Rive droite
\draw[popblue, very thick] (-3,3) -- (5,3) node[right] {Rive droite};
\draw[popblue!50, dashed] (-3,3.2) -- (5,3.2);
% Points
\coordinate (A) at (0,0);
\coordinate (B) at (4,0);
\coordinate (C) at (0,3);
% Triangle
\draw[popdark, thick] (A) -- node[below] {$d = AB$} (B) -- node[right] {$L = BC$} (C) -- node[left] {Largeur fleuve} (A);
% Angle en B
\draw (B) ++(-0.5,0) arc (180:143.13:0.5) node[midway, left] {$\alpha$};
% Points labels
\node[below left] at (A) {A};
\node[below right] at (B) {B};
\node[above left] at (C) {C};
% Angle droit en A
\draw (A) ++(0.3,0) -- ++(0,0.3) -- ++(-0.3,0);
% Flèches largeur
\draw[<->, poporange] (-0.5,0) -- (-0.5,3) node[midway, left] {Largeur $w$};
\end{tikzpicture}
\legende{Situation « Largeur d'un fleuve » : On mesure une distance $AB$ accessible sur la rive, l'angle $\widehat{ABC}$, on déduit la largeur $AC$ par trigonométrie.}
\end{popfigure}

\begin{popfigure}
\begin{tikzpicture}[scale=0.7, every node/.style={font=\small}]
% Sol
\draw[popmuted, very thick] (-1,0) -- (7,0);
% Arbre (fromager)
\draw[popgreen!70!black, line width=3pt] (2,0) -- (2,4.2) node[midway, right] {$h$};
% Feuillage
\draw[popgreen, fill=popgreen!20] (2,4.2) ellipse (1.2 and 0.8);
% Ombre (longueur ~5 pour angle 40 deg : h=4.2, L=5)
\draw[poporange, line width=2pt] (2,0) -- (7,0) node[midway, below] {$L$};
% Soleil
\draw[popgold, fill=popgold!30] (7.5,4) circle (0.4);
% Rayon soleil -> cime
\draw[popgold, dashed, ->] (7.5,4) -- (2,4.2);
% Angle élévation au bout de l'ombre (7,0)
\draw (7,0) ++(-0.6,0) arc (180:140:0.6) node[midway, above left] {$40^\circ$};
% Légendes
\node[below] at (2,0) {Pied de l'arbre};
\node[above] at (2,4.8) {Cime};
\end{tikzpicture}
\legende{Situation « Hauteur d'un fromager » : On mesure l'ombre $L$ et l'angle d'élévation du soleil $\alpha$, on calcule la hauteur $h = L \times \tan(\alpha)$.}
\end{popfigure}

\courssub{La pente en pourcentage : lien avec la tangente}

Au Cameroun, sur les routes de montagne (vers Dschang, Bafoussam, ou la route du Littoral), tu vois des panneaux « 10\% », « 12\% ». Ce n'est pas un angle, c'est un \textbf{pourcentage de pente}.

\begin{propriete}[Pourcentage de pente et angle d'inclinaison]
Sur un panneau routier, une pente de $p\%$ signifie que pour \textbf{100 m horizontaux}, on monte (ou descend) de \textbf{$p$ mètres}.
\[
\tan(\alpha) = \frac{\text{dénivelée}}{\text{horizontal}} = \frac{p}{100}
\]
L'angle d'inclinaison $\alpha$ (en degrés) s'obtient par :
\[
\alpha = \arctan\left(\frac{p}{100}\right)
\]
\textbf{Exemple :} Pente $10\% \Rightarrow \tan(\alpha)=0,1 \Rightarrow \alpha \approx 5,7^\circ$.
\end{propriete}

\begin{savaistu}[Métiers : Topographe / Géomètre-expert (MINDCAF)]
Les agents du cadastre (Ministère des Domaines, du Cadastre et des Affaires Foncières) utilisent des \textbf{stations totales}. Ces appareils mesurent au laser des distances (jusqu'à plusieurs km) et des angles (horizontal et vertical) avec une précision au millimètre et à la seconde d'arc. Le calculateur embarqué résout en temps réel des milliers de triangles rectangles (et sphériques) en 3D pour donner les coordonnées $(X, Y, Z)$ des bornes foncières. La trigonométrie, c'est leur quotidien !
\end{savaistu}

\courssub{Exemple résolu complet : Le fromager de la cour de l'école}

\begin{exoresolu}[Hauteur d'un fromager]
\textbf{Énoncé :} Pour estimer la hauteur d'un fromager dans la cour, un élève mesure la longueur de l'ombre au sol : $12,0\,\text{m}$. À ce moment, l'angle d'élévation du soleil est de $40^\circ$. Quelle est la hauteur de l'arbre ? (Donner l'arrondi au décimètre).

\tcblower
\textbf{Solution :}

\underline{\textit{1. Schéma et inventaire}}
\begin{itemize}
    \item Triangle rectangle : sommet soleil (S), pied arbre (P), cime (C). Angle droit en P.
    \item Données : $PC = 12,0\,\text{m}$ (côté adjacent à l'angle $40^\circ$), $\widehat{SPC}=40^\circ$.
    \item Cherché : $h = PS$ (côté opposé à l'angle $40^\circ$).
\end{itemize}

\underline{\textit{2. Choix du rapport trigonométrique}}
L'angle connu est $40^\circ$. On a l'adjacent ($12\,\text{m}$) et on cherche l'opposé ($h$).
\cle{Tangente} : $\tan(40^\circ) = \dfrac{\text{opposé}}{\text{adjacent}} = \dfrac{h}{12}$.

\underline{\textit{3. Calcul}}
\[
h = 12 \times \tan(40^\circ)
\]
Sur calculatrice (mode DEGRÉ) : $12 \times \tan(40) \approx 10,069...$

\underline{\textit{4. Vérification}}
$\tan(40^\circ) \approx 0,839$. $12 \times 0,839 \approx 10,07$. Cohérent.

\underline{\textit{5. Réponse}}
La hauteur du fromager est d'environ \textbf{10,1 m} (arrondi au décimètre).
\end{exoresolu}

\courssub{Arbre de décision : comment choisir son outil ?}

\begin{popfigure}
\begin{tikzpicture}[
    decision/.style={diamond, draw=popblue, thick, fill=popblue!10, text width=3.5cm, align=center, inner sep=2pt},
    process/.style={rectangle, draw=popgreen, thick, fill=popgreen!10, text width=3.5cm, align=center, inner sep=2pt, rounded corners},
    start/.style={ellipse, draw=poporange, thick, fill=poporange!20, text width=2.5cm, align=center},
    arrow/.style={->, thick, popdark},
    node distance=1.8cm and 2.5cm
]
\node[start] (start) {Données du problème};
\node[decision, below of=start] (q1) {Connais-tu un angle aigu ?};
\node[process, below left=1.5cm and 1cm of q1, align=center] (cas1){CAS 1 : 1 angle + 1 côté\\$\rightarrow$ Calculer 2ème angle ($90-\alpha$)\\$\rightarrow$ 2 côtés par trigonométrie};
\node[process, below right=1.5cm and 1cm of q1] (q2) {Connais-tu 2 côtés ?};
\node[process, below of=q2, align=center] (cas2){CAS 2 : 2 côtés\\$\rightarrow$ 3ème côté (Pythagore)\\$\rightarrow$ 2 angles (trigo inverse)};
\node[process, below=1.5cm of cas1, align=center] (fin1){Vérifier (somme angles, Pythagore)\\Phrase de réponse};
\node[process, below=1.5cm of cas2, align=center] (fin2){Vérifier (somme angles, Pythagore)\\Phrase de réponse};

\draw[arrow] (start) -- (q1);
\draw[arrow] (q1) -- node[left, near start] {\textbf{OUI}} (cas1);
\draw[arrow] (q1) -- node[right, near start] {\textbf{NON}} (q2);
\draw[arrow] (q2) -- node[right] {\textbf{OUI}} (cas2);
\draw[arrow] (cas1) -- (fin1);
\draw[arrow] (cas2) -- (fin2);
\end{tikzpicture}
\legende{Arbre de décision pour résoudre un triangle rectangle. Deux cas de figure exhaustifs au programme de 3ème.}
\end{popfigure}

\courssub{Exercices d'intégration (type BEPC)}

\begin{exercice}[Largeur du fleuve Wouri (modélisé)]
Depuis un point $A$ sur la rive, on repère un arbre $C$ sur l'autre rive. On mesure $AB = 50\,\text{m}$ le long de la rive (avec $B$ tel que $(AB) \perp (AC)$). L'angle $\widehat{ABC}$ mesure $32^\circ$.
Calculer la largeur du fleuve $AC$ (arrondi au mètre).
\end{exercice}

\begin{exercice}[Pente de la route de Dschang]
Un panneau indique une pente de $12\%$ sur 500 m de distance horizontale.
\begin{enumerate}
    \item Calculer l'angle d'inclinaison de la route (arrondi au dixième de degré).
    \item Calculer le dénivelé (différence d'altitude) sur ces 500 m.
\end{enumerate}
\end{exercice}

\begin{exercice}[Ombre du poteau CAMTEL]
Un poteau téléphonique vertical projette une ombre de $4,5\,\text{m}$ lorsque l'angle d'élévation du soleil est de $55^\circ$.
Calculer la hauteur du poteau (arrondi au centimètre).
\end{exercice}

\begin{corrige}[Exercice 1 : Largeur du fleuve]
\textbf{Données :} $AB=50\,\text{m}$, $\widehat{ABC}=32^\circ$, triangle $ABC$ rectangle en $A$.
\textbf{Cherché :} $AC$.
Dans $\triangle ABC$ rectangle en $A$, par rapport à l'angle $32^\circ$ en $B$ : $AB$ est adjacent, $AC$ est opposé.
$\tan(32^\circ) = \frac{AC}{50} \Rightarrow AC = 50 \times \tan(32^\circ) \approx 31,2\,\text{m}$.
\textbf{Réponse :} La largeur du fleuve est d'environ \textbf{31 m}.
\end{corrige}

\begin{corrige}[Exercice 2 : Pente de Dschang]
1. $\tan(\alpha) = \frac{12}{100} = 0,12 \Rightarrow \alpha = \arctan(0,12) \approx 6,8^\circ$.
2. Dénivelé $= 500 \times 0,12 = 60\,\text{m}$.
\textbf{Réponse :} L'angle est d'environ $6,8^\circ$ et le dénivelé est de $60\,\text{m}$.
\end{corrige}

\begin{corrige}[Exercice 3 : Poteau CAMTEL]
Triangle rectangle : poteau (hauteur $h$), ombre ($4,5\,\text{m}$), angle $55^\circ$ au bout de l'ombre.
$\tan(55^\circ) = \frac{h}{4,5} \Rightarrow h = 4,5 \times \tan(55^\circ) \approx 6,43\,\text{m}$.
\textbf{Réponse :} Le poteau mesure environ \textbf{6,43 m}.
\end{corrige}

\begin{aretenir}[Bilan de l'unité Trigonométrie]
\begin{itemize}
    \item Dans un triangle rectangle, $\cos = \frac{\text{adjacent}}{\text{hypoténuse}}$, $\sin = \frac{\text{opposé}}{\text{hypoténuse}}$, $\tan = \frac{\text{opposé}}{\text{adjacent}}$.
    \item Angles remarquables : $30^\circ, 45^\circ, 60^\circ$ (valeurs exactes à connaître).
    \item Pour trouver un côté : on utilise la relation directe ($\cos$, $\sin$, $\tan$).
    \item Pour trouver un angle : on utilise la relation inverse ($\arccos$, $\arcsin$, $\arctan$).
    \item Résoudre un triangle = trouver 3 côtés + 3 angles. Données minimales : (1 angle + 1 côté) ou (2 côtés).
    \item \textbf{Toujours :} Schéma $\rightarrow$ Inventaire $\rightarrow$ Choix angle $\rightarrow$ Calculs (mémoire calculatrice) $\rightarrow$ Vérification $\rightarrow$ Phrase réponse.
\end{itemize}
\end{aretenir}

\newpage

\coursec{Activité d'intégration}

\coursec{Trigonométrie dans le triangle rectangle}

\courssub{Repérage et vocabulaire dans le triangle rectangle}

Avant tout calcul, il faut savoir \emph{nommer} les côtés d'un triangle rectangle par rapport à un angle aigu choisi. C'est la clé pour ne pas se tromper de formule.

\begin{definition}[Côtés d'un triangle rectangle par rapport à un angle aigu]
Soit $ABC$ un triangle rectangle en $C$. On s'intéresse à l'angle aigu $\widehat{A}$ (ou $\widehat{B}$).
\begin{itemize}
\item L'\cle{hypoténuse} est le côté le plus long, \emph{toujours} opposé à l'angle droit (ici $[AB]$).
\item Le côté \cle{adjacent} à $\widehat{A}$ est le côté qui \emph{touche} l'angle $\widehat{A}$ et qui n'est pas l'hypoténuse (ici $[AC]$).
\item Le côté \cle{opposé} à $\widehat{A}$ est le côté qui ne touche \emph{pas} l'angle $\widehat{A}$ (ici $[BC]$).
\end{itemize}
\end{definition}

\begin{popfigure}
\begin{center}
\begin{tikzpicture}[scale=1.8]
\coordinate (A) at (0,0);
\coordinate (C) at (3.5,0);
\coordinate (B) at (3.5,2);
\draw[thick,popblue] (A)--(C)--(B)--cycle;
\draw[popdark] (3.2,0) rectangle (3.5,0.3);
% Étiquettes sommets
\node[below left] at (A) {$A$};
\node[below right] at (C) {$C$};
\node[above right] at (B) {$B$};
% Côtés
\node[below, popblue] at (1.75,0) {\textbf{Adjacent} $AC$};
\node[right, popblue] at (3.5,1) {\textbf{Opposé} $BC$};
\node[above left, popblue, rotate=30] at (1.75,1) {\textbf{Hypoténuse} $AB$};
% Angle A
\draw[poporange,thick] (0.3,0) arc (0:30:0.3);
\node[poporange] at (0.5,0.1) {$\widehat{A}$};
% Angle B
\draw[poporange,thick] (3.5,1.7) arc (90:120:0.3);
\node[poporange] at (3.2,1.9) {$\widehat{B}$};
\end{tikzpicture}
\end{center}
\legende{Repérage des côtés par rapport à l'angle $\widehat{A}$. Pour $\widehat{B}$, les rôles d'adjacent et d'opposé s'inversent.}
\end{popfigure}

\begin{methode}[Bien repérer les côtés (3 étapes)]
\begin{enumerate}
\item Tracer (ou visualiser) le triangle rectangle et marquer l'angle droit.
\item Identifier l'angle aigu de référence (celui dont on connaît la mesure ou qu'on cherche).
\item Nommer : \textbf{Hypoténuse} (face à l'angle droit) $\to$ \textbf{Adjacent} (collé à l'angle) $\to$ \textbf{Opposé} (le reste).
\end{enumerate}
\end{methode}

\begin{attention}[Piège classique]
Ne confonds pas \emph{adjacent} et \emph{opposé} ! L'adjacent \textbf{touche} l'angle (à côté de lui), l'opposé est \textbf{loin} de l'angle (en face de lui). L'hypoténuse ne change jamais de place.
\end{attention}

\courssub{Définition des rapports trigonométriques : Cosinus, Sinus, Tangente}

Dans tout triangle rectangle, les rapports de longueurs ne dépendent \emph{que} de l'angle, pas de la taille du triangle. C'est la propriété fondamentale de la trigonométrie.

\begin{propriete}[Définitions du cosinus, sinus, tangente]
Dans un triangle $ABC$ rectangle en $C$, pour l'angle aigu $\widehat{A}$ :
\[
\cos \widehat{A} = \dfrac{\text{côté adjacent}}{\text{hypoténuse}} = \dfrac{AC}{AB}, \qquad
\sin \widehat{A} = \dfrac{\text{côté opposé}}{\text{hypoténuse}} = \dfrac{BC}{AB}, \qquad
\tan \widehat{A} = \dfrac{\text{côté opposé}}{\text{côté adjacent}} = \dfrac{BC}{AC}.
\]
\begin{itemize}
\item Ces trois nombres sont \textbf{sans unité} (quotient de deux longueurs).
\item Pour tout angle aigu $\alpha$ : $0 < \cos \alpha < 1$, $0 < \sin \alpha < 1$, $\tan \alpha > 0$.
\item Relation fondamentale : $\tan \alpha = \dfrac{\sin \alpha}{\cos \alpha}$ et $\cos^2\alpha + \sin^2\alpha = 1$ (admis en 3ème, démontré en 2nde).
\end{itemize}
\end{propriete}

\begin{aretenir}[Mémo : SOH - CAH - TOA]
Pour retenir les formules, utilise la mnémotechnique :
\begin{center}
\begin{tabular}{c c c}
\textbf{S}inus = \textbf{O}pposé / \textbf{H}ypoténuse & \textbf{C}osinus = \textbf{A}djacent / \textbf{H}ypoténuse & \textbf{T}angente = \textbf{O}pposé / \textbf{A}djacent \\
\emph{SOH} & \emph{CAH} & \emph{TOA}
\end{tabular}
\end{center}
\end{aretenir}

\begin{exemplebox}[Calculer les rapports pour un triangle donné]
Dans un triangle $ABC$ rectangle en $C$, on a $AB = 10$ cm, $AC = 6$ cm, $BC = 8$ cm. Calculer $\cos \widehat{A}$, $\sin \widehat{A}$, $\tan \widehat{A}$.
\begin{align*}
\cos \widehat{A} &= \frac{AC}{AB} = \frac{6}{10} = 0{,}6 \\
\sin \widehat{A} &= \frac{BC}{AB} = \frac{8}{10} = 0{,}8 \\
\tan \widehat{A} &= \frac{BC}{AC} = \frac{8}{6} = \frac{4}{3} \approx 1{,}33
\end{align*}
\emph{Vérification :} $\cos^2 + \sin^2 = 0,36 + 0,64 = 1$. $\frac{\sin}{\cos} = \frac{0,8}{0,6} = \frac{4}{3} = \tan$. C'est cohérent.
\end{exemplebox}

\courssub{Angles remarquables (30°, 45°, 60°) et valeurs exactes}

Certains angles donnent des valeurs exactes (souvent avec des racines carrées) qu'il faut connaître par cœur pour le BEPC. Elles permettent de calculer \emph{sans calculatrice}.

\begin{experience}[Découverte des valeurs exactes]
\begin{enumerate}
\item \textbf{Angle de 45° :} Dans un triangle rectangle \emph{isocèle} (côtés de l'angle droit égaux), les angles aigus valent 45°. Si les côtés de l'angle droit valent 1, l'hypoténuse vaut $\sqrt{2}$ (Pythagore). En déduire $\cos 45^{\circ}$, $\sin 45^{\circ}$, $\tan 45^{\circ}$.
\item \textbf{Angles de 30° et 60° :} Dans un triangle équilatéral de côté 2, la hauteur coupe l'angle en deux (30°) et le côté en deux (1). La hauteur vaut $\sqrt{3}$ (Pythagore : $2^2 = 1^2 + h^2$). En déduire les valeurs pour 30° et 60°.
\end{enumerate}
\end{experience}

\begin{aretenir}[Valeurs exactes à connaître par cœur]
\begin{center}
\begin{tabular}{|c|c|c|c|}\hline
Angle $\alpha$ & $\cos \alpha$ & $\sin \alpha$ & $\tan \alpha$ \\ \hline
$30^{\circ}$ & $\dfrac{\sqrt{3}}{2}$ & $\dfrac{1}{2}$ & $\dfrac{\sqrt{3}}{3}$ \\ \hline
$45^{\circ}$ & $\dfrac{\sqrt{2}}{2}$ & $\dfrac{\sqrt{2}}{2}$ & $1$ \\ \hline
$60^{\circ}$ & $\dfrac{1}{2}$ & $\dfrac{\sqrt{3}}{2}$ & $\sqrt{3}$ \\ \hline
\end{tabular}
\end{center}
\begin{itemize}
\item Astuce : $\cos 30^{\circ} = \sin 60^{\circ}$, $\sin 30^{\circ} = \cos 60^{\circ}$. Le cosinus \emph{décroît}, le sinus \emph{croît} quand l'angle augmente de $0^{\circ}$ à $90^{\circ}$.
\item $\tan 45^{\circ} = 1$ car opposé = adjacent (triangle isocèle).
\end{itemize}
\end{aretenir}

\begin{attention}[Calculatrice vs Valeurs exactes]
Si l'énoncé demande une \textbf{valeur exacte}, laisse les racines ($\sqrt{2}, \sqrt{3}$) et les fractions. N'utilise la calculatrice que si on demande un \textbf{arrondi} (ex: « arrondir au dixième »).
\end{attention}

\courssub{Calculer une longueur dans un triangle rectangle}

On connaît un angle aigu et une longueur. On cherche une autre longueur.

\begin{methode}[Stratégie pour trouver une longueur]
\begin{enumerate}
\item \textbf{Repérer} : Triangle rectangle, angle de référence, côtés (Hyp, Adj, Opp) par rapport à cet angle.
\item \textbf{Choisir} le bon rapport : celui qui relie le \textbf{côté connu}, le \textbf{côté cherché} et l'\textbf{angle connu}.
\item \textbf{Écrire} la formule, \textbf{isoler} l'inconnue.
\item \textbf{Calculer} (calculatrice en mode \textbf{DEGRÉS}).
\item \textbf{Conclure} avec l'unité.
\end{enumerate}
\end{methode}

\begin{exoresolu}[Calcul de longueur — Toit d'une case]
Le toit d'une case forme un triangle rectangle $ABC$ en $C$. L'angle $\widehat{A} = 30^{\circ}$ et la hauteur du mur $BC = 3$ m. Calculer la longueur de la pente de toit $AB$ (hypoténuse).
\tcblower
\textbf{Repérage :} $\widehat{A}=30^{\circ}$, $BC$ (opposé) $= 3$ m, $AB$ (hypoténuse) $= ?$
\textbf{Choix :} $\sin 30^{\circ} = \dfrac{\text{opposé}}{\text{hypoténuse}} = \dfrac{BC}{AB}$.
\textbf{Calcul :} $\sin 30^{\circ} = \dfrac{1}{2} = \dfrac{3}{AB} \quad \Rightarrow \quad AB = 3 \times 2 = 6$ m.
\textbf{Conclusion :} La pente de toit mesure \textbf{6 m} (valeur exacte).
\end{exoresolu}

\begin{exoresolu}[Calcul de longueur — Pente de route]
Une route monte avec un angle de $12^{\circ}$ sur une distance horizontale de $500$ m. Quelle est la dénivellation (hauteur gagnée) ?
\tcblower
\textbf{Repérage :} Angle $12^{\circ}$, Adjacent $= 500$ m, Opposé $= ?$
\textbf{Choix :} $\tan 12^{\circ} = \dfrac{\text{opposé}}{\text{adjacent}}$.
\textbf{Calcul :} $\text{Opposé} = 500 \times \tan 12^{\circ} \approx 500 \times 0{,}2126 \approx 106{,}3$ m.
\textbf{Conclusion :} La dénivellation est d'environ \textbf{106,3 m}.
\end{exoresolu}

\courssub{Calculer un angle aigu dans un triangle rectangle}

On connaît deux longueurs. On cherche la mesure de l'angle.

\begin{methode}[Stratégie pour trouver un angle]
\begin{enumerate}
\item \textbf{Repérer} : Triangle rectangle, côtés connus (Hyp, Adj, Opp) par rapport à l'angle cherché.
\item \textbf{Choisir} le rapport qui utilise les \textbf{deux côtés connus}.
\item \textbf{Calculer} la valeur du rapport (division).
\item \textbf{Utiliser la touche réciproque} de la calculatrice : $\cos^{-1}$, $\sin^{-1}$ ou $\tan^{-1}$ (souvent \texttt{SHIFT} + \texttt{cos/sin/tan}).
\item \textbf{Vérifier} que la calculatrice est en mode \textbf{DEGRÉS} (affichage \texttt{D} ou \texttt{DEG}).
\item \textbf{Arrondir} selon la consigne (souvent au degré près ou au dixième de degré) et \textbf{conclure}.
\end{enumerate}
\end{methode}

\begin{exoresolu}[Calcul d'angle — Escalier]
Un escalier a une marche de hauteur $17$ cm (contremarche) et une profondeur de $28$ cm (giron). Quel est l'angle de l'escalier par rapport à l'horizontale ?
\tcblower
\textbf{Repérage :} Triangle rectangle. Opposé $= 17$ cm, Adjacent $= 28$ cm. Angle $\alpha$ cherché.
\textbf{Choix :} $\tan \alpha = \dfrac{17}{28}$.
\textbf{Calcul :} $\alpha = \tan^{-1}\left(\dfrac{17}{28}\right) \approx \tan^{-1}(0{,}6071) \approx 31{,}3^{\circ}$.
\textbf{Conclusion :} L'escalier forme un angle d'environ \textbf{31,3°} avec l'horizontale.
\end{exoresolu}

\begin{attention}[Erreurs fréquentes sur la calculatrice]
\begin{itemize}
\item Oublier de passer en mode \textbf{DEGRÉS} (résultat en radians ou grades, totalement faux).
\item Faire la division \emph{après} avoir appuyé sur $\tan^{-1}$ au lieu d'avant. \textbf{Ordre :} Division $\to$ $\tan^{-1}$.
\item Arrondir le résultat de la division \emph{avant} de faire $\tan^{-1}$ (perte de précision). Garde la fraction ou la valeur exacte de la calculatrice.
\end{itemize}
\end{attention}

\courssub{Résolution de triangle rectangle \& Problèmes concrets (Intégration)}

\courssub{Objectifs de l'activité}

Cette activité d'intégration clôt la leçon sur la trigonométrie dans le triangle rectangle. En travaillant en binôme sur un problème réel d'aménagement, tu vas mobiliser \emph{ensemble} tous les savoirs de l'unité :

\begin{itemize}
\item repérer, dans une situation concrète, un \cle{triangle rectangle} et nommer correctement ses côtés (\cle{hypoténuse}, côté \cle{adjacent}, côté \cle{opposé}) ;
\item choisir le bon \cle{rapport trigonométrique} (sinus, cosinus, tangente) selon les données ;
\item calculer la \cle{longueur d'un côté} connaissant un angle et un autre côté ;
\item calculer la \cle{mesure d'un angle aigu} connaissant deux côtés (à l'aide de la calculatrice) ;
\item modéliser, rédiger une démarche complète et argumenter une conclusion, comme on te le demandera au BEPC.
\end{itemize}

\courssub{Situation}

\textbf{L'accès au dispensaire de Bikok.}

Le village de Bikok, dans la région du Centre, vient de voir son dispensaire rénové grâce au financement de la commune. Mais un problème demeure : pour entrer dans le bâtiment, il faut monter une \cle{marche} de \SI{35}{cm} de haut. Les personnes âgées, les femmes enceintes et surtout les personnes en fauteuil roulant ne peuvent pas accéder seules au dispensaire.

Le maire de la commune souhaite faire construire une \cle{rampe d'accès} en béton. Il consulte la norme d'accessibilité, qui impose une condition :

\begin{center}
\fbox{\textbf{La pente de la rampe ne doit pas dépasser 10 \%.}}
\end{center}

\begin{definition}[Pente d'une rampe]
La \cle{pente} d'une rampe est le quotient de la dénivellation (la hauteur montée) par la distance horizontale parcourue :
\[ \text{pente} = \dfrac{\text{dénivellation}}{\text{distance horizontale}}. \]
Elle s'exprime souvent en pourcentage. Une pente de 10 \% signifie que l'on monte de \SI{10}{cm} lorsqu'on avance horizontalement de \SI{100}{cm}. Dans le triangle rectangle modélisant la rampe, la pente est exactement la \cle{tangente} de l'angle que fait la rampe avec le sol horizontal.
\end{definition}

On modélise la situation par un triangle $ABC$ rectangle en $C$ :
\begin{itemize}
\item $A$ est le pied de la rampe au sol, $B$ le haut de la marche (l'entrée du dispensaire) ;
\item $BC = \SI{35}{cm} = \SI{0,35}{m}$ est la hauteur de la marche (côté opposé à l'angle $\widehat{A}$) ;
\item $AC$ est l'emprise au sol de la rampe (côté adjacent à l'angle $\widehat{A}$) ;
\item $AB$ est la longueur de la rampe elle-même (hypoténuse).
\end{itemize}

\begin{popfigure}
\begin{center}
\begin{tikzpicture}[scale=1.5]
\coordinate (A) at (0,0);
\coordinate (C) at (5,0);
\coordinate (B) at (5,0.5);
\draw[thick,popblue] (A)--(C)--(B)--cycle;
\draw[popdark] (4.7,0) rectangle (5,0.3);
\draw[pattern=north east lines,pattern color=popmuted] (5,0) rectangle (5.6,0.5);
\node[right] at (5.6,0.25) {\footnotesize mur du dispensaire};
\node[below left] at (A) {$A$};
\node[below right] at (C) {$C$};
\node[above right] at (B) {$B$};
\node[below] at (2.5,0) {\footnotesize emprise au sol $AC$};
\node[right] at (5,0.25) {\footnotesize $BC = 0,35$ m};
\node[above left,rotate=5.7] at (2.5,0.25) {\footnotesize rampe $AB$};
\draw[poporange,thick] (0.8,0) arc (0:5.71:0.8);
\node[poporange] at (1.1,0.08) {\footnotesize $\widehat{A}$};
\end{tikzpicture}
\end{center}
\legende{Modélisation de la rampe par un triangle rectangle en $C$ (schéma non à l'échelle, angle $\widehat{A} \approx 5,7^{\circ}$).}
\end{popfigure}

Le maire dispose devant le dispensaire d'une cour dont la profondeur disponible est de \SI{3,20}{m} seulement. Ta classe de 3ème est chargée de l'étude technique.

\courssub{Consignes}

Travail en binôme. Chaque question doit être rédigée soigneusement : figure, données, formule utilisée, calcul, conclusion avec unité.

\begin{enumerate}
\item \textbf{Angle maximal autorisé.} On note $\alpha$ l'angle $\widehat{A}$ que fait la rampe avec le sol. Justifier que $\tan \alpha = \text{pente}$, puis déterminer, à l'aide de la calculatrice, la mesure maximale de $\alpha$ autorisée par la norme (arrondir au dixième de degré).
\item \textbf{Emprise au sol minimale.} Pour une pente exactement égale à 10 \%, calculer la longueur $AC$ de l'emprise au sol nécessaire pour franchir la marche de \SI{0,35}{m}. La cour de \SI{3,20}{m} suffit-elle ?
\item \textbf{Longueur de la rampe.} En utilisant l'angle $\alpha$ trouvé à la question 1 et le cosinus (ou le sinus), calculer la longueur $AB$ de la rampe (arrondir au centimètre). Vérifier le résultat avec le théorème de Pythagore.
\item \textbf{Une solution alternative.} Un entrepreneur propose une rampe moins pentue, à 7 \%, mais la cour est trop petite. Il imagine alors une rampe en deux tronçons identiques séparés par un palier horizontal de \SI{1,20}{m} (rampe en « U » le long du mur). Pour chaque tronçon montant de \SI{0,175}{m} avec une pente de 7 \% : calculer l'emprise au sol de chaque tronçon, puis la longueur totale au sol du dispositif (deux tronçons + palier). Cette solution tient-elle dans une cour de \SI{3,20}{m} de profondeur si les tronçons sont disposés le long de deux murs perpendiculaires ?
\item \textbf{Rapport technique.} Rédiger un court rapport (une demi-page) adressé au maire : schéma coté, résultats trigonométriques justifiés, solution retenue et argumentée. Préparer une présentation orale de 3 minutes.
\end{enumerate}

\courssub{Ressources à mobiliser}

\begin{aretenir}[Boîte à outils de la leçon]
Dans un triangle $ABC$ rectangle en $C$, pour l'angle aigu $\widehat{A}$ :
\[ \cos \widehat{A} = \dfrac{\text{côté adjacent}}{\text{hypoténuse}} = \dfrac{AC}{AB}, \qquad
\sin \widehat{A} = \dfrac{\text{côté opposé}}{\text{hypoténuse}} = \dfrac{BC}{AB}, \qquad
\tan \widehat{A} = \dfrac{\text{côté opposé}}{\text{côté adjacent}} = \dfrac{BC}{AC}. \]
\begin{itemize}
\item Pour trouver une \textbf{longueur} : choisir le rapport qui relie la longueur connue, la longueur cherchée et l'angle connu, puis isoler l'inconnue.
\item Pour trouver un \textbf{angle} : calculer le rapport de deux côtés, puis utiliser la touche réciproque de la calculatrice ($\tan^{-1}$, $\sin^{-1}$ ou $\cos^{-1}$, en mode degrés).
\item Le théorème de Pythagore ($AB^2 = AC^2 + BC^2$) sert de vérification.
\end{itemize}
\end{aretenir}

\begin{attention}[Réflexes à ne pas oublier]
\begin{itemize}
\item Vérifier que la calculatrice est en mode \textbf{degrés}.
\item Convertir toutes les longueurs dans la \textbf{même unité} avant de calculer (ici : le mètre).
\item Ne pas arrondir les résultats intermédiaires : garder les valeurs exactes ou les valeurs affichées par la calculatrice, et n'arrondir qu'à la fin.
\item Toujours conclure par une phrase répondant à la question posée, avec l'unité.
\end{itemize}
\end{attention}

\courssub{Démarche et corrigé commenté}

\begin{exoresolu}[L'accès au dispensaire de Bikok — corrigé complet]
Résoudre les questions 1 à 4 de l'activité.
\tcblower
\textbf{Question 1 : angle maximal autorisé.}

Dans le triangle $ABC$ rectangle en $C$, on a par définition :
\[ \tan \widehat{A} = \dfrac{BC}{AC} = \dfrac{\text{dénivellation}}{\text{distance horizontale}} = \text{pente}. \]
La norme impose une pente maximale de 10 \%, donc $\tan \alpha \leq 0{,}1$. La valeur limite vérifie $\tan \alpha = 0{,}1$. À la calculatrice (mode degrés) :
\[ \alpha = \tan^{-1}(0{,}1) \approx 5{,}7^{\circ}. \]
\textbf{Conclusion :} l'angle de la rampe avec le sol ne doit pas dépasser environ $5{,}7^{\circ}$.

\medskip
\textbf{Question 2 : emprise au sol minimale.}

On connaît le côté opposé $BC = 0{,}35$ m et l'angle $\alpha$ avec $\tan \alpha = 0{,}1$ ; on cherche le côté adjacent $AC$. On utilise la tangente :
\[ \tan \alpha = \dfrac{BC}{AC} \quad \Longleftrightarrow \quad AC = \dfrac{BC}{\tan \alpha} = \dfrac{0{,}35}{0{,}1} = 3{,}5 \text{ m}. \]
\textbf{Conclusion :} il faut une emprise au sol de \SI{3,5}{m}. Or la cour ne fait que \SI{3,20}{m} de profondeur : $3{,}5 > 3{,}2$, donc \textbf{la cour ne suffit pas} pour une rampe droite à 10 \%. Il faudra une solution alternative.

\medskip
\textbf{Question 3 : longueur de la rampe.}

On connaît $BC = 0{,}35$ m (opposé) et on cherche l'hypoténuse $AB$ : on utilise le sinus, avec $\alpha \approx 5{,}71^{\circ}$ (valeur non arrondie conservée à la calculatrice) :
\[ \sin \alpha = \dfrac{BC}{AB} \quad \Longleftrightarrow \quad AB = \dfrac{BC}{\sin \alpha} = \dfrac{0{,}35}{\sin(5{,}71^{\circ})} \approx \dfrac{0{,}35}{0{,}0995} \approx 3{,}52 \text{ m}. \]
\emph{Vérification par Pythagore :}
\[ AC^2 + BC^2 = 3{,}5^2 + 0{,}35^2 = 12{,}25 + 0{,}1225 = 12{,}3725 \quad \text{et} \quad \sqrt{12{,}3725} \approx 3{,}52 \text{ m} = AB. \]
Les deux méthodes donnent le même résultat : la rampe mesure environ \SI{3,52}{m}.

\medskip
\textbf{Question 4 : rampe en deux tronçons à 7 \%.}

Pour chaque tronçon, la dénivellation est $0{,}35 \div 2 = 0{,}175$ m et la pente est $0{,}07$. L'emprise au sol d'un tronçon est :
\[ AC_1 = \dfrac{0{,}175}{0{,}07} = 2{,}5 \text{ m}. \]
(Au passage, l'angle de chaque tronçon vaut $\tan^{-1}(0{,}07) \approx 4{,}0^{\circ}$, plus doux que la norme : c'est encore mieux pour les usagers.)

Le dispositif complet comporte deux tronçons de \SI{2,5}{m} et un palier de \SI{1,20}{m}. Comme les tronçons sont disposés le long de deux murs \emph{perpendiculaires} (rampe en « U »), la profondeur occupée dans la cour n'est que celle d'un seul tronçon, soit \SI{2,5}{m}, la largeur occupée le long de l'autre mur étant $2{,}5 + 1{,}2 = 3{,}7$ m le long du mur (espace disponible supposé suffisant).
\[ 2{,}5 < 3{,}2 \quad \Rightarrow \quad \textbf{la solution en « U » tient dans la cour.} \]
La longueur de chaque tronçon incliné est $\dfrac{0{,}175}{\sin(4{,}0^{\circ})} \approx 2{,}51$ m, soit environ \SI{5,02}{m} de rampe inclinée au total, plus le palier.

\medskip
\textbf{Question 5 (éléments de rapport).} Le rapport au maire doit contenir : le schéma coté du triangle $ABC$ ; le calcul de l'angle maximal ($5{,}7^{\circ}$) ; le constat qu'une rampe droite exige \SI{3,5}{m} au sol, incompatible avec la cour de \SI{3,20}{m} ; la proposition de la rampe en « U » à 7 \% (emprise \SI{2,5}{m} par tronçon, palier de \SI{1,20}{m}), plus confortable et conforme ; enfin une conclusion argumentée recommandant cette solution.
\end{exoresolu}

\courssub{Bilan de l'activité}

Cette activité t'a permis de mettre en évidence plusieurs idées essentielles de la leçon :

\begin{itemize}
\item \textbf{La trigonométrie modélise le réel.} Une rampe, un escalier, un toit, une route en côte : dès qu'une situation fait intervenir une hauteur, une distance horizontale et un angle, on peut la ramener à un triangle rectangle et la traiter avec $\cos$, $\sin$, $\tan$.
\item \textbf{Le choix du rapport est la clé.} On choisit sinus, cosinus ou tangente selon les côtés connus et cherchés : ici la tangente pour relier pente, hauteur et emprise ; le sinus pour obtenir la longueur de la rampe.
\item \textbf{On sait calculer dans les deux sens.} Connaissant deux côtés, on retrouve l'angle avec les touches réciproques de la calculatrice ($\tan^{-1}(0{,}1) \approx 5{,}7^{\circ}$) ; connaissant un angle et un côté, on retrouve les autres côtés.
\item \textbf{Pythagore reste un outil de vérification} précieux pour contrôler un résultat obtenu par la trigonométrie.
\item \textbf{Les mathématiques servent à décider.} Le calcul a montré qu'une solution (rampe droite) était impossible et a permis d'en justifier une autre (rampe en « U ») : rédiger une démarche et argumenter une conclusion sont des compétences exigées au BEPC, comme dans la vie professionnelle.
\end{itemize}

\begin{savaistu}[Les rampes autour de nous]
Les normes d'accessibilité des bâtiments publics (écoles, hôpitaux, mairies) imposent des pentes maximales, souvent entre 5 \% et 10 \%, précisément pour qu'un fauteuil roulant puisse monter sans danger. Les ingénieurs et les maçons utilisent quotidiennement la tangente d'un angle pour tracer rampes, escaliers et toitures. À Yaoundé comme à Bikok, un simple triangle rectangle suffit à rendre un bâtiment accessible à tous.
\end{savaistu}

\newpage

\coursec{Méthodes et exercices résolus}

\coursec{Trigonométrie dans le triangle rectangle}

\courssub{Repérage et vocabulaire dans le triangle rectangle}

\begin{definition}[Triangle rectangle]
Un triangle est \textbf{rectangle} s'il possède un angle droit ($90^{\circ}$). Le côté opposé à l'angle droit s'appelle l'\textbf{hypoténuse} ; c'est toujours le \emph{plus grand} côté. Les deux autres côtés sont les \textbf{côtés de l'angle droit}.
\end{definition}

\begin{experience}[Repérer les côtés par rapport à un angle aigu]
Dans un triangle rectangle $ABC$ rectangle en $A$, considérons l'angle aigu $\widehat{B}$.
\begin{itemize}
\item L'\textbf{hypoténuse} est $[BC]$ (côté opposé à l'angle droit). Elle ne change jamais.
\item Le \textbf{côté opposé} à $\widehat{B}$ est $[AC]$ (il ne touche pas à l'angle $\widehat{B}$).
\item Le \textbf{côté adjacent} à $\widehat{B}$ est $[AB]$ (il touche à l'angle $\widehat{B}$ et n'est pas l'hypoténuse).
\end{itemize}
\emph{Exercice :} Sur la figure ci-dessous, repère l'hypoténuse, le côté opposé et le côté adjacent par rapport à l'angle $\widehat{C}$.
\end{experience}

\begin{popfigure}
\begin{tikzpicture}[scale=1.2, >=Stealth]
\coordinate (A) at (0,0);
\coordinate (B) at (4,0);
\coordinate (C) at (0,3);
\draw[popdark, thick] (A) -- node[below] {$c$} (B) -- node[right] {$a$} (C) -- node[left] {$b$} cycle;
\draw[popblue, thick] (A) rectangle ++(0.3,0.3); % angle droit
% Angles
\draw[poporange, thick] (B) ++(180:0.6) arc (180:180-36.87:0.6) node[midway, left=2pt] {$\widehat{B}$};
\draw[poporange, thick] (C) ++(-90:0.6) arc (-90:-90+53.13:0.6) node[midway, right=2pt] {$\widehat{C}$};
% Labels
\node[below left] at (A) {$A$};
\node[below right] at (B) {$B$};
\node[above left] at (C) {$C$};
% Annotations for angle B
\draw[<->, popgreen, thick] (0.2, -0.5) -- node[below] {Adjacent à $\widehat{B}$} (3.8, -0.5);
\draw[<->, poppink, thick] (-0.5, 0.2) -- node[left] {Opposé à $\widehat{B}$} (-0.5, 2.8);
\draw[<->, popblue, thick] (4.2, 0) -- node[right] {Hypoténuse} (4.2, 3);
\end{tikzpicture}
\legende{Repérage des côtés par rapport à l'angle $\widehat{B}$ dans le triangle $ABC$ rectangle en $A$.}
\end{popfigure}

\begin{aretenir}[Vocabulaire essentiel]
\begin{itemize}
\item \textbf{Hypoténuse} : côté opposé à l'angle droit (le plus long).
\item \textbf{Côté opposé} : côté qui ne touche pas l'angle aigu considéré.
\item \textbf{Côté adjacent} : côté qui touche l'angle aigu considéré (et qui n'est pas l'hypoténuse).
\item \textbf{Angle aigu} : angle strictement compris entre $0^{\circ}$ et $90^{\circ}$.
\end{itemize}
\end{aretenir}

\courssub{Définition des rapports trigonométriques : Cosinus, Sinus, Tangente}

\begin{propriete}[Rapports trigonométriques dans un triangle rectangle]
Soit $ABC$ un triangle rectangle en $A$ et $\widehat{B}$ un de ses angles aigus. On définit trois rapports de longueurs :
\[
\cos \widehat{B} = \dfrac{\text{côté adjacent}}{\text{hypoténuse}} = \dfrac{AB}{BC}, \qquad
\sin \widehat{B} = \dfrac{\text{côté opposé}}{\text{hypoténuse}} = \dfrac{AC}{BC}, \qquad
\tan \widehat{B} = \dfrac{\text{côté opposé}}{\text{côté adjacent}} = \dfrac{AC}{AB}.
\]
Ces rapports ne dépendent \textbf{que} de la mesure de l'angle $\widehat{B}$, pas de la taille du triangle.
\end{propriete}

\begin{aretenir}[Moyen mnémotechnique : \cle{SOH-CAH-TOA}]
\begin{center}
\begin{tabular}{|c|c|c|}
\hline
\rowcolor{popblueL} \textbf{S}inus & \textbf{C}osinus & \textbf{T}angente \\
\hline
$\sin = \dfrac{\textbf{O}pposé}{\textbf{H}ypoténuse}$ & $\cos = \dfrac{\textbf{A}djacent}{\textbf{H}ypoténuse}$ & $\tan = \dfrac{\textbf{O}pposé}{\textbf{A}djacent}$ \\
\hline
\end{tabular}
\end{center}
\emph{Astuce :} \textbf{SOH} - \textbf{CAH} - \textbf{TOA} (se lit « So-Ca-Toa »).
\end{aretenir}

\begin{exemplebox}[Calculer les rapports dans un triangle donné]
Soit $DEF$ un triangle rectangle en $D$ tel que $DE = 3$~cm, $DF = 4$~cm et $EF = 5$~cm.
Par rapport à $\widehat{E}$ : adjacent $= DE = 3$, opposé $= DF = 4$, hypoténuse $= EF = 5$.
\[
\cos \widehat{E} = \frac{3}{5} = 0{,}6, \quad
\sin \widehat{E} = \frac{4}{5} = 0{,}8, \quad
\tan \widehat{E} = \frac{4}{3} \approx 1{,}33.
\]
Par rapport à $\widehat{F}$ : adjacent $= DF = 4$, opposé $= DE = 3$, hypoténuse $= EF = 5$.
\[
\cos \widehat{F} = \frac{4}{5} = 0{,}8, \quad
\sin \widehat{F} = \frac{3}{5} = 0{,}6, \quad
\tan \widehat{F} = \frac{3}{4} = 0{,}75.
\]
\emph{Remarque :} $\sin \widehat{E} = \cos \widehat{F}$ et $\cos \widehat{E} = \sin \widehat{F}$ car $\widehat{E}$ et $\widehat{F}$ sont complémentaires.
\end{exemplebox}

\courssub{Angles remarquables ($30^{\circ}$, $45^{\circ}$, $60^{\circ}$) et valeurs exactes}

\begin{experience}[D'où viennent ces valeurs ?]
\begin{enumerate}
\item \textbf{Angle $45^{\circ}$ :} Dans un carré de côté $1$, la diagonale mesure $\sqrt{2}$ (Pythagore). Le triangle isocèle rectangle a des angles $45^{\circ}$, $45^{\circ}$, $90^{\circ}$.
$\sin 45^{\circ} = \cos 45^{\circ} = \dfrac{1}{\sqrt{2}} = \dfrac{\sqrt{2}}{2}$, $\tan 45^{\circ} = 1$.
\item \textbf{Angles $30^{\circ}$ et $60^{\circ}$ :} Dans un triangle équilatéral de côté $2$, la hauteur coupe l'angle en deux ($30^{\circ}$) et le côté en deux ($1$). La hauteur vaut $\sqrt{3}$ (Pythagore : $2^2 = 1^2 + h^2$).
$\sin 30^{\circ} = \frac{1}{2}$, $\cos 30^{\circ} = \frac{\sqrt{3}}{2}$, $\tan 30^{\circ} = \frac{\sqrt{3}}{3}$.
$\sin 60^{\circ} = \frac{\sqrt{3}}{2}$, $\cos 60^{\circ} = \frac{1}{2}$, $\tan 60^{\circ} = \sqrt{3}$.
\end{enumerate}
\end{experience}

\begin{propriete}[Valeurs exactes des angles remarquables]
\begin{center}
\begin{tabular}{|c|c|c|c|}
\hline
\rowcolor{popblueL} Angle & $30^{\circ}$ & $45^{\circ}$ & $60^{\circ}$ \\
\hline
$\sin$ & $\dfrac{1}{2}$ & $\dfrac{\sqrt{2}}{2}$ & $\dfrac{\sqrt{3}}{2}$ \\
\hline
$\cos$ & $\dfrac{\sqrt{3}}{2}$ & $\dfrac{\sqrt{2}}{2}$ & $\dfrac{1}{2}$ \\
\hline
$\tan$ & $\dfrac{\sqrt{3}}{3}$ & $1$ & $\sqrt{3}$ \\
\hline
\end{tabular}
\end{center}
\emph{Astuce mémoire :} Pour le sinus, on écrit $\dfrac{\sqrt{1}}{2}$, $\dfrac{\sqrt{2}}{2}$, $\dfrac{\sqrt{3}}{2}$ pour $30^{\circ}$, $45^{\circ}$, $60^{\circ}$ ; le cosinus est la même liste dans l'ordre inverse.
\end{propriete}

\courssub{Calculer une longueur dans un triangle rectangle}

\begin{methode}[Utiliser les rapports trigonométriques pour calculer une longueur]
Pour calculer la longueur d'un côté d'un triangle rectangle lorsqu'on connaît un angle aigu et la longueur d'un autre côté :
\begin{enumerate}
\item \textbf{Nommer} le triangle rectangle et préciser où est l'angle droit (condition d'application).
\item \textbf{Repérer} par rapport à l'angle aigu connu : l'hypoténuse, le côté opposé et le côté adjacent.
\item \textbf{Choisir le bon rapport} grâce à \cle{SOH-CAH-TOA} : celui qui fait intervenir le côté \emph{connu} et le côté \emph{cherché}.
\item \textbf{Écrire l'égalité} avec les longueurs, puis isoler la longueur inconnue (produit en croix).
\item \textbf{Calculer} à la calculatrice (mode \emph{degrés} vérifié) et \textbf{conclure} par une phrase avec l'unité et l'arrondi demandé.
\end{enumerate}
\end{methode}

\begin{exoresolu}[Hauteur d'un pylône Orange]
Un technicien d'Orange Cameroun veut connaître la hauteur d'un pylône. Il se place à $25$~m du pied $P$ du pylône, au point $T$. L'angle sous lequel il voit le sommet $S$ est $\widehat{PTS} = 38^{\circ}$. Le triangle $PTS$ est rectangle en $P$. Calculer la hauteur $PS$ du pylône (arrondir au décimètre).
\tcblower
\textbf{Étape 1 :} Le triangle $PTS$ est rectangle en $P$ : on peut utiliser la trigonométrie.

\textbf{Étape 2 :} Par rapport à l'angle $\widehat{PTS} = 38^{\circ}$ :
\begin{itemize}
\item le côté \textbf{opposé} est $[PS]$ (longueur cherchée) ;
\item le côté \textbf{adjacent} est $[PT]$ avec $PT = 25$~m (longueur connue) ;
\item l'hypoténuse est $[TS]$ (inutile ici).
\end{itemize}

\textbf{Étape 3 :} Le rapport qui relie le côté opposé et le côté adjacent est la \cle{tangente} :
\[ \tan \widehat{PTS} = \dfrac{PS}{PT}. \]

\textbf{Étape 4 :} On isole $PS$ (produit en croix) :
\[ PS = PT \times \tan \widehat{PTS} = 25 \times \tan 38^{\circ}. \]

\textbf{Étape 5 :} À la calculatrice (mode degrés) : $\tan 38^{\circ} \approx 0{,}7812856...$, donc
\[ PS \approx 25 \times 0{,}7812856 \approx 19{,}532...~\text{m}. \]

\textbf{Conclusion :} le pylône mesure environ $\boxed{19{,}5~\text{m}}$ de hauteur (arrondi au décimètre).
\end{exoresolu}

\courssub{Calculer la mesure d'un angle aigu dans un triangle rectangle}

\begin{methode}[Calculer un angle connaissant deux côtés]
Pour calculer la mesure d'un angle aigu d'un triangle rectangle lorsqu'on connaît deux côtés :
\begin{enumerate}
\item Vérifier que le triangle est \textbf{rectangle} et l'écrire.
\item Repérer les deux côtés connus par rapport à l'angle cherché (opposé, adjacent, hypoténuse).
\item Choisir le rapport trigonométrique adapté (SOH-CAH-TOA) et écrire l'égalité (ex : $\cos \widehat{A} = \frac{\text{adjacent}}{\text{hypoténuse}}$).
\item Calculer la valeur numérique du rapport (fraction ou décimal).
\item Utiliser la touche \textbf{réciproque} de la calculatrice : $\cos^{-1}$ (ou Arccos), $\sin^{-1}$, $\tan^{-1}$, en mode degrés.
\item Conclure par une phrase, avec un arrondi au degré (ou au dixième) selon la consigne.
\end{enumerate}
\end{methode}

\begin{exoresolu}[Pente d'une rampe d'accès au marché de Mokolo]
Pour accéder au marché, une rampe $[AB]$ de $6$~m de long permet de franchir un dénivelé vertical $BC = 1{,}5$~m. Le triangle $ABC$ est rectangle en $C$. Calculer la mesure de l'angle $\widehat{BAC}$ que fait la rampe avec l'horizontale (arrondir au degré).
\tcblower
\textbf{Étape 1 :} Le triangle $ABC$ est rectangle en $C$ : la trigonométrie s'applique.

\textbf{Étape 2 :} Par rapport à l'angle $\widehat{BAC}$ (angle cherché) :
\begin{itemize}
\item l'\textbf{hypoténuse} est $[AB]$ (rampe) avec $AB = 6$~m (connu) ;
\item le côté \textbf{opposé} est $[BC]$ (dénivelé) avec $BC = 1{,}5$~m (connu) ;
\item le côté adjacent $[AC]$ n'est pas utilisé ici.
\end{itemize}

\textbf{Étape 3 :} Le rapport reliant l'opposé et l'hypoténuse est le \cle{sinus} :
\[ \sin \widehat{BAC} = \dfrac{BC}{AB} = \dfrac{1{,}5}{6} = 0{,}25. \]

\textbf{Étape 4 :} À la calculatrice (mode degrés) :
\[ \widehat{BAC} = \sin^{-1}(0{,}25) \approx 14{,}4775^{\circ}. \]

\textbf{Conclusion :} la rampe fait un angle d'environ $\boxed{14^{\circ}}$ avec l'horizontale (arrondi au degré).
\emph{Note :} Dans la version précédente, l'énoncé plaçait le dénivelé sur le côté adjacent, ce qui donnait un angle de $76^{\circ}$ (rampe quasi verticale). Ici, le dénivelé est l'opposé, ce qui donne une pente réaliste pour une rampe d'accès.
\end{exoresolu}

\courssub{Résolution de triangle rectangle \& Problèmes concrets (Intégration)}

\begin{methode}[Calculs exacts avec les angles remarquables ($30^{\circ}$, $45^{\circ}$, $60^{\circ}$)]
\begin{enumerate}
\item Identifier l'angle remarquable dans le triangle.
\item Remplacer le rapport trigonométrique par sa \textbf{valeur exacte} (tableau de la section 3).
\item Résoudre l'équation obtenue (produit en croix) et donner la réponse sous \textbf{forme exacte} (avec $\sqrt{2}$ ou $\sqrt{3}$), puis éventuellement une valeur approchée.
\end{enumerate}
\end{methode}

\begin{exoresolu}[Longueur exacte d'une ombre]
Un arbre vertical $[AB]$ de hauteur $AB = 8$~m projette une ombre $[BC]$ sur le sol horizontal. Les rayons du soleil font un angle de $60^{\circ}$ avec le sol : $\widehat{ACB} = 60^{\circ}$, et le triangle $ABC$ est rectangle en $B$. Déterminer la \textbf{valeur exacte} de la longueur $BC$ de l'ombre, puis une valeur approchée au centimètre.
\tcblower
\textbf{Étape 1 :} Le triangle $ABC$ est rectangle en $B$ : on applique la trigonométrie.

\textbf{Étape 2 :} Par rapport à l'angle $\widehat{ACB} = 60^{\circ}$ : le côté opposé est $[AB]$ ($AB = 8$~m, connu) et le côté adjacent est $[BC]$ (cherché).

\textbf{Étape 3 :} On utilise la \cle{tangente} :
\[ \tan 60^{\circ} = \dfrac{AB}{BC} \quad \Rightarrow \quad BC = \dfrac{AB}{\tan 60^{\circ}} = \dfrac{8}{\sqrt{3}}. \]

\textbf{Étape 4 :} On rationalise le dénominateur (forme exacte) :
\[ BC = \dfrac{8}{\sqrt{3}} = \dfrac{8\sqrt{3}}{3}~\text{m}. \]

Valeur approchée : $BC \approx \dfrac{8 \times 1{,}732}{3} \approx 4{,}62$~m.

\textbf{Conclusion :} la longueur exacte de l'ombre est $\boxed{\dfrac{8\sqrt{3}}{3}~\text{m}}$, soit environ $4{,}62$~m.
\end{exoresolu}

\begin{methode}[Résoudre complètement un triangle rectangle]
\emph{Résoudre} un triangle rectangle, c'est déterminer \textbf{toutes} les longueurs et \textbf{tous} les angles inconnus.
\begin{enumerate}
\item Faire une \textbf{figure à main levée} et y reporter toutes les données.
\item Pour les angles : utiliser le fait que les deux angles aigus sont \textbf{complémentaires} (somme $= 90^{\circ}$).
\item Pour les longueurs : utiliser un rapport trigonométrique (si un angle est connu) ou le \textbf{théorème de Pythagore} (si deux côtés sont connus).
\item Vérifier la cohérence : l'hypoténuse est le \emph{plus grand} côté ; un sinus ou un cosinus est entre $0$ et $1$.
\item Rédiger chaque étape avec une justification (nom de la propriété).
\end{enumerate}
\end{methode}

\begin{exoresolu}[Résolution complète]
Soit $EFG$ un triangle rectangle en $E$ tel que $FG = 10$~cm et $\widehat{EFG} = 35^{\circ}$. Résoudre ce triangle (arrondir les longueurs au millimètre).
\tcblower
\textbf{Calcul de l'angle $\widehat{EGF}$ :} Dans le triangle $EFG$ rectangle en $E$, les angles aigus sont complémentaires :
\[ \widehat{EGF} = 90^{\circ} - \widehat{EFG} = 90^{\circ} - 35^{\circ} = 55^{\circ}. \]

\textbf{Calcul de $EF$ :} Par rapport à $\widehat{EFG} = 35^{\circ}$, $[EF]$ est le côté adjacent et $[FG]$ l'hypoténuse. On utilise le cosinus :
\[ \cos 35^{\circ} = \dfrac{EF}{FG} \quad \Rightarrow \quad EF = 10 \times \cos 35^{\circ} \approx 10 \times 0{,}81915 \approx 8{,}19~\text{cm}. \]

\textbf{Calcul de $EG$ :} Par rapport à $\widehat{EFG}$, $[EG]$ est le côté opposé. On utilise le sinus :
\[ \sin 35^{\circ} = \dfrac{EG}{FG} \quad \Rightarrow \quad EG = 10 \times \sin 35^{\circ} \approx 10 \times 0{,}57357 \approx 5{,}74~\text{cm}. \]

\textbf{Vérification (Pythagore) :} $EF^2 + EG^2 \approx 8{,}19^2 + 5{,}74^2 \approx 67{,}08 + 32{,}95 \approx 100{,}03 \approx 10^2 = FG^2$. C'est cohérent.

\textbf{Conclusion :} $\widehat{EGF} = 55^{\circ}$, $EF \approx 8{,}19$~cm et $EG \approx 5{,}74$~cm.
\end{exoresolu}

\begin{methode}[Modéliser une situation concrète par un triangle rectangle]
\begin{enumerate}
\item \textbf{Lire} attentivement et repérer ce qui est vertical, horizontal, perpendiculaire.
\item \textbf{Schématiser} la situation par un triangle rectangle : nommer les points, coder l'angle droit, reporter les données (attention aux unités !).
\item \textbf{Traduire} la question en calcul de longueur ou d'angle.
\item \textbf{Appliquer} la trigonométrie (ou Pythagore) en rédigeant rigoureusement.
\item \textbf{Conclure} par une phrase répondant à la question, avec l'unité, et vérifier le réalisme du résultat.
\end{enumerate}
\end{methode}

\begin{exoresolu}[Descente d'un avion Camair-Co vers Douala]
Un avion de Camair-Co amorce sa descente vers l'aéroport de Douala. Il vole à une altitude constante de $900$~m et se trouve à $6$~km (distance horizontale) du point d'atterrissage $A$. On note $B$ la position de l'avion et $H$ le point au sol à la verticale de l'avion. Le triangle $ABH$ est rectangle en $H$.
\begin{enumerate}
\item Faire un schéma de la situation.
\item Calculer l'angle de descente $\widehat{BAH}$ (arrondir au degré).
\item Calculer la distance $AB$ que l'avion va parcourir pendant la descente (arrondir au mètre).
\end{enumerate}
\tcblower
\textbf{1. Schéma :} On convertit les unités : $6$~km $= 6000$~m.
$AH = 6000$~m (distance horizontale), $BH = 900$~m (altitude), angle droit en $H$.

\begin{popfigure}
\begin{tikzpicture}[scale=0.8, >=Stealth]
\coordinate (A) at (0,0);
\coordinate (H) at (6,0);
\coordinate (B) at (6,0.9);
\draw[popdark, thick] (A) -- node[below] {$6000$ m} (H) -- node[right] {$900$ m} (B) -- node[above left, pos=0.6] {$AB$ ?} cycle;
\draw[popblue, thick] (H) rectangle ++(-0.3,0.3); % angle droit
\draw[poporange, thick] (A) ++(0:0.8) arc (0:8.53:0.8) node[midway, right=2pt] {$\alpha$};
\node[below left] at (A) {$A$ (atterrissage)};
\node[below right] at (H) {$H$};
\node[above right] at (B) {$B$ (avion)};
\draw[dashed, popmuted] (H) -- ++(0, -0.5) node[below] {Verticale};
\draw[dashed, popmuted] (A) -- ++(6.5, 0) node[right] {Horizontale};
\end{tikzpicture}
\legende{Schéma de la descente (non à l'échelle : altitude exagérée pour la lisibilité).}
\end{popfigure}

\textbf{2. Angle de descente $\widehat{BAH}$ :} Par rapport à $\alpha = \widehat{BAH}$, on connaît l'opposé $BH = 900$ et l'adjacent $AH = 6000$. On utilise la tangente :
\[ \tan \alpha = \dfrac{BH}{AH} = \dfrac{900}{6000} = 0{,}15. \]
À la calculatrice (mode degrés) : $\alpha = \tan^{-1}(0{,}15) \approx 8{,}5307^{\circ}$.
Arrondi au degré : $\boxed{\alpha \approx 9^{\circ}}$.

\textbf{3. Distance $AB$ (hypoténuse) :} \emph{Méthode recommandée : Théorème de Pythagore (exact, sans arrondi d'angle).}
\[ AB = \sqrt{AH^2 + BH^2} = \sqrt{6000^2 + 900^2} = \sqrt{36\,000\,000 + 810\,000} = \sqrt{36\,810\,000} \approx 6067,1~\text{m}. \]
\emph{Méthode alternative (trigonométrie) :} $\sin \alpha = \frac{BH}{AB} \Rightarrow AB = \frac{900}{\sin 8{,}5307^{\circ}} \approx 6067$~m (en gardant la valeur exacte de l'angle dans la calculatrice).

\textbf{Conclusion :} L'avion descend sous un angle d'environ $\boxed{9^{\circ}}$ et parcourt environ $\boxed{6067~\text{m}}$ (soit $6{,}07$~km). Cet angle faible est réaliste pour un avion de ligne.
\end{exoresolu}

\begin{methode}[Rédiger une solution de trigonométrie au BEPC]
Une copie bien rédigée rapporte des points même si le résultat final est faux. Les réflexes :
\begin{enumerate}
\item \textbf{Annoncer les hypothèses} : « Le triangle $ABC$ est rectangle en $A$, donc on peut utiliser les rapports trigonométriques. »
\item \textbf{Citer la formule} avec les sommets du triangle : $\cos \widehat{ABC} = \dfrac{AB}{BC}$, pas une formule vague.
\item \textbf{Substituer} les valeurs numériques avant de calculer.
\item \textbf{Isoler} proprement l'inconnue (produit en croix écrit explicitement).
\item \textbf{Arrondir} uniquement à la fin, comme demandé, et écrire l'unité.
\item \textbf{Conclure} par une phrase complète : « Donc $BC \approx 5{,}3$~cm. »
\end{enumerate}
\end{methode}

\begin{exoresolu}[Rédaction complète type BEPC]
$MNP$ est un triangle rectangle en $M$ tel que $MN = 4{,}5$~cm et $NP = 7{,}5$~cm.
\begin{enumerate}
\item Calculer la mesure de l'angle $\widehat{MNP}$ (arrondir au degré).
\item En déduire la mesure de l'angle $\widehat{MPN}$.
\end{enumerate}
\tcblower
\textbf{1.} Le triangle $MNP$ est rectangle en $M$, donc on peut utiliser les rapports trigonométriques.
Par rapport à l'angle $\widehat{MNP}$ :
\begin{itemize}
\item le côté adjacent est $[MN]$ avec $MN = 4{,}5$~cm ;
\item l'hypoténuse est $[NP]$ avec $NP = 7{,}5$~cm.
\end{itemize}
On utilise donc le cosinus :
\[ \cos \widehat{MNP} = \dfrac{MN}{NP} = \dfrac{4{,}5}{7{,}5} = 0{,}6. \]
À la calculatrice (mode degrés) :
\[ \widehat{MNP} = \cos^{-1}(0{,}6) \approx 53{,}13^{\circ}. \]
Donc $\boxed{\widehat{MNP} \approx 53^{\circ}}$.

\textbf{2.} Dans le triangle $MNP$ rectangle en $M$, les deux angles aigus sont complémentaires :
\[ \widehat{MPN} = 90^{\circ} - \widehat{MNP} \approx 90^{\circ} - 53^{\circ} = 37^{\circ}. \]
Donc $\boxed{\widehat{MPN} \approx 37^{\circ}}$.
\end{exoresolu}

\begin{attention}[Les 5 erreurs qui coûtent des points au BEPC]
\begin{enumerate}
\item \textbf{Oublier de justifier que le triangle est rectangle.} La trigonométrie (sinus, cosinus, tangente) ne s'applique \emph{que} dans un triangle rectangle : écris toujours « le triangle ... est rectangle en ..., donc ... ».
\item \textbf{Calculatrice en mauvais mode.} Si la calculatrice est en radians, $\cos 60$ ne donne pas $0{,}5$ ! Vérifie le mode \emph{degrés} avant tout calcul (teste avec $\cos 60^{\circ} = 0{,}5$).
\item \textbf{Confondre opposé et adjacent.} Le côté opposé et le côté adjacent dépendent de \emph{l'angle choisi} : repère-les toujours par rapport à l'angle utilisé. Seule l'hypoténuse ne change jamais (côté opposé à l'angle droit).
\item \textbf{Erreur de produit en croix.} De $\sin \widehat{A} = \dfrac{BC}{AB}$, on tire $BC = AB \times \sin \widehat{A}$ et non $BC = \dfrac{\sin \widehat{A}}{AB}$. Écris l'étape d'isolement de l'inconnue.
\item \textbf{Arrondir trop tôt ou oublier la conclusion.} Garde les valeurs exactes de la calculatrice jusqu'au bout, arrondis seulement au résultat final comme demandé, et termine par une phrase avec l'unité. Un sinus ou un cosinus supérieur à $1$ doit te mettre la puce à l'oreille : c'est impossible, cherche l'erreur !
\end{enumerate}
\end{attention}

\begin{savaistu}[La trigonométrie au Cameroun : du cadastre au GPS]
Le mot « trigonométrie » vient du grec \emph{trigonon} (triangle) et \emph{metron} (mesure). Au Cameroun, les géomètres-topographes du MINDAF (Ministère des Domaines, du Cadastre et des Affaires Foncières) l'utilisent quotidiennement pour le bornage des terrains : en mesurant un angle et une distance (méthode du rayonnement), ils calculent les coordonnées des bornes. Les antennes-relais MTN/Orange sont implantées grâce à des calculs d'azimuts et de portées (triangles rectangles verticaux). Et ton téléphone portable ? Le GPS utilise la \emph{trilatération} (intersection de sphères), généralisation 3D de la trigonométrie, pour te situer à quelques mètres près !
\end{savaistu}

\newpage

\coursec{Exercices}

\courssub{Je vérifie mes connaissances}

\begin{exercice}[QCM : les bons rapports]
Pour chaque question, une seule réponse est exacte. On considère un triangle $ABC$ rectangle en $A$.
\begin{enumerate}
\item $\cos(\widehat{ABC})$ est égal à :
\begin{enumerate}
\item[a)] $\dfrac{AB}{BC}$ \qquad b)] $\dfrac{AC}{BC}$ \qquad c)] $\dfrac{AC}{AB}$
\end{enumerate}
\item $\sin(\widehat{ACB})$ est égal à :
\begin{enumerate}
\item[a)] $\dfrac{AC}{BC}$ \qquad b)] $\dfrac{AB}{AC}$ \qquad c)] $\dfrac{AB}{BC}$
\end{enumerate}
\item $\tan(\widehat{ABC})$ est égal à :
\begin{enumerate}
\item[a)] $\dfrac{AB}{AC}$ \qquad b)] $\dfrac{AC}{AB}$ \qquad c)] $\dfrac{AB}{BC}$
\end{enumerate}
\item Si $\widehat{ABC} = 60°$, alors $\cos(\widehat{ABC})$ vaut :
\begin{enumerate}
\item[a)] $\dfrac{\sqrt{3}}{2}$ \qquad b)] $\dfrac{1}{2}$ \qquad c)] $\dfrac{\sqrt{2}}{2}$
\end{enumerate}
\end{enumerate}
\end{exercice}

\begin{exercice}[Vrai ou faux ? Justifier]
On considère un triangle $MNP$ rectangle en $M$. Répondre par vrai ou faux en justifiant chaque réponse.
\begin{enumerate}
\item Le côté $[NP]$ est l'hypoténuse du triangle.
\item Par rapport à l'angle $\widehat{MNP}$, le côté $[MP]$ est le côté adjacent.
\item $\sin(\widehat{MNP}) = \cos(\widehat{MPN})$.
\item On peut avoir $\cos(\widehat{MNP}) = 1,4$.
\end{enumerate}
\end{exercice}

\begin{exercice}[Vocabulaire et définitions]
Soit $RST$ un triangle rectangle en $S$.
\begin{enumerate}
\item Citer l'hypoténuse de ce triangle.
\item Par rapport à l'angle $\widehat{SRT}$, citer le côté adjacent et le côté opposé.
\item Compléter les égalités : $\cos(\widehat{SRT}) = \dfrac{\ldots}{\ldots}$ ; $\sin(\widehat{SRT}) = \dfrac{\ldots}{\ldots}$ ; $\tan(\widehat{SRT}) = \dfrac{\ldots}{\ldots}$.
\item Rappeler la relation fondamentale reliant $\cos^2(\alpha)$ et $\sin^2(\alpha)$ pour un angle aigu $\alpha$.
\end{enumerate}
\end{exercice}

\begin{exercice}[Calculs directs]
\begin{enumerate}
\item À l'aide de la calculatrice, donner la valeur arrondie au millième de : $\cos(35°)$ ; $\sin(58°)$ ; $\tan(20°)$.
\item Déterminer, à $1°$ près, la mesure de l'angle aigu $\alpha$ tel que : $\cos(\alpha) = 0,5$ ; $\sin(\alpha) = 0,8$ ; $\tan(\alpha) = 1$.
\item Donner les valeurs exactes de : $\cos(30°)$ ; $\sin(45°)$ ; $\tan(60°)$.
\end{enumerate}
\end{exercice}

\courssub{Je m'entraîne}

\begin{exercice}[Calculer une longueur avec le cosinus]
Soit $ABC$ un triangle rectangle en $A$ tel que $BC = 10$ cm et $\widehat{ABC} = 40°$.
\begin{enumerate}
\item Faire une figure à main levée en codant les données.
\item Calculer $AB$, arrondi au millimètre.
\item Calculer $AC$, arrondi au millimètre.
\end{enumerate}
\end{exercice}

\begin{exercice}[Calculer une longueur avec la tangente]
Soit $EFG$ un triangle rectangle en $E$ tel que $EF = 7$ cm et $\widehat{EFG} = 55°$.
\begin{enumerate}
\item Quel rapport trigonométrique permet de calculer $EG$ directement ? Justifier.
\item Calculer $EG$, arrondi au millimètre.
\item Calculer $FG$, arrondi au millimètre.
\end{enumerate}
\end{exercice}

\begin{exercice}[Calculer un angle aigu]
Soit $IJK$ un triangle rectangle en $I$ tel que $IJ = 4$ cm et $IK = 3$ cm.
\begin{enumerate}
\item Calculer $JK$ en justifiant.
\item Calculer la mesure de l'angle $\widehat{IJK}$, arrondie au degré.
\item En déduire la mesure de l'angle $\widehat{IKJ}$ sans calculatrice.
\end{enumerate}
\end{exercice}

\begin{exercice}[Valeurs exactes et demi-équilatéral]
Soit $ABC$ un triangle rectangle en $A$ tel que $BC = 8$ cm et $\widehat{ABC} = 60°$.
\begin{enumerate}
\item Calculer les valeurs exactes de $AB$ et $AC$.
\item Vérifier que $AB^2 + AC^2 = BC^2$.
\item Le triangle $ABC$ est-il un demi-triangle équilatéral ? Justifier.
\end{enumerate}
\end{exercice}

\courssub{Je me prépare au BEPC}

\begin{exercice}[L'échelle contre le mur — type BEPC]
Une échelle de $5$ m est appuyée contre un mur vertical. Son pied est situé à $2$ m du pied du mur. On modélise la situation par un triangle $ABC$ rectangle en $B$, où $B$ est le pied du mur, $A$ le pied de l'échelle et $C$ le point de contact de l'échelle avec le mur.
\begin{enumerate}
\item Faire une figure en vraie grandeur en choisissant l'échelle $1/100$.
\item Calculer l'angle $\widehat{BAC}$ que l'échelle fait avec le sol, arrondi au degré.
\item Calculer la hauteur $BC$ atteinte par l'échelle sur le mur, arrondie au centimètre.
\item Pour des raisons de sécurité, l'angle entre l'échelle et le sol doit être compris entre $65°$ et $75°$. L'échelle est-elle correctement installée ? Justifier.
\item On recule le pied de l'échelle jusqu'à $3$ m du mur. Calculer la nouvelle mesure de l'angle $\widehat{BAC}$, arrondie au degré.
\end{enumerate}
\end{exercice}

\begin{exercice}[Le pylône haubané — type BEPC]
Dans un quartier de Yaoundé, un pylône électrique vertical est maintenu par un câble (hauban) fixé au sol à $8$ m du pied du pylône. Le câble fait un angle de $60°$ avec le sol horizontal. On note $P$ le pied du pylône, $S$ son sommet et $A$ le point d'ancrage du câble au sol.
\begin{enumerate}
\item Faire une figure codée de la situation.
\item Calculer la valeur exacte de la hauteur $PS$ du pylône, puis en donner une valeur arrondie au décimètre.
\item Calculer la valeur exacte de la longueur $AS$ du câble.
\item L'entreprise d'électrification dispose de câbles vendus par rouleaux de $20$ m. Un rouleau suffit-il pour haubaner ce pylône ? Justifier.
\item Un second hauban est fixé au même sommet $S$, de l'autre côté du pylône, et fait un angle de $45°$ avec le sol. Calculer la distance entre le pied $P$ du pylône et le point d'ancrage de ce second câble (valeur exacte).
\end{enumerate}
\end{exercice}

\begin{exercice}[La largeur du fleuve et la route du Mont Fébé — problème de synthèse]
\textbf{Partie A : mesurer sans traverser.}\\
Depuis un point $A$ au bord d'un fleuve, un élève aperçoit un arbre $B$ sur l'autre rive, exactement en face de lui. Il plante un piquet $C$ sur la même rive que $A$, de sorte que $AC = 20$ m et $\widehat{CAB} = 90°$. Il mesure ensuite $\widehat{ACB} = 35°$.
\begin{enumerate}
\item Faire une figure à main levée de la situation.
\item Expliquer pourquoi le triangle $ABC$ est rectangle et préciser en quel point.
\item Calculer la largeur $AB$ du fleuve, arrondie au décimètre.
\item Calculer la distance $BC$ qui sépare le piquet de l'arbre, arrondie au décimètre.
\end{enumerate}
\textbf{Partie B : la pente du Mont Fébé.}\\
Une route de Yaoundé qui monte vers le Mont Fébé a une pente moyenne de $12\,\%$, c'est-à-dire que pour $100$ m parcourus à l'horizontale, on s'élève de $12$ m. On modélise la route par l'hypoténuse d'un triangle rectangle.
\begin{enumerate}
\item[5.] Calculer l'angle d'inclinaison moyen de la route avec l'horizontale, arrondi au degré.
\item[6.] Un cycliste parcourt $500$ m sur la route (distance réelle sur la chaussée). Calculer le dénivelé positif (la hauteur gagnée), arrondi au mètre.
\item[7.] Calculer la distance horizontale correspondant à ces $500$ m de route, arrondie au mètre.
\end{enumerate}
\end{exercice}

\newpage

\coursec{Corrigés des exercices}

\courssub{Repérage et vocabulaire dans le triangle rectangle}

\begin{corrige}[Exercice 1 — QCM : les bons rapports]
Dans le triangle $ABC$ rectangle en $A$, l'hypoténuse est $[BC]$.
\begin{enumerate}
\item Pour l'angle $\widehat{ABC}$ : le côté adjacent est $[AB]$ et l'hypoténuse est $[BC]$, donc $\cos(\widehat{ABC}) = \dfrac{AB}{BC}$. \textbf{Réponse a).}
\item Pour l'angle $\widehat{ACB}$ : le côté opposé est $[AB]$ et l'hypoténuse est $[BC]$, donc $\sin(\widehat{ACB}) = \dfrac{AB}{BC}$. \textbf{Réponse c).}
\item Pour l'angle $\widehat{ABC}$ : le côté opposé est $[AC]$ et le côté adjacent est $[AB]$, donc $\tan(\widehat{ABC}) = \dfrac{AC}{AB}$. \textbf{Réponse b).}
\item $\cos(60°) = \dfrac{1}{2}$ (valeur exacte à connaître). \textbf{Réponse b).}
\end{enumerate}
\end{corrige}

\begin{corrige}[Exercice 2 — Vrai ou faux ? Justifier]
Le triangle $MNP$ est rectangle en $M$ : l'hypoténuse est le côté opposé à l'angle droit, c'est $[NP]$.
\begin{enumerate}
\item \textbf{Vrai.} L'hypoténuse est le côté opposé à l'angle droit $\widehat{NMP}$ ; c'est donc $[NP]$.
\item \textbf{Faux.} Par rapport à l'angle $\widehat{MNP}$ (sommet $N$), le côté adjacent est $[MN]$ (le côté de l'angle autre que l'hypoténuse). Le côté $[MP]$ est le côté \emph{opposé} à cet angle.
\item \textbf{Vrai.} On a $\sin(\widehat{MNP}) = \dfrac{MP}{NP}$ (opposé sur hypoténuse) et $\cos(\widehat{MPN}) = \dfrac{MP}{NP}$ (adjacent sur hypoténuse). Les deux rapports sont égaux. (C'est le cas général : les angles aigus d'un triangle rectangle sont complémentaires, et le sinus de l'un égale le cosinus de l'autre.)
\item \textbf{Faux.} Pour tout angle aigu, le cosinus est le quotient d'un côté de l'angle droit par l'hypoténuse, qui est le plus grand côté : ce quotient est donc toujours compris entre $0$ et $1$. On ne peut pas avoir $\cos(\widehat{MNP}) = 1,4 > 1$.
\end{enumerate}
\end{corrige}

\begin{corrige}[Exercice 3 — Vocabulaire et définitions]
\begin{enumerate}
\item L'hypoténuse est le côté opposé à l'angle droit $\widehat{RST}$ : c'est $[RT]$.
\item Par rapport à l'angle $\widehat{SRT}$ (sommet $R$) : le côté \cle{adjacent} est $[RS]$ et le côté \cle{opposé} est $[ST]$.
\item $\cos(\widehat{SRT}) = \dfrac{RS}{RT}$ ; \quad $\sin(\widehat{SRT}) = \dfrac{ST}{RT}$ ; \quad $\tan(\widehat{SRT}) = \dfrac{ST}{RS}$.
\item Pour tout angle aigu $\alpha$ : $\cos^2(\alpha) + \sin^2(\alpha) = 1$.
\end{enumerate}
\end{corrige}

\courssub{Définition des rapports trigonométriques : Cosinus, Sinus, Tangente}

\begin{corrige}[Exercice 4 — Calculs directs]
\begin{enumerate}
\item Avec la calculatrice (en mode degrés) :
\[ \cos(35°) \approx 0,819 \ ; \qquad \sin(58°) \approx 0,848 \ ; \qquad \tan(20°) \approx 0,364. \]
\item On utilise les touches $\cos^{-1}$, $\sin^{-1}$, $\tan^{-1}$ :
\begin{itemize}
\item $\cos(\alpha) = 0,5$ donne $\alpha = 60°$ (valeur exacte connue) ;
\item $\sin(\alpha) = 0,8$ donne $\alpha \approx 53°$ (car $\sin^{-1}(0,8) \approx 53,13°$) ;
\item $\tan(\alpha) = 1$ donne $\alpha = 45°$.
\end{itemize}
\item Valeurs exactes : $\cos(30°) = \dfrac{\sqrt{3}}{2}$ ; \quad $\sin(45°) = \dfrac{\sqrt{2}}{2}$ ; \quad $\tan(60°) = \sqrt{3}$.
\end{enumerate}
\end{corrige}

\courssub{Angles remarquables (30°, 45°, 60°) et valeurs exactes}

\begin{corrige}[Exercice 8 — Valeurs exactes et demi-équilatéral]
\begin{enumerate}
\item Dans le triangle $ABC$ rectangle en $A$, par rapport à $\widehat{ABC} = 60°$ :
\[ AB = BC \times \cos(60°) = 8 \times \dfrac{1}{2} = 4 \text{ cm} \ ; \qquad AC = BC \times \sin(60°) = 8 \times \dfrac{\sqrt{3}}{2} = 4\sqrt{3} \text{ cm}. \]
Donc $\boxed{AB = 4 \text{ cm}}$ et $\boxed{AC = 4\sqrt{3} \text{ cm}}$ (valeurs exactes).
\item Vérifions :
\[ AB^2 + AC^2 = 4^2 + (4\sqrt{3})^2 = 16 + 16 \times 3 = 16 + 48 = 64 = 8^2 = BC^2. \]
L'égalité de Pythagore est bien vérifiée.
\item \textbf{Oui.} Considérons le triangle équilatéral $BCD$ de côté $8$ cm, où $A$ est le milieu de $[BD]$. La hauteur issue de $C$ coupe $[BD]$ en son milieu $A$, avec $BA = 4$ cm, et les angles valent $\widehat{ABC} = 60°$ et $\widehat{ACB} = 30°$. Le triangle $ABC$ obtenu a exactement les dimensions calculées ($AB = 4$, $AC = 4\sqrt{3}$, $BC = 8$) : c'est bien la moitié d'un triangle équilatéral de côté $8$ cm, partagé par une hauteur. On retrouve au passage que la hauteur d'un triangle équilatéral de côté $a$ vaut $\dfrac{a\sqrt{3}}{2}$.
\end{enumerate}
\end{corrige}

\courssub{Calculer une longueur dans un triangle rectangle}

\begin{corrige}[Exercice 5 — Calculer une longueur avec le cosinus]
\begin{enumerate}
\item Figure à main levée : triangle $ABC$ avec l'angle droit en $A$, l'hypoténuse $[BC]$ de longueur $10$ cm, et l'angle de $40°$ codé en $B$.
\item Par rapport à l'angle $\widehat{ABC} = 40°$, le côté $[AB]$ est adjacent et $[BC]$ est l'hypoténuse. On utilise le cosinus :
\[ \cos(\widehat{ABC}) = \dfrac{AB}{BC} \quad \text{donc} \quad AB = BC \times \cos(40°) = 10 \times \cos(40°). \]
Comme $\cos(40°) \approx 0,766$, on obtient $AB \approx 7,66$ cm, soit $\boxed{AB \approx 7,7 \text{ cm}}$ (arrondi au millimètre : $7,7$ cm $= 77$ mm ; plus précisément $AB \approx 76,6$ mm).
\item Par rapport à $\widehat{ABC}$, le côté $[AC]$ est opposé. On utilise le sinus :
\[ AC = BC \times \sin(40°) = 10 \times \sin(40°) \approx 10 \times 0,643 = 6,43 \text{ cm}. \]
Donc $\boxed{AC \approx 6,4 \text{ cm}}$ (soit environ $64$ mm).
\end{enumerate}
\textbf{Vérification :} $AB^2 + AC^2 \approx 7,66^2 + 6,43^2 \approx 58,7 + 41,3 = 100 = BC^2$. Le théorème de Pythagore est bien vérifié.
\end{corrige}

\begin{corrige}[Exercice 6 — Calculer une longueur avec la tangente]
\begin{enumerate}
\item Par rapport à l'angle $\widehat{EFG} = 55°$ (sommet $F$) : on connaît le côté adjacent $EF = 7$ cm et on cherche le côté opposé $EG$. Le rapport qui relie côté opposé et côté adjacent est la \cle{tangente} :
\[ \tan(\widehat{EFG}) = \dfrac{EG}{EF}. \]
\item On en déduit :
\[ EG = EF \times \tan(55°) = 7 \times \tan(55°) \approx 7 \times 1,428 = 9,997 \text{ cm}. \]
Donc $\boxed{EG \approx 10,0 \text{ cm}}$ (arrondi au millimètre).
\item Pour $FG$ (l'hypoténuse), on utilise le cosinus de $\widehat{EFG}$ :
\[ \cos(55°) = \dfrac{EF}{FG} \quad \text{donc} \quad FG = \dfrac{EF}{\cos(55°)} = \dfrac{7}{\cos(55°)} \approx \dfrac{7}{0,574} \approx 12,2 \text{ cm}. \]
Donc $\boxed{FG \approx 12,2 \text{ cm}}$.
\end{enumerate}
\textbf{Vérification :} $EF^2 + EG^2 \approx 49 + 99,9 = 148,9$ et $FG^2 \approx 12,2^2 = 148,8$ : cohérent (aux arrondis près).
\end{corrige}

\courssub{Calculer un angle aigu dans un triangle rectangle}

\begin{corrige}[Exercice 7 — Calculer un angle aigu]
\begin{enumerate}
\item Le triangle $IJK$ est rectangle en $I$, donc d'après le théorème de Pythagore :
\[ JK^2 = IJ^2 + IK^2 = 4^2 + 3^2 = 16 + 9 = 25 \quad \text{donc} \quad JK = \sqrt{25} = 5 \text{ cm}. \]
\item Par rapport à l'angle $\widehat{IJK}$ (sommet $J$) : le côté opposé est $[IK] = 3$ cm et l'hypoténuse est $[JK] = 5$ cm. On utilise le sinus :
\[ \sin(\widehat{IJK}) = \dfrac{IK}{JK} = \dfrac{3}{5} = 0,6. \]
À la calculatrice : $\widehat{IJK} = \sin^{-1}(0,6) \approx 36,87°$, donc $\boxed{\widehat{IJK} \approx 37°}$.
\item Les deux angles aigus d'un triangle rectangle sont complémentaires :
\[ \widehat{IKJ} = 90° - \widehat{IJK} \approx 90° - 37° = 53°. \]
Donc $\boxed{\widehat{IKJ} \approx 53°}$, sans calculatrice.
\end{enumerate}
\end{corrige}

\courssub{Résolution de triangle rectangle \& Problèmes concrets (Intégration)}

\begin{corrige}[Exercice 9 — L'échelle contre le mur — type BEPC]
\textbf{Barème indicatif (sur 10) :} figure 1 pt ; question 2 : 2 pts ; question 3 : 2 pts ; question 4 : 2 pts ; question 5 : 2 pts ; rédaction et unités 1 pt.
\begin{enumerate}
\item À l'échelle $1/100$, $5$ m sont représentés par $5$ cm et $2$ m par $2$ cm. On trace un segment $[AB]$ horizontal de $2$ cm (le sol), la perpendiculaire en $B$ (le mur), puis on reporte $AC = 5$ cm au compas depuis $A$ pour placer $C$ sur le mur.
\begin{popfigure}
\begin{tikzpicture}[scale=0.5, >=Stealth]
\draw[popblue, thick] (0,0) coordinate (A) -- (2,0) coordinate (B) node[midway, below] {$2$ m};
\draw[popblue, thick] (2,0) -- (2,4.583) coordinate (C) node[midway, right] {$BC$};
\draw[poporange, thick, dashed] (0,0) -- (2,4.583) node[midway, above left, sloped] {$5$ m};
\draw (2,0) rectangle +(0.3,0.3);
\pic[draw, popmuted, angle radius=8mm, "$\alpha$", angle eccentricity=1.3] {angle = C--A--B};
\node[below left] at (A) {$A$ (pied)};
\node[below right] at (B) {$B$ (base mur)};
\node[above right] at (C) {$C$ (haut)};
\end{tikzpicture}
\legende{L'échelle contre le mur : triangle $ABC$ rectangle en $B$.}
\end{popfigure}
\item Dans le triangle $ABC$ rectangle en $B$, par rapport à l'angle $\widehat{BAC}$ : le côté adjacent est $[AB] = 2$ m et l'hypoténuse est $[AC] = 5$ m. On utilise le cosinus :
\[ \cos(\widehat{BAC}) = \dfrac{AB}{AC} = \dfrac{2}{5} = 0,4 \quad \text{donc} \quad \widehat{BAC} = \cos^{-1}(0,4) \approx 66,4°. \]
Donc $\boxed{\widehat{BAC} \approx 66°}$.
\item On utilise le sinus (ou Pythagore) :
\[ BC = AC \times \sin(\widehat{BAC}) \approx 5 \times \sin(66,4°) \approx 5 \times 0,9165 \approx 4,58 \text{ m}. \]
(Par Pythagore : $BC = \sqrt{5^2 - 2^2} = \sqrt{21} \approx 4,58$ m : même résultat.) Donc $\boxed{BC \approx 4,58 \text{ m}}$.
\item L'angle mesuré vaut environ $66°$, et $65° \leq 66° \leq 75°$ : l'angle est bien compris entre $65°$ et $75°$. \textbf{L'échelle est correctement installée.}
\item Avec $AB = 3$ m : $\cos(\widehat{BAC}) = \dfrac{3}{5} = 0,6$, donc $\widehat{BAC} = \cos^{-1}(0,6) \approx 53,1°$, soit $\boxed{\widehat{BAC} \approx 53°}$. L'échelle est alors moins inclinée (angle plus faible) : le pied de l'échelle risque de glisser vers l'extérieur.
\end{enumerate}
\textbf{Points de vigilance :} citer le triangle rectangle et le rapport utilisé avant chaque calcul ; vérifier que la calculatrice est en mode degrés ; ne pas confondre le côté adjacent ($[AB]$, le long du sol) avec le côté opposé ($[BC]$, le mur) ; conclure chaque question par une phrase avec l'unité.
\end{corrige}

\begin{corrige}[Exercice 10 — Le pylône haubané — type BEPC]
\textbf{Barème indicatif (sur 10) :} figure 1 pt ; question 2 : 2 pts ; question 3 : 2 pts ; question 4 : 1,5 pt ; question 5 : 2,5 pts ; rédaction 1 pt.
\begin{enumerate}
\item Figure : segment $[PA]$ horizontal de longueur $8$ (le sol), segment $[PS]$ vertical (le pylône), angle droit en $P$, angle de $60°$ codé en $A$ entre le sol et le câble $[AS]$.
\begin{popfigure}
\begin{tikzpicture}[scale=0.35, >=Stealth]
% Premier hauban (60°)
\draw[popblue, thick] (0,0) coordinate (P) -- (8,0) coordinate (A) node[midway, below] {$8$ m};
\draw[popblue, thick] (0,0) -- (0,13.856) coordinate (S) node[midway, left] {$PS = 8\sqrt{3}$};
\draw[poporange, thick, dashed] (8,0) -- (0,13.856) node[midway, above right, sloped] {$AS = 16$ m};
\draw (0,0) rectangle +(0.5,0.5);
\pic[draw, popmuted, angle radius=12mm, "$60°$", angle eccentricity=1.4] {angle = P--A--S};
\node[below left] at (P) {$P$};
\node[below right] at (A) {$A$};
\node[above left] at (S) {$S$ (sommet)};
% Second hauban (45°)
\draw[popgreen, thick, dashed] (0,13.856) -- (-13.856,0) coordinate (A2) node[midway, above left, sloped] {$A'S$};
\draw[popblue, thick] (0,0) -- (-13.856,0) node[midway, below] {$PA' = 8\sqrt{3}$};
\pic[draw, popmuted, angle radius=12mm, "$45°$", angle eccentricity=1.4] {angle = P--A2--S};
\node[below left] at (A2) {$A'$};
\end{tikzpicture}
\legende{Pylône haubané : triangles rectangles $PAS$ ($60°$) et $PA'S$ ($45°$).}
\end{popfigure}
\item Le triangle $PAS$ est rectangle en $P$. Par rapport à l'angle $\widehat{PAS} = 60°$ : on connaît le côté adjacent $PA = 8$ m et on cherche le côté opposé $PS$. On utilise la tangente :
\[ PS = PA \times \tan(60°) = 8\sqrt{3} \text{ m}. \]
Valeur exacte : $\boxed{PS = 8\sqrt{3} \text{ m}}$, soit $PS \approx 13,9$ m au décimètre (car $8\sqrt{3} \approx 13,856$).
\item Pour le câble $AS$ (hypoténuse), on utilise le cosinus :
\[ \cos(60°) = \dfrac{PA}{AS} \quad \text{donc} \quad AS = \dfrac{PA}{\cos(60°)} = \dfrac{8}{\tfrac{1}{2}} = 16 \text{ m}. \]
Valeur exacte : $\boxed{AS = 16 \text{ m}}$.
\item Le câble mesure $16$ m et un rouleau contient $20$ m. Comme $16 < 20$, \textbf{un rouleau suffit} (il reste même $4$ m de câble).
\item Soit $A'$ le second point d'ancrage, de l'autre côté du pylône. Le triangle $PA'S$ est rectangle en $P$, avec $\widehat{PA'S} = 45°$ et le côté opposé $PS = 8\sqrt{3}$ m. On cherche le côté adjacent $PA'$ :
\[ \tan(45°) = \dfrac{PS}{PA'} \quad \text{donc} \quad PA' = \dfrac{PS}{\tan(45°)} = \dfrac{8\sqrt{3}}{1} = 8\sqrt{3} \text{ m}. \]
La distance cherchée est $\boxed{PA' = 8\sqrt{3} \text{ m} \approx 13,9 \text{ m}}$. (Remarque : à $45°$, le triangle est rectangle isocèle, donc $PA' = PS$.)
\end{enumerate}
\textbf{Points de vigilance :} bien distinguer valeur exacte ($8\sqrt{3}$ m) et valeur arrondie ($13,9$ m) — les deux sont demandées ; pour la question 5, l'angle de $45°$ est au point d'ancrage $A'$, pas au sommet ; justifier l'utilisation de la tangente en nommant les côtés (opposé/adjacent).
\end{corrige}

\begin{corrige}[Exercice 11 — La largeur du fleuve et la route du Mont Fébé — problème de synthèse]
\textbf{Barème indicatif (sur 10) :} Partie A : 5 pts (figure 1 pt, question 2 : 1 pt, question 3 : 1,5 pt, question 4 : 1,5 pt) ; Partie B : 5 pts (question 5 : 1,5 pt, question 6 : 2 pts, question 7 : 1,5 pt).

\textbf{Partie A.}
\begin{enumerate}
\item Figure à main levée : $A$ et $C$ sur la même rive avec $AC = 20$ m, $B$ sur l'autre rive en face de $A$, angle droit codé en $A$, angle de $35°$ codé en $C$.
\begin{popfigure}
\begin{tikzpicture}[scale=0.3, >=Stealth]
\draw[popblue, thick] (0,0) coordinate (A) -- (20,0) coordinate (C) node[midway, below] {$20$ m};
\draw[popblue, thick] (0,0) -- (0,14) coordinate (B) node[midway, left] {$AB$};
\draw[poporange, thick, dashed] (20,0) -- (0,14) node[midway, above right, sloped] {$BC$};
\draw (0,0) rectangle +(0.8,0.8);
\pic[draw, popmuted, angle radius=15mm, "$35°$", angle eccentricity=1.4] {angle = A--C--B};
\node[below left] at (A) {$A$};
\node[below right] at (C) {$C$};
\node[above left] at (B) {$B$ (autre rive)};
\end{tikzpicture}
\legende{Mesure de la largeur du fleuve : triangle $ABC$ rectangle en $A$.}
\end{popfigure}
\item On a planté le piquet de sorte que $\widehat{CAB} = 90°$ : le triangle $ABC$ est donc \textbf{rectangle en $A$}, d'hypoténuse $[BC]$.
\item Par rapport à l'angle $\widehat{ACB} = 35°$ : le côté adjacent est $[AC] = 20$ m et le côté opposé est $[AB]$ (la largeur du fleuve). On utilise la tangente :
\[ AB = AC \times \tan(35°) = 20 \times \tan(35°) \approx 20 \times 0,7002 \approx 14,0 \text{ m}. \]
La largeur du fleuve est $\boxed{AB \approx 14,0 \text{ m}}$.
\item Pour $BC$ (hypoténuse), on utilise le cosinus :
\[ \cos(35°) = \dfrac{AC}{BC} \quad \text{donc} \quad BC = \dfrac{AC}{\cos(35°)} = \dfrac{20}{\cos(35°)} \approx \dfrac{20}{0,8192} \approx 24,4 \text{ m}. \]
Donc $\boxed{BC \approx 24,4 \text{ m}}$.
\end{enumerate}

\textbf{Partie B.}
\begin{enumerate}
\item[5.] Une pente de $12\,\%$ signifie que, dans le triangle rectangle modélisant la route, le côté opposé à l'angle d'inclinaison $\alpha$ vaut $12$ m quand le côté adjacent vaut $100$ m. D'où :
\[ \tan(\alpha) = \dfrac{12}{100} = 0,12 \quad \text{donc} \quad \alpha = \tan^{-1}(0,12) \approx 6,8°. \]
L'angle d'inclinaison moyen est $\boxed{\alpha \approx 7°}$.
\item[6.] Le cycliste parcourt $500$ m sur l'hypoténuse. Le dénivelé $h$ est le côté opposé à l'angle $\alpha$ :
\[ h = 500 \times \sin(\alpha) \approx 500 \times \sin(6,8°) \approx 500 \times 0,1185 \approx 59 \text{ m}. \]
Le dénivelé positif est $\boxed{h \approx 59 \text{ m}}$.
\item[7.] La distance horizontale $d$ est le côté adjacent :
\[ d = 500 \times \cos(\alpha) \approx 500 \times \cos(6,8°) \approx 500 \times 0,9930 \approx 496 \text{ m}. \]
La distance horizontale est $\boxed{d \approx 496 \text{ m}}$.
\end{enumerate}
\textbf{Vérification (Pythagore) :} $496^2 + 59^2 = 246\,016 + 3\,481 = 249\,497 \approx 500^2 = 250\,000$ : cohérent aux arrondis près.

\textbf{Points de vigilance :} dans la Partie A, l'angle de $35°$ est au piquet $C$, pas en $A$ — c'est ce qui permet de calculer $AB$ avec la tangente ; dans la Partie B, la pente en pourcentage se traduit par une \emph{tangente} (rapport dénivelé/horizontal), et non par un sinus ; les $500$ m du cycliste sont parcourus sur la chaussée, donc sur l'\emph{hypoténuse} ; arrondir seulement à la fin des calculs.
\end{corrige}

\newpage

\coursec{Fiche bilan}

\begin{popfigure}
\begin{center}\tcbox[colback=popink,colframe=popink,coltext=white,arc=9pt,boxsep=4pt,left=6pt,right=6pt]{\bfseries\small Trigonométrie dans le triangle rectangle}\end{center}
\vspace{-2pt}
\begin{tcbraster}[raster columns=3,raster equal height=rows,raster column skip=6pt,raster row skip=6pt,enhanced,blankest]
\begin{tcolorbox}[enhanced,colback=white,colframe=popblue,boxrule=0.9pt,arc=3pt,title={1. Vocabulaire},fonttitle=\bfseries\scriptsize,coltitle=white,colbacktitle=popblue,left=4pt,right=4pt,top=2pt,bottom=2pt,boxsep=1pt,toptitle=1.5pt,bottomtitle=1.5pt,halign=left]
\begin{itemize}[nosep,leftmargin=*,itemsep=1pt,font=\scriptsize]
\item Hypoténuse
\item Côté adjacent
\item Côté opposé
\end{itemize}
\end{tcolorbox}
\begin{tcolorbox}[enhanced,colback=white,colframe=poporange,boxrule=0.9pt,arc=3pt,title={2. Rapports trigonométriques},fonttitle=\bfseries\scriptsize,coltitle=white,colbacktitle=poporange,left=4pt,right=4pt,top=2pt,bottom=2pt,boxsep=1pt,toptitle=1.5pt,bottomtitle=1.5pt,halign=left]
\begin{itemize}[nosep,leftmargin=*,itemsep=1pt,font=\scriptsize]
\item $\cos = \frac{\text{adj}}{\text{hyp}}$
\item $\sin = \frac{\text{opp}}{\text{hyp}}$
\item $\tan = \frac{\text{opp}}{\text{adj}}$
\end{itemize}
\end{tcolorbox}
\begin{tcolorbox}[enhanced,colback=white,colframe=popgreen,boxrule=0.9pt,arc=3pt,title={3. Angles remarquables},fonttitle=\bfseries\scriptsize,coltitle=white,colbacktitle=popgreen,left=4pt,right=4pt,top=2pt,bottom=2pt,boxsep=1pt,toptitle=1.5pt,bottomtitle=1.5pt,halign=left]
\begin{itemize}[nosep,leftmargin=*,itemsep=1pt,font=\scriptsize]
\item $30^\circ, 45^\circ, 60^\circ$
\item Valeurs exactes
\item Tableau de valeurs
\end{itemize}
\end{tcolorbox}
\begin{tcolorbox}[enhanced,colback=white,colframe=poppurple,boxrule=0.9pt,arc=3pt,title={4. Calculer une longueur},fonttitle=\bfseries\scriptsize,coltitle=white,colbacktitle=poppurple,left=4pt,right=4pt,top=2pt,bottom=2pt,boxsep=1pt,toptitle=1.5pt,bottomtitle=1.5pt,halign=left]
\begin{itemize}[nosep,leftmargin=*,itemsep=1pt,font=\scriptsize]
\item Formule directe
\item Produit en croix
\item Calculatrice en mode DEG
\end{itemize}
\end{tcolorbox}
\begin{tcolorbox}[enhanced,colback=white,colframe=popgold,boxrule=0.9pt,arc=3pt,title={5. Calculer un angle},fonttitle=\bfseries\scriptsize,coltitle=white,colbacktitle=popgold,left=4pt,right=4pt,top=2pt,bottom=2pt,boxsep=1pt,toptitle=1.5pt,bottomtitle=1.5pt,halign=left]
\begin{itemize}[nosep,leftmargin=*,itemsep=1pt,font=\scriptsize]
\item Fonctions réciproques
\item $\cos^{-1}, \sin^{-1}, \tan^{-1}$
\item Arrondi au degré
\end{itemize}
\end{tcolorbox}
\begin{tcolorbox}[enhanced,colback=white,colframe=poppink,boxrule=0.9pt,arc=3pt,title={6. Résolution \& Problèmes},fonttitle=\bfseries\scriptsize,coltitle=white,colbacktitle=poppink,left=4pt,right=4pt,top=2pt,bottom=2pt,boxsep=1pt,toptitle=1.5pt,bottomtitle=1.5pt,halign=left]
\begin{itemize}[nosep,leftmargin=*,itemsep=1pt,font=\scriptsize]
\item 2 côtés connus
\item 1 côté + 1 angle
\item Situations concrètes
\end{itemize}
\end{tcolorbox}
\end{tcbraster}

\legende{Carte mentale : Trigonométrie dans le triangle rectangle (3ème)}
\end{popfigure}

\coursec{Bilan de la leçon : Trigonométrie dans le triangle rectangle}

\courssub{Définitions et vocabulaire clés}

\begin{definition}[Triangle rectangle et vocabulaire]
\begin{itemize}
    \item Un \cle{triangle rectangle} est un triangle possédant un angle droit ($90^\circ$). Le côté opposé à l'angle droit est l'\cle{hypoténuse} : c'est toujours le plus grand des trois côtés.
    \item Pour un angle aigu $\alpha$ de ce triangle : le \cle{côté adjacent} à $\alpha$ est le côté de l'angle droit qui touche $\alpha$ ; le \cle{côté opposé} à $\alpha$ est le côté qui fait face à $\alpha$ (il ne touche pas $\alpha$).
\end{itemize}
\end{definition}

\begin{definition}[Rapports trigonométriques]
Dans un triangle rectangle, pour un angle aigu $\alpha$ :
\[
\cos\alpha = \frac{\text{côté adjacent}}{\text{hypoténuse}},\qquad
\sin\alpha = \frac{\text{côté opposé}}{\text{hypoténuse}},\qquad
\tan\alpha = \frac{\text{côté opposé}}{\text{côté adjacent}}.
\]
Moyen mnémotechnique : \textbf{SOH--CAH--TOA} (Sinus = Opposé/Hypoténuse, Cosinus = Adjacent/Hypoténuse, Tangente = Opposé/Adjacent).
\end{definition}

\begin{propriete}[Relations entre les rapports]
Pour tout angle aigu $\alpha$ :
\[
\tan\alpha = \frac{\sin\alpha}{\cos\alpha}
\qquad\text{et}\qquad
\cos^2\alpha + \sin^2\alpha = 1.
\]
La seconde relation découle du théorème de Pythagore ; elle sert de vérification et sera très utile en classe de Seconde.
\end{propriete}

\begin{attention}[Le vocabulaire dépend de l'angle choisi]
Le côté adjacent et le côté opposé se définissent \emph{par rapport à l'angle aigu utilisé}. Si l'on change d'angle aigu, les rôles des deux côtés de l'angle droit s'échangent ! L'hypoténuse, elle, ne change jamais.
\end{attention}

\courssub{Angles remarquables : valeurs exactes à connaître par cœur}

\begin{aretenir}[Tableau des valeurs remarquables]
\begin{center}
\begin{tabularx}{\linewidth}{|c|X|X|X|}\hline
\rowcolor{popblueL}
$\alpha$ & $\cos\alpha$ & $\sin\alpha$ & $\tan\alpha$ \\ \hline
$30^\circ$ & $\dfrac{\sqrt{3}}{2}$ & $\dfrac{1}{2}$ & $\dfrac{\sqrt{3}}{3}$ \\ \hline
$45^\circ$ & $\dfrac{\sqrt{2}}{2}$ & $\dfrac{\sqrt{2}}{2}$ & $1$ \\ \hline
$60^\circ$ & $\dfrac{1}{2}$ & $\dfrac{\sqrt{3}}{2}$ & $\sqrt{3}$ \\ \hline
\end{tabularx}
\end{center}
\vspace{0.2cm}
\emph{Astuce mnémotechnique} : pour le sinus, on écrit $\frac{\sqrt{0}}{2}, \frac{\sqrt{1}}{2}, \frac{\sqrt{2}}{2}, \frac{\sqrt{3}}{2}, \frac{\sqrt{4}}{2}$ pour les angles $0^\circ, 30^\circ, 45^\circ, 60^\circ, 90^\circ$. Le cosinus suit l'ordre inverse. On vérifie bien que $\cos^2 60^\circ + \sin^2 60^\circ = \frac{1}{4} + \frac{3}{4} = 1$.
\end{aretenir}

\courssub{Méthodes express}

\begin{methode}[Calculer une longueur (un angle aigu et un côté connus)]
\begin{enumerate}
    \item Identifier l'angle aigu $\alpha$ utilisé et nommer les côtés (hypoténuse, adjacent, opposé) \emph{par rapport à} $\alpha$.
    \item Choisir le rapport ($\cos$, $\sin$ ou $\tan$) qui relie le côté \textbf{cherché} au côté \textbf{connu}.
    \item Écrire l'égalité trigonométrique, par exemple $\cos\alpha = \dfrac{\text{adj}}{\text{hyp}}$.
    \item Isoler l'inconnue par un produit en croix, puis calculer avec la calculatrice en mode \textbf{DEG} (degrés).
    \item Conclure par une phrase avec l'unité : « Donc $AB \approx 5,4$~cm ».
\end{enumerate}
\end{methode}

\begin{methode}[Calculer un angle aigu (deux côtés connus)]
\begin{enumerate}
    \item Nommer les côtés par rapport à l'angle cherché $\alpha$.
    \item Choisir le rapport qui utilise les \textbf{deux côtés connus}.
    \item Calculer la valeur numérique de ce rapport (division).
    \item Appliquer la fonction réciproque de la calculatrice : $\alpha = \cos^{-1}(\ldots)$, $\alpha = \sin^{-1}(\ldots)$ ou $\alpha = \tan^{-1}(\ldots)$ (touches \texttt{SHIFT} + \texttt{cos}, \texttt{sin} ou \texttt{tan}).
    \item Arrondir au degré près (ou au dixième de degré selon la consigne) et conclure.
\end{enumerate}
\end{methode}

\begin{methode}[Résoudre complètement un triangle rectangle]
Données possibles : deux côtés, ou bien un côté et un angle aigu.
\begin{enumerate}
    \item Faire un schéma annoté avec toutes les données.
    \item Calculer le troisième angle : les deux angles aigus sont complémentaires, donc angle inconnu $= 90^\circ - \text{angle aigu connu}$.
    \item Calculer les côtés manquants avec les rapports trigonométriques (le théorème de Pythagore peut servir de vérification).
    \item Récapituler les résultats : les 3 côtés et les 3 angles.
\end{enumerate}
\end{methode}

\begin{exoresolu}[Calcul d'une longueur]
Soit $ABC$ un triangle rectangle en $A$ tel que $\widehat{ABC} = 35^\circ$ et $BC = 8$~cm. Calculer $AB$ (arrondir au dixième).
\tcblower
\textbf{Solution.} Dans le triangle $ABC$ rectangle en $A$, l'hypoténuse est $BC$. Par rapport à l'angle $\widehat{ABC} = 35^\circ$, le côté $AB$ est le côté adjacent. On utilise donc le cosinus :
\[
\cos \widehat{ABC} = \frac{AB}{BC}
\quad\Longrightarrow\quad
AB = BC \times \cos 35^\circ = 8 \times \cos 35^\circ.
\]
La calculatrice (en mode DEG) donne $\cos 35^\circ \approx 0,819$, d'où $AB \approx 8 \times 0,819 \approx 6,6$~cm.\\
\textbf{Conclusion :} $AB \approx 6,6$~cm. Vérification : $AB < BC$, cohérent car l'hypoténuse est le plus grand côté.
\end{exoresolu}

\begin{exoresolu}[Calcul d'un angle aigu]
Soit $MNP$ un triangle rectangle en $M$ tel que $MN = 5$~cm et $MP = 7$~cm. Calculer la mesure de l'angle $\widehat{MPN}$ (arrondir au degré).
\tcblower
\textbf{Solution.} Par rapport à l'angle $\widehat{MPN}$ : le côté adjacent est $MP = 7$~cm et le côté opposé est $MN = 5$~cm. On utilise la tangente :
\[
\tan \widehat{MPN} = \frac{MN}{MP} = \frac{5}{7} \approx 0,714.
\]
Avec la fonction réciproque : $\widehat{MPN} = \tan^{-1}\left(\dfrac{5}{7}\right) \approx 36^\circ$.\\
\textbf{Conclusion :} $\widehat{MPN} \approx 36^\circ$.
\end{exoresolu}

\begin{exoresolu}[Problème concret : hauteur d'un mât]
Pour installer une antenne MTN dans un quartier de Yaoundé, un technicien se place à $20$~m du pied d'un mât vertical. Il vise le sommet du mât sous un angle de $40^\circ$ avec l'horizontale. Quelle est la hauteur du mât (arrondir au dixième de mètre) ?
\tcblower
\textbf{Solution.} La situation se modélise par un triangle rectangle : le mât est le côté opposé à l'angle de $40^\circ$, la distance au sol ($20$~m) est le côté adjacent. On utilise la tangente :
\[
\tan 40^\circ = \frac{\text{hauteur}}{20}
\quad\Longrightarrow\quad
\text{hauteur} = 20 \times \tan 40^\circ \approx 20 \times 0,839 \approx 16,8~\text{m}.
\]
\textbf{Conclusion :} le mât mesure environ $16,8$~m de haut.
\end{exoresolu}

\courssub{Pièges à éviter (check-list de rédaction)}

\begin{attention}[Erreurs fréquentes au BEPC]
\begin{itemize}
    \item \textbf{Mode calculatrice} : toujours vérifier que la calculatrice est en mode \texttt{DEG} (degrés), et non \texttt{RAD} ou \texttt{GRAD}.
    \item \textbf{Nommer les côtés} \emph{par rapport à l'angle utilisé} : adjacent et opposé s'échangent si l'on change d'angle aigu !
    \item \textbf{Inversion du rapport} : de $\cos\alpha = \dfrac{\text{adj}}{\text{hyp}}$ on tire $\text{hyp} = \dfrac{\text{adj}}{\cos\alpha}$, et non $\text{adj} \times \cos\alpha$.
    \item \textbf{Arrondis} : ne pas arrondir les valeurs intermédiaires ; n'arrondir que le résultat final.
    \item \textbf{Unités} : longueurs en cm, m, km\ldots ; angles en degrés ($^\circ$). Ne jamais oublier l'unité dans la phrase de conclusion.
    \item \textbf{Cohérence} : l'hypoténuse doit rester le plus grand côté, et un angle aigu est toujours compris entre $0^\circ$ et $90^\circ$.
\end{itemize}
\end{attention}

\begin{aretenir}[Checklist avant le BEPC]
\begin{itemize}
    \item $\square$ Je sais nommer hypoténuse, côté adjacent, côté opposé par rapport à un angle aigu.
    \item $\square$ Je connais par cœur les définitions de $\cos$, $\sin$, $\tan$ (SOH--CAH--TOA).
    \item $\square$ Je connais les valeurs exactes de $\cos$, $\sin$, $\tan$ pour $30^\circ$, $45^\circ$, $60^\circ$.
    \item $\square$ Je sais choisir le bon rapport trigonométrique selon les données du problème.
    \item $\square$ Je sais calculer une longueur inconnue (produit en croix, calculatrice en mode DEG).
    \item $\square$ Je sais calculer un angle aigu inconnu (fonctions réciproques $\cos^{-1}$, $\sin^{-1}$, $\tan^{-1}$).
    \item $\square$ Je sais résoudre complètement un triangle rectangle (3 côtés, 3 angles).
    \item $\square$ Je sais modéliser un problème concret (ombre, pente, hauteur d'arbre, câble, antenne\ldots) par un triangle rectangle.
    \item $\square$ Je rédige ma démarche : schéma, choix du rapport, formule, calcul, phrase de conclusion avec unité.
    \item $\square$ Je vérifie la cohérence de mes résultats (hypoténuse $>$ chaque côté, angles aigus $< 90^\circ$, somme des angles $= 180^\circ$).
    \item $\square$ Je gère mon temps : exercices types (calcul direct, problème) en moins de 10 minutes chacun.
\end{itemize}
\end{aretenir}

\findecours

\end{document}
